\chapter*{概率专题练习}

\begin{enumerate}
    \item 事件的分类
    \begin{enumerate}
        \item 以下事件是随机事件的是\dotfill{(\kh{C})}
        \xx{下雨屋顶湿}{秋后柳叶黄}{有水就有鱼}{水结冰体积变大}
        \begin{jieda}
            由题意得 ${A},{B},{D}$ 的概率为 1,所以是必然事件，C 的概率不为 0 , 也不为 1 , 所以它是随机事件, 故 C 正确.
        \end{jieda}
        
        \item 将一根长为 $a$ 的铁丝随意截成三段，这三段铁丝构成一个三角形, 此事件是\dotfill{(\kh{C})}
        \xx{必然事件}{不可能事件}{随机事件}{不能判定}
        
        \item 给出下列四个命题，其中正确命题的序号是\blkda{\ding{172}\ding{173}\ding{175}}
        \begin{dingautolist}{172}
            \item “三个球全部放入两个盒子，其中必有一个盒子有一个以上的球”是必然事件
            \item “当 $x$ 为某一实数时，可使 ${x}^{2} < 0$ ”是不可能事件
            \item “明天上海要下雨”是必然事件
            \item “从 100 个灯泡中取出 5 个， 5 个都是次品”是随机事件
        \end{dingautolist}
    \end{enumerate}
    
    \item 事件与样本空间
    \begin{enumerate}
        \item 高一(1)班计划从 $A,B,C,D,E$ 这五名班干部中选两人代表班级参加一次活动，则样本空间中样本点的个数为\blkda{10}.
        \begin{jieda}
            $C_5^2 = 10$
        \end{jieda}
        \item 随机事件“连续掷一颗筛子直到出现 5 点停止，观察掷的次数” 的样本空间是\dotfill{(\kh{D})}
        \xx{5}{1 到 6 的正整数}{6}{一切正整数}
        \item 若随机试验的样本空间为 $\Omega  = \{ 0,1,2\}$ ，则下列说法不正确的是\dotfill{(\kh{D})}
        \xx{事件 $P = \{ 1,2\}$ 是随机事件}{事件 $Q = \{ 0,1,2\}$ 是必然事件}{事件 $M = \{  - 1, - 2\}$ 是不可能事件}{事件 $\{  - 1,0\}$ 是随机事件}
        \begin{jieda}
            随机试验的样本空间为 $\Omega  = \{ 0,1,2\}$ ,
            
            则事件 $P = \{ 1,2\}$ 是随机事件，故 A 正确；
            
            事件 $Q = \{ 0,1,2\}$ 是必然事件,故 ${B}$ 正确;
            
            事件 $M = \{  - 1, - 2\}$ 是不可能事件,故 ${C}$ 正确;
            
            事件 $\{  - 1,0\}$ 是不可能事件,故 ${D}$ 错误.
        \end{jieda}
        \item 一个家庭有两个小孩，则样本空间可以是\blkda{\ding{172}\ding{174}}
        \begin{dingautolist}{172}
            \item \{(男, 女), (男, 男), (女, 女)\}
            \item \{(男, 女), (女, 男)\}
            \item \{(男, 男), (男, 女), (女, 男), (女, 女)\}
            \item \{(男, 男), (女, 女)\}
        \end{dingautolist}
    \end{enumerate}

    \item 写出基本事件
    \begin{enumerate}
        \item 有 4 张卡片，上面分别写有数字 1, 2, 3, 4, 从这 4 张卡片中随机抽取 2 张，则取出的 2 张卡片上的数字之和为奇数的所有基本事件数为\blkda{4}
        \begin{jieda}
            取出的 2 张卡片上的数字之和为奇数所包含的基本事件为:\{1,2\}, \{1,4\}, \{3,2\}, \{3,4\}
        \end{jieda}
        
        \item 先后抛掷 2 枚均匀的一分、二分的硬币，观察落地后硬币的正反面情况，则下列事件包含 3 个基本事件的是\dotfill{(\kh{A})}
        \xx{“至少一枚硬币正面向上”}{“只有一枚硬币正面向上”}{“两枚硬币都是正面向上”}{“两枚硬币一枚正面向上，另一枚反面向上”}
        \begin{jieda}
            先后抛掷 2 枚均匀的一分、二分的硬币, 所包含的基本事件有\{正, 正\}、\{正, 反\}、\{反, 正\}、\{反, 反\},
            
            “至少一枚硬币正面向上” 包含的基本事件有 \{正, 正\}、\{正, 反\}、\{反，正\} 共三个，故 A 正确；
            
            “只有一枚硬币正面向上”。任含的基本事件有\{正，反\}、\{反，正\}共两个，故 B 错；
            
            “两枚硬币都是正面向上”包含的基本事件有\{正, 正 \} 共一个，故 C 错；
            
            “两枚硬币一枚正面向上，另一枚反面向上” 包含的基本事件有\{正，反\}、\{反，正\}共两个，故 D 错.
        \end{jieda}
        
        \item 有 5 根木棍，其长度分别为 2, 3, 4, 5, 6，从这 5 根木棍中任取 3 根，首尾相接能构成三角形的有\blkda{7}个.
        \begin{jieda}
            根据三角形任意两边之和大于第三边，所以能组成三角形的有: 2, 3, 4; 2, 4, 5; 2, 5, 6; 3, 4, 5; 3, 4, 6; 3, 5, 6; 4, 5, 6, 共 7 个
        \end{jieda}
        
        \item 一个箱子中装有编号分别为 1、2、3、4、5的 5 个小球， 5 个小球除编号外其他均无异，现有事件 $A$ 为“从箱中任取 3 个小球观察其编号”，问事件 $A$ 的样本点数有\blkda{10}个.
        \begin{jieda}
            因为事件 $A$ 为“从箱中任取 3 个小球观察其编号”，则事件 $A$ 包含的样本点有 $\left( {1,2,3}\right) ,\left( {1,2,4}\right) ,\left( {1,2,5}\right) ,\left( {1,3,4}\right) ,\left( {1,3,5}\right) ,\left( {1,4,5}\right) ,\left( {2,3,4}\right)$, $\left( {2,3,5}\right) ,\left( {2,4,5}\right) ,\left( {3,4,5}\right)$ 共 10 种.
        \end{jieda}
    \end{enumerate}

    \item 计算古典概型问题的概率
    \begin{enumerate}
        \item 掷一个骰子，观察朝上的面的点数，设事件 $M =$ “点数为奇数”，事件 $N{ = }$ “点数为3的整数倍”,若 $P\left( M\right) ,P\left( N\right)$ 分别表示事件 $M,N$ 发生的概率,则\dotfill{(\kh{B})}
        \xx{$P\left( M\right)  = \dfrac{1}{3},P\left( N\right)  = \dfrac{1}{2}$}{$P\left( M\right)  = \dfrac{1}{2},P\left( N\right)  = \dfrac{1}{3}$}{$P\left( M\right)  = P\left( N\right)  = \dfrac{1}{2}$}{$P\left( M\right)  = P\left( N\right)  = \dfrac{1}{3}$}
        \begin{jieda}
            首先，我们知道投掷的点数有1,2,3,4,5,6，对于 $M$ ,符合条件的有1,3,5,对于 $N$ ,符合条件的有3,6，故 $P\left( M\right)  = \dfrac{3}{6} = \dfrac{1}{2},P\left( N\right)  = \dfrac{2}{6} = \dfrac{1}{3}$ ，故 ${B}$ 正确.
        \end{jieda}

        \item 某公园有东、南、西、北共 4 个大门供游客出入，小军、 小明从不同的大门进入公园游玩，游玩结束后，他们随机地从其中一个大门离开，则他们恰好从同一个大门出去的概率是\blkda{$\dfrac{1}{4}$}
        
        \item 在山西的某个旅游景点内有刀削面、油炸糕、糖火烧、炕馍、莜面这 5 种传统小吃. 某游客从中随机选择 3 种品尝, 则该游客选择了油炸糕和莜面品尝的概率为\blkda{$\dfrac{3}{10}$}
        \begin{jieda}
            $\dfrac{C_3^1}{C_5^3} = \dfrac{3}{10}$
        \end{jieda}
        
        \item 某汽车站每天均有 3 辆开往省城的分为上、中、下等级的客车, 某天袁先生准备在该汽车站乘车前往省城办事，但他不知道客车的车况，也不知道发车顺序. 为了尽可能乘坐上等车，他采取如下策略:先放过一辆，如果第二辆比第一辆好则上第二辆，否则上第三辆，则他乘坐上等车的概率为\blkda{$\dfrac{1}{2}$}
        \begin{jieda}
            根据题意, 所有可能的客车通过顺序的情况为

            (上, 中, 下), (上, 下, 中), (中, 上, 下)，(中, 下, 上) , (下, 中, 上) , (下, 上, 中) , 共 6 种,
            
            其中该人可以乘坐上等车的情况有(中，上，下)，(中，下，上)，(下，上，中)，共 3 种， 则其概率为 $\dfrac{3}{6} = \dfrac{1}{2}$ .
        \end{jieda}
        \item 从大小相同，编号为 1,2,3,4,5 的 5 个小球中，选取 3 个小球，求下列事件的概率:
        \begin{enumerate}
            \item 编号为 1,2 的小球同时被取到的概率
            \item 所取到的三个小球的编号之和为偶数的概率
        \end{enumerate}
        \begin{jieda}
            \begin{enumerate}
                \item 从编号为 1,2,3,4,5 的 5 个小球中任意取出 3 个，样本空间为 $\Omega  = \{ {123},{124},{125},{134},{135},{145},{234},{235},{245},{345}\}$

                记事件 $A$ 为“编号为1,2的小球同时被取到”，则 $A = \{ {123},{124},{125}\}$ ,故 $P\left( A\right)  = \dfrac{3}{10}$ .

                \item 记事件 $B$ 为“所取到的三个小球编号之和为偶数”，则 $B = \{ {123},{125},{134},{145},{235},{345}\}$ ,故 $P\left( B\right)  = \dfrac{6}{10} = \dfrac{3}{5}$
            \end{enumerate}
        \end{jieda}
        
        \item 已知某医疗队共有医生 20 人，护士 30 人，现在要用分层随机抽样的方法从中选取 5 人组建一个救援小组.
        \begin{enumerate}
            \item 求救援小组中医生和护士的人数
            \item 若从救援小组中随机选取 2 人担任组长，求医生和护士各有 1 人被选中的概率
        \end{enumerate}
        \begin{jieda}
            \begin{enumerate}
                \item 由题可知救援小组中医生的人数为 $\dfrac{20}{{20} + {30}} \times  5 = 2$ ，护士的人数为 $\dfrac{30}{{20} + {30}} \times  5 = 3$
                \item 由(1)可知救援小组中有 2 名医生，记为 $A$ ， $B$ ；有 3 名护士，记为 $a$ ， $b$ ， $c$ .
                
                从中随机选取 2 人担任组长,所有的结果为 $\left( {A,B}\right) ,\left( {A,a}\right) ,\left( {A,b}\right) ,\left( {A,c}\right) ,\left( {B,a}\right)$ ,
                
                $\left( {B,b}\right) ,\left( {B,c}\right) ,\left( {a,b}\right) ,\left( {a,c}\right) ,\left( {b,c}\right)$ ,共有 10 种可能的结果.
                
                记事件 $M$ 为 “医生和护士各有 1 人被选中担任组长”，依题意可知事件 $M$ 包含的样本点
                
                有 $\left( {A,a}\right) ,\left( {A,b}\right) ,\left( {A,c}\right) ,\left( {B,a}\right) ,\left( {B,b}\right) ,\left( {B,c}\right)$ ,共有 6 种可能的结果.
                
                故 $P\left( M\right)  = \dfrac{6}{10} = \dfrac{3}{5}$
            \end{enumerate}
        \end{jieda}
        
        \item 在一个盒子中有除颜色外完全相同的 3 个球，蓝球，红球，绿球各 1 个, 从中随机地取出 1 个球, 观察其颜色后放回, 然后随机取出 1 个球.
        \begin{enumerate}
            \item 请用适当的符号表示试验的可能结果，写出试验的样本空间
            \item 求“第一次取出的是红球”的概率
            \item 求“第一次取出的是红球且两次取出的球颜色不同”的概率
        \end{enumerate}
        \begin{jieda}
            \begin{enumerate}
                \item 样本空间 $\Omega  = \{$ (蓝球，蓝球)，(蓝球，红球)，(蓝球，绿球)，(红球，蓝球)， (红球, 红球), (红球, 绿球) (绿球, 蓝球) (绿球, 红球) (绿球, 绿球) \}.
                \item 设“第一次取出的是红球”为事件 $A$ .
                
                事件 $A$ 包含的样本点有 (红球,蓝球),(红球,红球)，(红球，绿球)，共 3 个，
                
                由( 1 )知基本事件总数 $n = 9$ ，所以 $P\left( A\right)  = \dfrac{3}{9} = \dfrac{1}{3}$

                \item 设“第一次取出的是红球且两次取出的球颜色不同”为事件 $B$ .
                
                事件 $B$ 包含的样本点有 (红球,蓝球),(红球,绿球),共 2 个, 由(1)知基本事件总数 $n = 9$，所以 $P\left( B\right)  = \dfrac{2}{9}$
            \end{enumerate}
        \end{jieda}

        \item 一只不透明的袋子中装有 2 个白球，3 个红球，这些球除颜色外都相同.
        \begin{enumerate}
            \item 搅匀后从中任意摸出 2 个球，求这 2 个都球是白球的概率
            \item 搅匀后从中任意摸出 1 个球，记录颜色后放回，搅匀，再从中任意摸出 1 个球，求 2 次摸到的球恰好是 1 个白球和 1 个红球的概率.
        \end{enumerate}
        \begin{jieda}
            \begin{enumerate}
                \item 将 3 个红球记为红 1 ，红 2 ，红 3 ，2 个白球记为白 1 ，白 2 ，
                
                则任意摸出 2 个球的样本空间 $\Omega$ 有: 红 1 红 2，红 1 红 3，红 1 白 1，红 1 白 2，红 2 红 3，红 2 白 1，红 2 白 2，红 3 白 1，红 3 白 2，白 1 白 2 共 10 个样本点，

                其中 2 球均为白球事件的样本点只有 1 个，因此 2 个球都是白球概率为 $\dfrac{1}{10}$

                \item 搅匀后从中任意摸出 1 个球，记录下颜色后放回袋子中并搅匀，再从中任意摸出 1 个球，将 3 个红球记为红 1 ，红 2 ，红 3 ，2 个白球记为白 1 ，白 2 ，列表如图所示:

                \begin{center}
                
                \begin{tabular}{|c|c|c|c|c|c|}
                \hline
                第 2 次 \& 第 1 次 & 红 1 & 红 2 & 红 3 & 白 1 & 白 2 \\
                \cline{1-6}
                红 1 & (红 1, 红 1) & (红 1, 红 2) & (红 1, 红 3) & (红 1, 白 1) & (红 1, 白 2) \\
                \cline{1-6}
                红 2 & (红 2, 红 1) & (红 2, 红 2) & (红 2, 红 3) & (红 2, 白 1) & (红 2, 白 2) \\
                \cline{1-6}
                红 3 & (红 3, 红 1) & (红 3, 红 2) & (红 3, 红 3) & (红 3, 白 1) & (红 3, 白 2) \\
                \cline{1-6}
                白 1 & (白 1, 红 1) & (白 1, 红 2) & (白 1, 红 3) & (白 1,白 1) & (白 1,白 2) \\
                \cline{1-6}
                白 2 & (白2, 红1) & (白2, 红2) & (白 2, 红 3) & (白 2,白 1) & (白 2, 白 2) \\
                \cline{1-6}
                \hline
                \end{tabular}
                \end{center}
                
                所以搅匀后从中任意摸出 1 个球，记录下颜色后放回袋子中并搅匀，再从中任意摸出 1 个球事件的样本空间共有 25 个样本点, 它们出现的可能性相同,
                
                其中满足事件 $B =$ " 2 次摸到的球恰好是 1 个白球和 1 个红球"的样本点有 12 个，所以 $P\left( B\right)  = \dfrac{12}{25}$
                
            \end{enumerate}
        \end{jieda}
        
        \item 新高考数学试题中有多项选择题，要求为:“在每小题给出的四个选项中, 有多项符合题目要求（必然有多个但不是全部正确选项）, 全部选对的得 5 分, 部分选对的得 2 分, 有选错的得 0 分. ”已知某多项选择题的正确答案是 BCD.
        \begin{enumerate}
            \item 若甲同学不会做该题，但他想猜对得 5 分，就随机填写了一个答案，求他得 5 分的概率
            \item 若乙同学也不会做该题，他只想得 2 分，就按单项选择题处理，随机填写了一个答案，求他得 2 分的概率
        \end{enumerate}
        \begin{jieda}
            \begin{enumerate}
                \item 该事件的样本空间 $\Omega  = \{ {{AB}},{{AC}},{{AD}},{{BC}},{{BD}},{{CD}},{{ABC}},{{ABD}},{{ACD}},{{BCD}}\}$ ，共 10 个样本点， 且每个样本点的发生是等可能的，故可以用古典概型计算概率.
                
                记事件 $A$ 表示“甲同学得 5 分”，则 $A = \{ {{BCD}}\}$ ，含有 1 个样本点，所以 $P\left( A\right)  = \dfrac{1}{10}$ .

                \item 该事件的样本空间 $\Omega  = \{ {A},{B},{C},{D}\}$ ，共 4 个样本点，每个样本点的发生是等可能的，故可以用古典概型计算概率，记事件 $B$ 表示“乙同学得 2 分”，则 $B = \{ {B},{C},{D}\}$ ，含有 3 个样本点，所以 $P\left( B\right)  = \dfrac{3}{4}$
            \end{enumerate}
        \end{jieda}
        
        \item 某电脑公司现有 $A$ ， $B$ ， $C$ 三种型号的甲品牌电脑和 $D$ ， $E$ 两种型号的乙品牌电脑, 某中学要从甲、乙两种品牌电脑中各随机选购一种型号的电脑.
        \begin{enumerate}
            \item 写出所有选购方案
            \item 如果(1)中各种选购方案被选中的可能性相同，那么 $A$ 型号电脑被选中的概率是多少？(直接写出结果即可)
        \end{enumerate}
        \begin{jieda}
            \begin{enumerate}
                \item 选购方案为 $\left( {A,D}\right) ,\left( {A,E}\right) ,\left( {B,D}\right) ,\left( {B,E}\right) ,\left( {C,D}\right) ,\left( {C,E}\right)$
                \item $A$ 型号电脑被选中的情形为 $\left( {A,D}\right) ,\left( {A,E}\right)$ ，即含 2 个样本点，所以 $A$ 型号电脑被选中的概率为 $P = \dfrac{2}{6} = \dfrac{1}{3}$ .
            \end{enumerate}
        \end{jieda}
        
        \item 在一个不透明的袋子中装有大小、质地完全相同的 1 个红球和 1 个白球,每次从袋子中随机摸出 1 个球,观察其颜色后放回. 甲连续摸球 2 次,乙连续摸球 4 次. 用 $a$ 表示摸出红球, $b$ 表示摸出白球.
        \begin{enumerate}
            \item 分别写出甲和乙的摸球试验的样本空间及其包含样本点的个数
            \item 设 $A =$ “甲恰有一次摸出红球”， $B =$ “乙恰有两次摸出红球”，比较 $P\left( A\right)$ 与 $P\left( B\right)$ 的大小.
        \end{enumerate}
        \begin{jieda}
            \begin{enumerate}
                \item 甲摸球试验的样本空间: $\{ \left( {a,a}\right) ,\left( {a,b}\right) ,\left( {b,a}\right) ,\left( {b,b}\right) \}$ ，样本点 4 个.
                
                乙摸球试验的样本空间:
                
                $\{ \left( {a,a,a,a}\right) ,\left( {a,a,a,b}\right) ,\left( {a,a,b,a}\right) ,\left( {a,b,a,a}\right) ,\left( {b,a,a,a}\right) ,\left( {a,a,b,b}\right) ,\left( {a,b,b,a}\right) ,\left( {b,b,a,a}\right) ,$
                
                $\left( {a,b,a,b}\right) ,\left( {b,a,b,a}\right) ,\left( {b,a,a,b}\right) ,\left( {a,b,b,b}\right) ,\left( {b,a,b,b}\right) ,\left( {b,b,a,b}\right) ,\left( {b,b,b,a}\right) ,\left( {b,b,b,b}\right) \}$ , 样本点 16 个.

                \item 由(1)得 $P\left( A\right)  = \dfrac{2}{4} = \dfrac{1}{2},P\left( B\right)  = \dfrac{6}{16} = \dfrac{3}{8}$ ，所以 $P\left( A\right)  > P\left( B\right)$ .
            \end{enumerate}
        \end{jieda}
        
        \item 为了防止注册账号被他人非法登录，某系统在账号登录前，要先输入验证码. 已知该系统登入设置的每个验证码均由有序数字串 ${abc}$ 组成，其中 $a,b,c \in  \{ 1,2,3\}$ ，某人非法登录一个账号，任选一组验证码输入.
        \begin{enumerate}
            \item 求这个人输入的验证码恰有两位正确的概率
            \item 若这个人通过技术获得了验证码的第一位数，求这个人输入的验证码正确的概率
        \end{enumerate}
        \begin{jieda}
            \begin{enumerate}
                \item 由题可知，所有的验证码包括 $3\times3\times3 = 27$ 种.
                
                恰有两位正确的验证码有$C_3^2C_2^1 = 6$种， 故这个人输入的验证码恰有两位正确的概率为 $\dfrac{6}{27} = \dfrac{2}{9}$ .

                \item 这个人输入的验证码有 $3\times3 = 9$ 种，则这个人输入的验证码正确的概率为 $\dfrac{1}{9}$
            \end{enumerate}
        \end{jieda}
        \item 在某项大型活动中,甲、乙等五名志愿者被随机地分到 $A\text{、}B\text{、}C\text{、}D$ 四个不同的岗位服务,每个岗位至少有一名志愿者.
        \begin{enumerate}
            \item 求甲、乙两人同时参加 $A$ 岗位服务的概率
            \item 求甲、乙两人不在同一个岗位服务的概率
            \item 求五名志愿者中仅有一人参加 $A$ 岗位服务的概率
        \end{enumerate}
        \begin{jieda}
        由题意可知, 有两人在同一岗位服务, 因此, 第一步先从 5 个人中选两个人作为一组，第二步再将这两个人和其他 3 个人一起进行排列,共有 ${{C}}_{5}^{2}{{P}}_{4}^{4} = {240}$ (种) 安排方法,即样本空间，$\Omega$ 中样本点个数为 ${{C}}_{5}^{2}{{P}}_{4}^{4} = {240}$ .
            \begin{enumerate}
                \item 记“甲、乙两人同时参加 $A$ 岗位服务”为事件 $X$ ,则由计数原理可得 $\left| X\right|  = {{P}}_{3}^{3} = 6$ ,所以,事件 $X$ 的概率为 $P\left( X\right)  = \dfrac{\left| X\right| }{\left| \Omega \right| } = \dfrac{{{P}}_{3}^{3}}{{{C}}_{5}^{2}{{P}}_{4}^{4}} = \dfrac{1}{40}$ .
                
                \item 记“甲、乙两人不在同一个岗位服务”为事件 $Y$ ,则由计数原理可得 $\left| Y\right|  = {{C}}_{5}^{2}{{P}}_{4}^{4} - {{P}}_{4}^{4}$，所以事件 $Y$ 的概率为 $P\left( Y\right)  = \dfrac{\left| Y\right| }{\left| \Omega \right| } = \dfrac{{{C}}_{5}^{2}{{P}}_{4}^{4} - {{P}}_{4}^{4}}{{{C}}_{5}^{2}{{P}}_{4}^{4}} = \dfrac{9}{10}$
                
                \item 记“仅有一人参加 $A$ 岗位服务”为事件 $Z$ ,则由计数原理可得 $\left| Z\right|  = {{C}}_{5}^{2}{{P}}_{4}^{4} - {{C}}_{5}^{2}{{P}}_{3}^{3}$，所以事件 $Z$ 的概率为 $P\left( Z\right)  = \dfrac{\left| Z\right| }{\left| \Omega \right| } = \dfrac{{{C}}_{5}^{2}{{P}}_{4}^{4} - {{C}}_{5}^{2}{{P}}_{3}^{3}}{{{C}}_{5}^{2}{{P}}_{4}^{4}} = \dfrac{3}{4}$ .
            \end{enumerate}
        \end{jieda}
    \end{enumerate}
    
    \item 游戏的公平性问题
    \begin{enumerate}
        \item 小明与小华两人玩游戏，则下列游戏不公平的是\dotfill{(\kh{B})}
        \xx{抛掷一枚骰子，向上的点数为奇数，小明获胜，向上的点数为偶数，小华获胜}{同时抛掷两枚硬币，恰有一枚正面向上，小明获胜，两枚都正面向上，小华获胜}{从一副不含大小王的扑克牌中抽一张，扑克牌是红色，小明获胜，扑克牌是黑色，小华获胜}{小明、小华两人各写一个数字 6 或 8 ，如果两人写的数字相同，小明获胜，否则小华获胜}
        
        \item 用掷两枚硬币做胜负游戏，规定:两枚硬币同时出现正面或同时出现反面算甲胜, 一个正面、一个反面算乙胜. 这个游戏公平吗?
        \begin{jieda}
            这个游发是公平的,理由:抛掷两枚硬币共有 4 种等可能结果:(正，正)，(正，反)，(反，正)，(反，反)，所以甲、乙获胜的概率都是 $\dfrac{1}{2}$
        \end{jieda}
        \item 下面的三个游戏都是在袋子中装球, 然后从袋子中不放同地取球, 分别计算三个游戏中甲获胜的概率，你认为哪个游戏是公平的？
        
        \begin{center}
        \begin{tabular}{|c|c|c|c|}
        \hline
        \phantom{X} & 游戏 1 & 游戏 2 & 游戏 3 \\
        \cline{1-4}
        袋子中球的数量和颜色 & 1 个红球和 1 个白球 & 2 个红球和 2 个白球 & 3 个红球和 1 个白球 \\
        \cline{1-4}
        取球规则 & 取 1 个球 & 依次取出 2 个球 & 依次取出 2 个球 \\
        \cline{1-4}
        {获胜规则} & 取到红球 $\rightarrow$ 甲胜 & 两个球同色 $\rightarrow$ 甲胜 & 两个球同色 $\rightarrow$ 已胜 \\
        \cline{2-4}
         & 取到白球 $\rightarrow$ 乙胜 & 两个球不同色 $\rightarrow$ 乙胜 & 两个球不同色 $\rightarrow$ 乙胜 \\
        \cline{1-4}
        \hline
        \end{tabular}
        \end{center}
        \begin{jieda}
            游戏 1 中,甲获胜的概率为 $\dfrac{1}{2}$
            
            游戏 2 中，甲获胜的视率为 $\dfrac{2}{6} = \dfrac{1}{3}$ .
            
            游戏 3 中，甲获胜的概率为 $\dfrac{3}{6} = \dfrac{1}{2}$ ，
            
            所以游戏 1 和游戏 3 是公平的.
        \end{jieda}

        \item 一天，甲拿出一个装有三张卡片的盒子(一张卡片的两面都是绿色，一张卡片的两面都是蓝色，还有一张卡片一面是绿色，另一面是蓝色)，跟乙说玩一个游戏，规则是: 甲将盒子里的卡片顺序打乱后，由乙随机抽出一张卡片放在桌子上，然后卡片朝下的面的颜色决定胜负， 如果朝下的面的颜色与朝上的面的颜色一致, 则甲赢, 否则甲输.乙对游戏的公平性提出了质疑, 但是甲说: “当然公平! 你看, 如果朝上的面的颜色为绿色, 则这张卡片不可能两面都是蓝色, 因此朝下的面要么是绿色，要么是蓝色，因此，你赢的概率为 $\dfrac{1}{2}$ ，我赢的概率也是 $\dfrac{1}{2}$ ，怎么不公平？”分析这个游戏是否公平.
        \begin{jieda}
            把卡片六个面的颜色记为 ${9}_{1},{9}_{2},{9}_{3},{B}_{1},{B}_{2},{B}_{3}$ ,
            
            其中，G 表示绿色，B 表示蓝色； ${9}_{3}$ 和 ${B}_{3}$ 是两面颜色不一样的那张卡片的颜色.
            
            游戏所有的结果可以用如图表示.
            
            朝上的面
            
            \begin{center}
            \includegraphics[width=0.2\textwidth]{bodyxb/src/bo_d67fmgv7aajc739ac530_13_304_758_329_117_0.jpg}
            \end{center}
            
            
            朝下的面 $\begin{array}{llllll} {9}_{2} & {9}_{1} & {B}_{3} & {B}_{2} & {B}_{1} & {9}_{3} \end{array}$
            
            不难看出, 此时, 样本空间中共有 6 个样本点, 朝上的面与朝下的面颜色不一致的情况只有 2 种, 因此乙赢的概率为 $\dfrac{2}{6} = \dfrac{1}{3}$
            
            因此, 这个游戏不公平.
        \end{jieda}
        
        \item 已知 $n$ 是一个三位正整数，若 $n$ 的个位数字大于十位数字，十位数字大于百位数字，则称 $n$ 为 “三位递增数” (如 135,256,345 等)现要从甲乙两名同学中，选出一个参加某市组织的数学竞赛，选取的规则如下:从由 1, 2, 3, 4, 5, 6 组成的所有“三位递增数”中随机抽取 1 个数，且只抽取 1 次，若抽取的“三位递增数”是偶数，则甲参加数学竞赛；否则，乙参加数学竞赛.
        \begin{enumerate}
            \item 由1,2,3,4,5,6可组成多少“三位递增数”? 并一一列举出来.
            \item 这种选取规则对甲乙两名学生公平吗？并说明理由.
        \end{enumerate}
        \begin{jieda}
            \begin{enumerate}
                \item 由题意知，所有由 1,2,3,4,5,6 组成的“三位递增数共有 20 个.分别是123,124,125,126,134,135,136,145,146,156,234,235,236,245,246,256,345,346, 356, 456.

                \item 不公平由(1)知，所有由 1，2，3，4，5，6 组成的“三位递增数”有 20 个，记“甲参加数学竞赛”为事件 A，记“乙参加数学竞赛”为事件 B. 则事件 A 含有基本事件有:124，134，234，126，136，146，156，236，246, 256, 346, 356, 456 共 13 个.
                
                由古典概型计算公式, 得$P\left( A\right)  = \dfrac{\text{ 事件 }A\text{ 含有的基本事件的个数 }}{\text{ 试验所有基本事件的总数 }} = \dfrac{13}{20}$
                
                又 $A$ 与 $B$ 对立,所以 $P\left( B\right)  = 1 - P\left( A\right)  = 1 - \dfrac{13}{20} = \dfrac{7}{20}$ ,
                
                所以 $P\left( A\right)  > P\left( B\right)$ . 故选取规则对甲、乙两名学生不公平.
            \end{enumerate}
        \end{jieda}
    \end{enumerate}

    \item 频率的计算
    \begin{enumerate}
        \item 某地一种植物一年生长的高度如下表:
        
        \begin{center}
        \begin{tabular}{|c|c|c|c|c|c|}
        \hline
        高度/cm & [10, 20) & [20, 30) & [30, 40) & [40, 50) & [50,60] \\
        \cline{1-6}
        频数 & 20 & 30 & 80 & 40 & 30 \\
        \cline{1-6}
        \hline
        \end{tabular}
        \end{center}
        
        则该植物一年生长在$[30, 40)$ 内的频率是\blkda{$0.40$}
        \item 某产品的设计长度为 20cm，规定误差不超过 0.5cm 为合格品. 今对一批产品进行测量, 测量结果如下表:

        \begin{center}
        \begin{tabular}{|c|c|c|c|}
        \hline
        长度(cm) & 19.5 以下 & 19.5\textasciitilde 20.5 & 20.5 以上 \\
        \cline{1-4}
        件数 & 5 & 68 & 7 \\
        \cline{1-4}
        \hline
        \end{tabular}
        \end{center}则这批产品的不合格率为\blkda{$\dfrac{3}{20}$}
        \begin{jieda}
            因为某产品的设计长度为 ${20}{{cm}}$ ，规定误差不超过 ${0.5}{{cm}}$ 为合格品，且 ${20} - {0.5} = {19.5},{20} +  {0.5} = {20.5}$ ,
            
            所以长度在 ${19.5} \sim  {20.5}$ 的为合格品,
            
            所以这批产品的不合格率为 $\dfrac{5 + 7}{5 + {68} + 7} = \dfrac{12}{80} = \dfrac{3}{20}$
        \end{jieda}
        
        \item 已知一个容量为 20 的样本，其数据具体如下:
        
        $\begin{array}{llllllllll} {10} & 8 & 6 & {10} & {13} & 8 & {10} & {12} & {11} & 7 \end{array}$
        
        $\begin{array}{llllllllll} 8 & 9 & {11} & 9 & {12} & 9 & {10} & {11} & {12} & {11} \end{array}$
        
        那么频率为 0.4 的范围是\blkda{$9.5 \sim 11.5$}
        
        \item 在第 112 个“三八”国际妇女节到来之际，我们年级组给每位女教师送了火红的玫瑰和热情洋溢的贺卡，其中我最喜欢的一句话英文翻译是:“You stand on the podium as an oath and tell the ordinary as a fairly tale."统计如下表:
        
        \begin{center}
        \begin{tabular}{|c|c|c|c|c|c|c|c|c|c|c|c|c|c|c|c|c|}
        \hline
        字母 & $y$ & $o$ & $u$ & $s$ & $a$ & $p$ & $d$ & $e$ & $t$ & $f$ & $l$ & $i$ & $n$ & $m$ & $r$ & $h$ \\
        \cline{1-17}
        频数 & 3 & 5 & 2 & 3 & 10 & 1 & 4 & 4 & 6 & 1 & 4 & 3 & 5 & 1 & 3 & 3 \\
        \cline{1-17}
        \hline
        \end{tabular}
        \end{center}
        
        累计频数为 58 ，则字母“ $a$ ”出现的频率约为\dotfill{(\kh{B})}
        \xx{$10$}{$0.172$}{$0.123$}{$0.153$}
        \begin{jieda}
            由表格数据可知: 字母“ $a$ ”对应的频数为 10, $\therefore$ 字母“ $a$ ”出现的频率约为 $\dfrac{10}{58} \approx  {0.172}$.
        \end{jieda}
    \end{enumerate}
    \item 辨析概率与频率的关系
    \begin{enumerate}
        \item 下列说法一定正确的是\dotfill{(\kh{B})}
        \xx{一名篮球运动员，号称“百发百中”，若罚球三次，不会出现三投都不中的情况}{随机事件发生的概率与试验次数无关}{若买彩票中奖的概率为万分之一，则买一万元的彩票一定会中奖一元}{一个骰子掷一次得到 2 的概率是 $\dfrac{1}{6}$ ，则掷 6 次一定会出现一次 2}
        
        \item 下列说法正确的是\dotfill{(\kh{D})}
        \xx{一个人打靶，打了 10 发子弹，有 7 发子弹中靶，因此这个人中靶的概率是 $\dfrac{7}{10}$}{一个同学做掷硬币试验，掷了 6 次，一定有 3 次正面向上}{某地发行彩票，其回报率为 47\%，有人花了 100 元钱买彩票，一定会有 47 元的回报}{大量试验后，可以用频率近似估计概率}
        
        \item 抛一枚硬币的试验中, 下列对“伯努利大数定律”的理解正确的是\dotfill{(\kh{B})}
        \xx{大量的试验中，出现正面的频率为 0.5}{不管试验多少次，出现正面的概率始终为 0.5}{试验次数增大，出现正面的经验概率为 0.5}{试验次数每增加一次，下一次出现正面的频率一定比它前一次更接近于 0.5}

        
        \item 在调查运动员是否服用过兴奋剂的时候，给出两个问题作答，无关紧要的问题是:“你的身份证号码的尾数是奇数吗？”敏感的问题是:“你服用过兴奋剂吗？”然后要求被调查的运动员掷一枚硬币，如果出现正面，就回答第一个问题，否则回答第二个问题. 由于回答哪一个问题只有被测试者自己知道，所以应答者一般乐意如实地回答问题. 如我们把这种方法用于 300 个被调查的运动员，得到 80 个“是”的回答，则这群人中服用过兴奋剂的百分率大约为\dotfill{(\kh{B})}
        \xx{$4.33\%$}{$3.33\%$}{$3.44\%$}{$4.44\%$}
        \begin{jieda}
            因为抛硬币出现正面朝上的概率为 $\dfrac{1}{2}$ . 大约有 150 人回答第一个问题，又身份证号码的尾数是奇数或偶数是等可能的.

            在回答第一个问题的 150 人中大约有一半人，即 75 人回答了“是”，
            
            共有 80 个“是”的回答，故回答服用过兴奋剂的人有 5 人，
            
            因此我们估计这群人中，服用过兴奋剂的百分率大约为 $\dfrac{5}{150} \approx  {3.33}\%$ .
        \end{jieda}
        
        \item 下列说法正确的有\blkda{\ding{172}}
        \begin{dingautolist}{172}
            \item 随机事件 $A$ 的概率是频率的稳定值，频率是概率的估计值
            \item 某人打靶，射击 10 次，击中 7 次，那么此人中靶的概率 0.7
            \item 一位同学做掷硬币试验，掷 6 次，一定有 3 次正面朝上
            \item 某地发行福利彩票，回报率为 47\%，有人花了 100 元钱买彩票，一定会有 47 元的回报
        \end{dingautolist}

        \item 下列说法正确的是\dotfill{(\kh{A})}
        \xx{随着试验次数的增大, 随机事件发生的频率会逐渐稳定于该随机事件发生的概率}{某种福利彩票的中奖概率为 $\dfrac{1}{1000}$ ，买 1000 张这种彩票一定能中奖}{连续 100 次掷一枚硬币，结果出现了 49 次反面，则掷一枚硬币出现反面的概率为 $\dfrac{49}{100}$}{某市气象台预报“明天本市降水概率为 70\%”，指的是:该市气象台专家中，有 70\%认为明天会降水， 30\%认为明天不会降水}
    \end{enumerate}
    
    \item 实际问题的概率解释
    \begin{enumerate}
        \item 在天气预报中，有“降水概率预报”，例如预报“明天降水的概率为 80\%"，这是指\dotfill{(\kh{B})}
        \xx{明天该地区有 80\%的地方降水，有 20\%的地方不降水}{明天该地区降水的可能性为 80\%}{气象台的专家中有 80\%的人认为会降水，另外有 20\%的专家认为不降水}{明天该地区有 80\%的时间降水，其他时间不降水}
        
        \item 气象台预报“本市明天降雨概率是 70\%”，下列说法正确的是\dotfill{(\kh{C})}
        \xx{本市明天将有 70\%的地区降雨}{本市有天将有 70\%的时间降雨}{明天出行不带雨具淋雨的可能性很大}{明天出行不带雨具肯定要淋雨}
        
        \item 某种彩票的中奖概率为 $\dfrac{1}{100000}$ ，则以下理解正确的是\dotfill{(\kh{B})}
        \xx{购买这种彩票 100000 张，一定能中奖一次}{购买这种彩票 100000 张，可能一次也没中奖}{购买这种彩票 1 张，一定不能中奖}{购买这种彩票 100000 张，至少能中奖一次}
        
        \item 已知某厂生产的某批产品的合格率为 90\%，现从该批次产品中抽出 100 件产品检查, 则下列说法正确的是\dotfill{(\kh{D})}
        \xx{合格产品少于 90 件}{合格产品多于 90 件}{合格产品正好是 90 件}{合格产品可能是 90 件}
        
        \item 老师讲一道数学题，李峰能听懂的概率是 0.8 ，是指\dotfill{(\kh{C})}
        \xx{老师每讲一题，该题有 80\%的部分能听懂，20\%的部分听不懂}{老师在讲的 10 道题中，李峰能听懂 8 道}{李峰听懂老师所讲这道题的可能性为 80\%}{以上解释都不对}
        
        \item 根据某市疾控中心的健康监测，该市在校中学生的近视率约为 78.7\%.某眼镜厂商要到中学给近视学生配送滴眼液，每人一瓶，已知该校学生总数为 600 人，则眼镜商应带滴眼液的瓶数为\dotfill{(\kh{C})}
        \xx{600}{787}{不少于 473}{不多于 473}
        \begin{jieda}
            依题意, 该市在校中学生的近视率约为 78.7\%.
            
            故 600 人中大约有 ${600} \times  {78.7}\%  \approx  {472}$，故眼镜商应带滴眼液的瓶数应不少于 473 瓶
        \end{jieda}
        \item 下列说法正确的是\dotfill{(\kh{B})}
        \xx{某厂一批产品的次品率为 $\dfrac{1}{10}$ ，则任意抽取其中 10 件产品一定会发现一件次品}{掷一枚硬币，连续出现 5 次正面向上，第六次出现反面向上的概率与正面向上的概率仍然都为 0.5}{某医院治疗一种疾病的治愈率为 10\%，那么前 9 个病人都没有治愈，第 10 个人就一定能治愈}{气象部门预报明天下雨的概率是 90\%，说明明天该地区 90\%的地方要下雨，其余 10\%的地方不会下雨}
        \begin{jieda}
            A. 产品的次品率是通过大量的产品通过实验得到的数据，题目中的产品个数很少，故不正确； 
            
            B. 掷硬币正面或反面朝上的概率是通过大量实验得到的准确的值，和实验次数无关，故正确；
            
            C.解释同 A 选项, 也不正确;
            
            D. 事件的概率是大量实验后得到的结果, 是准确的值, 和实验次数无关, 但是 D 选项的说法体现的不是概率的概念, 故不正确.
        \end{jieda}
        \item 已知使用一剂某种药物治愈某种疾病的概率为 90\%，则下列说法正确的是\dotfill{(\kh{C})}
        \xx{如果有 100 个这种病人各使用一剂这样的药物，那么有 90 人会被治愈}{如果一个患有这种疾病的病人使用两剂这样的药物就一定会被治愈}{使用一剂这种药物治愈这种疾病的可能性是 90\%}{以上说法都不对}
    \end{enumerate}
    
    \item 事件的关系和运算
    \begin{enumerate}
        \item 掷一枚骰子，设事件 $A =$ \{出现的点数不小于 5\}， $B =$ \{出现的点数为偶数 $\}$ ,则事件 $A$ 与事件 $B$ 的关系是\dotfill{(\kh{B})}
        \xx{$A \subseteq  B$}{$A \cap  B = \{$ 出现的点数为 6 $\}$}{事件 $A$ 与 $B$ 互斥}{事件 $A$ 与 $B$ 是对立事件}
        \begin{jieda}
            $A = \{$ 出现的点数不小于 5 $\}  = \{$ 出现的点数为 5,6 \}， $B = \{$ 出现的点数为偶数 $\}  = \{$ 出现的点数为 2,4,6 \}，则 $A \cap  B =$ \{出现的点数为 6\},故 B 正确, A 错误;

            因为事件 $A$ 与事件 $B$ 可以同时发生,故事件 $A$ 与 $B$ 不是互斥事件,也不是对立事件,故 ${C},{D}$ 错误, 故选: B.
        \end{jieda}

        
        \item 抛掷一枚质地均匀的骰子，设事件 $M =$ “点数不大于 2”，事件 $N =$ “点数大于 1”，则下列结论中正确的是\dotfill{(\kh{D})}
        \xx{$M$ 是不可能事件}{$N$ 是必然事件}{$M \cap  N$ 是不可能事件}{$M \cup  N$ 是必然事件}

        
        \item 从装有两个红球和三个黑球的口袋里任取两个球，那么互斥而不对立的两个事件是\dotfill{(\kh{A})}
        \xx{“恰好有一个黑球”与“恰好有两个黑球”}{“至少有一个黑球”与“都是红球”}{“至少有一个黑球”与“至少有一个红球”}{“至少有一个黑球”与“都是黑球”}
        \begin{jieda}
            对于 A，恰好有一个黑球的事件与恰好有两个黑球的事件不能同时发生，可以同时不发生， 因此“恰好有一个黑球”与“恰好有两个黑球”是互斥而不对立的两个事件，A 是；
            
            对于 ${B}$ ,至少有一个黑球的事件与都是红球的事件是对立事件, ${B}$ 不是;
            
            对于 ${C}$ ,至少有一个黑球的事件与至少有一个红球的事件可以同时发生,不互斥, ${C}$ 不是;
            
            对于 ${D}$ ,至少有一个黑球的事件与都是黑球的事件可以同时发生,不互斥, ${D}$ 不是.
        \end{jieda}
        
        \item 在一次随机试验中，彼此互斥的事件 $A,B,C,D$ 的概率分别是0.2,0.2,0.3,0.3,则下列说法正确的是\dotfill{(\kh{D})}
        \xx{$A \cup B$ 与 $C$ 是互斥事件，也是对立事件}{$B \cup C$ 与 $D$ 是互斥事件,也是对立事件}{$A \cup C$ 与 $B \cup D$ 是互斥事件，但不是对立事件}{$A$ 与 $B \cup C \cup D$ 是互斥事件,也是对立事件}
        \begin{jieda}
            由于 $A,B,C,D$ 彼此互斥且 ${0.2} + {0.2} + {0.3} + {0.3} = 1$ ,则 $A + B + C + D$ 是一个必然事件，故有：
            
            任何一个事件与其余 3 个事件的和事件必然是对立事件;
            
            任何两个事件的和事件与其余两个事件的和事件也是对立事件.
            
            所以 $A \neq  B$ 与 $C$ 是互斥事件，但不是对立事件；
            
            $B \cup C$ 与 $D$ 是互斥事件，但不是对立事件；
            
            $A \cup C$ 与 $B \cup D$ 是互斥事件，也是对立事件；
            
            $A$ 与 $B \cup C \cup D$ 是互斥事件,也是对立事件.
        \end{jieda}
        
        \item 甲、乙两人对同一个靶各射击一次，设事件 $A =$ “甲击中靶”，事件 $B =$ “乙击中靶”，事件 $E =$ “靶未被击中”，事件 $F =$ “靶被击中”，事件 $G =$ “恰一人击中靶”，对下列关系式 ( $\bar{A}$ 表示 $A$ 的对立事件, $\bar{B}$ 表示 $B$ 的对立事件)正确的是\blkda{\ding{172}\ding{174}\ding{176}\ding{177}}
        \begin{dingautolist}{172}
            \item $E = \overline{A} \cap \overline{B}$
            \item $F = {A \cap B}$
            \item $F = A \cup B$
            \item $G = A \cup B$
            \item $G = \bar{A}B \cup A\bar{B}$
            \item $P\left( F\right)  = 1 - P\left( E\right)$
            \item $P\left( F\right)  = P\left( A\right)  + P\left( B\right)$
        \end{dingautolist}
    \end{enumerate}
    
    \item 概率的基本性质及其应用
    \begin{enumerate}
        \item 设 $A,B$ 为两个互斥事件，且 $P\left( A\right)  > 0,P\left( B\right)  > 0$ ，则下列各式一定正确的是\dotfill{(\kh{B})}
        \xx{$P\left( {AB}\right)  = P\left( A\right) P\left( B\right)$}{$P\left( {A \cup  B}\right)  = P\left( A\right)  + P\left( B\right)$}{$P\left( {AB}\right)  = P\left( A\right)  + P\left( B\right)$}{$P\left( {A \cup  B}\right)  = P\left( A\right) P\left( B\right)$}
        \begin{jieda}
            因为 $A,B$ 为两个互斥事件, $P\left( A\right)  > 0,P\left( B\right)  > 0$，所以 $A \cap  B = \varnothing$ ,即 $P\left( {AB}\right)  = 0$ ,且 $P\left( {A \cup  B}\right)  = P\left( A\right)  + P\left( B\right)$ .
        \end{jieda}

        \item 若事件 $A,B$ 为两个互斥事件，且 $P\left( A\right)  > 0,P\left( B\right)  > 0$ ，有以下四个结论，其中正确的结论是\blkda{\ding{172}\ding{174}\ding{175}}
        \begin{dingautolist}{172}
            \item $P\left( {AB}\right)  = 0$
            \item $P\left( {\bar{A}B}\right)  = \left\lbrack  {1 - P\left( A\right) }\right\rbrack  P\left( B\right)$
            \item $P\left( {\overset{ - }{A} \cup  \overset{ - }{B}}\right)  = 1$
            \item $P\left( {A \cup  B}\right)  = P\left( A\right)  + P\left( B\right)$
        \end{dingautolist}
        
        \item 已知事件 $A\text{ 、 }B\text{ 、 }C$ 两两互斥,若 $P\left( A\right)  = \dfrac{1}{6},P\left( C\right)  = \dfrac{1}{3},P(A \cup  B) = \dfrac{11}{30}$ ,则 $P\left( {B \cup  C}\right)  =$\blkda{$\dfrac{8}{15}$}
        \begin{jieda}
            因为$A\text{ 、 }B\text{ 、 }C$ 两两互斥，所以$P(A \cup B) = P(A) + P(B)$，因此$P(B) = \dfrac{1}{5}$

            故$P(B \cup C) = P(B) + P(C) = \dfrac{8}{15}$
        \end{jieda}
    
        \item 设 $A,B$ 是一个随机试验中的两个事件,且 $P\left( A\right)  = \dfrac{1}{2},P\left( \bar{B}\right)  = \dfrac{7}{12}$ , $P\left( {\bar{A}B \cup A\bar{B}}\right)  = \dfrac{1}{4}$ 则 $P\left( {A \cup B}\right)  =$\blkda{$\dfrac{7}{12}$}
        \begin{jieda}
            $P\left( A\right)  = \dfrac{1}{2},P\left( \bar{B}\right)  = \dfrac{7}{12}$ ,故 $P\left( \bar{A}\right)  = \dfrac{1}{2},P\left( B\right)  = \dfrac{5}{12}$
        
            因为 $\bar{A}B$ 与 $A\bar{B}$ 为互斥事件,故 $P\left( {\bar{A}B \cup  A\bar{B}}\right)  = P\left( {\bar{A}B}\right)  + P\left( {A\bar{B}}\right)  = \dfrac{1}{4}$

            又因为$\left\{ \begin{array}{l}
                 P\left( {\bar{A}B}\right)  + P\left( {AB}\right)  = P\left( B\right)  \\
                 P\left( {A\bar{B}}\right)  + P\left( {AB}\right)  = P\left( A\right) 
            \end{array} \right.$，所以有 $P\left( B\right)  - P\left( {AB}\right)  + P\left( A\right)  - P\left( {AB}\right)  = \dfrac{5}{12} + \dfrac{1}{2} - {2P}\left( {AB}\right)  = \dfrac{1}{4}$，故 $P\left( {AB}\right)  = \dfrac{1}{3}$

            故 $P\left( {A \cup B}\right)  = P\left( A\right)  + P\left( B\right)  - P\left( {AB}\right)  = \dfrac{1}{2} + \dfrac{5}{12} - \dfrac{1}{3} = \dfrac{7}{12}$.
        \end{jieda}
    \end{enumerate}
    
    \item 独立事件的判断
    \begin{enumerate}
        \item 设事件 $A$ 与事件 $B$ 满足: $P\left( A\right)  = \dfrac{2}{3},P\left( B\right)  = \dfrac{3}{4},P\left( {AB}\right)  = \dfrac{1}{2}$ ， 则下列说法正确的是\dotfill{(\kh{C})}
        \xx{事件 $\bar{A}$ 与事件 $B$ 不是相互独立事件}{事件 $A$ 与事件 $\bar{B}$ 不是相互独立事件}{事件 $A$ 与事件 $B$ 是相互独立事件}{事件 $\bar{A}$ 与事件 $\bar{B}$ 不是相互独立事件}
        \begin{jieda}
            因为$P(AB) = P(A)P(B)$，所以$A$和$B$独立
        \end{jieda}

        
        \item 已知 $A,B$ 是两个概率大于 0 的随机事件，则下列说法错误的是\dotfill{(\kh{D})}
        \xx{若 $A,B$ 是对立事件，则 $A,B$ 是互斥事件}{若事件 $A,B$ 相互独立，则 $A$ 与 $\bar{B}$ 也相互独立}{若事件 $A,B$ 相互独立，则 $A$ 与 $B$ 不互斥}{若事件 $A,B$ 互斥,则 $A$ 与 $B$ 相互独立}
        \begin{jieda}
            A. 两个事件是对立事件，则一定是互斥事件，故 A 正确；
            
            B. 若事件 $A,B$ 相互独立，则 $A$ 与 $\bar{B}$ 也相互独立，故 B 正确；
            
            C. 若事件 $A,B$ 相互独立，则 $A$ 与 $B$ 可以同时发生，不互斥，故 C 正确；
            
            D. 若事件 $A,B$ 互斥，则 $A$ 与 $B$ 不能同时发生，即事件 $A\left( B\right)$ 是否发生，对另一个事件 $B\left( A\right)$ 是有影响的，所以两个事件不相互独立，故 D 错误.
        \end{jieda}

        \item 下列事件 $A,B$ 是相互独立事件的是\dotfill{(\kh{A})}
        \xx{一枚硬币掷两次， $A$ 表示“第一次为正面”， $B$ 表示“第二次为反面”}{袋中有 2 个白球,2 个黑球,不放回地摸球两次,每次摸一球, $A$ 表示“第一次摸到白球”, $B$ 表示“第二次摸到白球”}{掷一枚骰子, $A$ 表示“出现点数为奇数”, $B$ 表示“出现点数为偶数”}{对于同一枚灯泡，$A$ 表示“一个灯泡能用 1000 小时”， $B$ 表示“一个灯泡能用 2000 小时”}
        
        \item 有 4 个大小质地相同的小球，分别标有数字1,2,3,4，从中不放回的随机抽取两次，每次取一个球. 甲表示事件“第一次取出的球的数字是奇数”，乙表示事件“第一次取出的球的数字是偶数”，丙表示事件“两次取出的球的数字之和为 4”，丁表示事件“两次取出的球的数字之和为 5”， 则\dotfill{(\kh{C})}
        \xx{甲和乙相互独立}{甲和丙相互独立}{甲和丁相互独立}{丁和丙相互独立}
        \begin{jieda}
            甲和乙是对立事件，不独立；丙和丁是互斥事件，不独立
            
            丙发生甲必发生，故甲和丙不独立；$P(\text{甲}) = \dfrac{1}{2}, P(\text{丙}) = \dfrac{1}{6}$，$P(\text{甲}\cap\text{丙}) = \dfrac{1}{6}$
        \end{jieda}

        \item 已知甲盒中有 3 个大小和质地相同的小球，标号为 1,3,4，乙盒中有 3 个大小和质地相同的小球,标号为3,4,6,现从甲、乙两盒中分别随机摸出 1 个小球,记事件 $A =$ “摸到的两个小球标号相同”，事件 $B =$ “摸到的两个小球标号之和为奇数”，则\dotfill{(\kh{D})}
        \xx{事件 $A$ 和 $B$ 相等}{事件 $A$ 和 $B$ 互相对立}{事件 $A$ 和 $B$ 相互独立}{事件 $A$ 和 $B$ 互斥}

        \item 一个不透明的盒子中装有大小和质地都相同的编号分别为 1, 2, 3, 4, 5, 6 的 6 个小球,从中任意摸出两个球. 设事件 ${A}_{1} =$ “摸出的两个球的编号之和不超过 6”,事件 ${A}_{2} =$ “摸出的两个球的编号都大于 3 ”，事件 ${A}_{3} =$ “摸出的两个球中有编号为 4 的球”，则\dotfill{(\kh{D})}
        \xx{事件 ${A}_{1}$ 与事件 ${A}_{2}$ 是相互独立事件}{事件 ${A}_{1}$ 与事件 ${A}_{3}$ 是对立事件}{事件 ${A}_{1} \cup  {A}_{2}$ 与事件 ${A}_{3}$ 是互斥事件}{事件 ${A}_{1} \cap  {A}_{3}$ 与事件 ${A}_{2} \cap  {A}_{3}$ 是互斥事件}
        \begin{jieda}
            事件 ${A}_{1}$ 与事件 ${A}_{2}$互斥，不独立.

            若摸出两个球分别为$(1,4)$，则事件 ${A}_{1}$ 与事件 ${A}_{3}$均发生，故两者不是对立事件，同理事件 ${A}_{1} \cup  {A}_{2}$ 与事件 ${A}_{3}$ 不是互斥事件
        \end{jieda}
        
        \item 袋子中有 5 个相同的球，分别标有数字1,2,3,4,5，从中随机取出两个球，设事件 $A =$ “取出的球的数字之积为奇数”，事件 $B =$ “取出的球的数字之积为偶数”，事件 $C =$ "取出的球的数字之和为偶数"，事件 $D =$ "取出的球的数字之和大于 5"，则下列说法错误的是\dotfill{(\kh{C})}
        \xx{事件 $A$ 与 $B$ 是互斥事件}{事件 $A$ 与 $B$ 是对立事件}{事件 $C$ 与 $D$ 相互独立}{事件 $C$ 与 $D$ 不是互斥事件}
        \begin{jieda}
            袋子中有 5 个相同的球,分别标有数字1,2,3,4,5,
            
            从中随机取出两个球的试验样本空间包含的样本点为:(1,2)，(1,3)，(1,4)，(1,5)，(2,3)，(2,4)， $\left( {2,5}\right) ,\left( {3,4}\right) ,\left( {3,5}\right) ,\left( {4,5}\right)$ 共 10 个,
            
            其中事件 $A$ 包含的样本点为:(1,3)，(1,5)，(3,5)共 3 个，故 $P\left( A\right)  = \dfrac{3}{10}$
            
            事件 $B$ 包含的样本点为:(1,2)，(1,4)，(2,3)，(2,4)，(2,5)，(3,4)，(4,5)共 7 个，故 $P\left( B\right)  = \dfrac{7}{10}$
            
            事件 $C$ 包含的样本点为:(1,3)，(1,5)，(2,4)，(3,5)共 4 个，故 $P\left( C\right)  = \dfrac{4}{10} = \dfrac{2}{5}$
            
            事件 $D$ 包含的样本为:(1,5)，(2,4)，(2,5)，(3,4)，(3,5)，(4,5)共 6 个，故 $P\left( D\right)  = \dfrac{6}{10} = \dfrac{3}{5}$ ， 因为事件 $A \cap  B = \varnothing ,A \cup  B = \Omega$ ,故事件 $A$ 与 $B$ 互斥且对立,故 ${A},{B}$ 正确；
            
            因为 $P\left( {CD}\right)  = \dfrac{3}{10} \neq  \dfrac{2}{5} \times  \dfrac{3}{5} = P\left( C\right) P\left( D\right)$ ,所以 $C$ 与 $D$ 不相互独立,故 ${C}$ 错误.
            
            因为 $C \cap  D = \{ \left( {1,5}\right) ,\left( {2,4}\right) ,\left( {3,5}\right) \}$ ,所以 $C$ 与 $D$ 不互斥,故 ${D}$ 正确.
        \end{jieda}
        
        \item 假定生男生女是等可能的，设事件A:一个家庭中既有男孩又有女孩； 事件 $B$ : 一个家庭中最多有一个女孩.针对下列两种情形: \ding{172} 家庭中有 2 个小孩； \ding{173} 家庭中有 3 个小孩，下面说法正确是\dotfill{(\kh{B})}
        \xx{\ding{172} 中事件 $A$ 与事件 $B$ 相互独立、 \ding{173} 中的事件 $A$ 与事件 $B$ 相互独立}{\ding{172} 中事件 $A$ 与事件 $B$ 不相互独立、 \ding{173} 中的事件 $A$ 与事件 $B$ 相互独立}{\ding{172} 中事件 $A$ 与事件 $B$ 相互独立、 \ding{173} 中的事件 $A$ 与事件 $B$ 不相互独立}{\ding{172} 中事件 $A$ 与事件 $B$ 不相互独立、 \ding{173} 中的事件 $A$ 与事件 $B$ 不相互独立}
        \begin{jieda}
            \begin{dingautolist}{172}
                \item $P(A) = \dfrac{1}{2}, P(B) = \dfrac{3}{4}, P(AB) = \dfrac{1}{2}$

                \item $P(A) = \dfrac{3}{4}, P(B) = \dfrac{1}{2}, P(AB) = \dfrac{3}{8}$
            \end{dingautolist}

            \ps 无论知道不知道$B$，$A$发生的条件概率与无条件概率一致.
        \end{jieda}
    \end{enumerate}
    
    \item 相互独立事件的概率
    \begin{enumerate}
        \item 已知事件 $A,B$ 相互独立,且 $P\left( A\right)  = \dfrac{1}{3},P\left( B\right)  = \dfrac{2}{3}$ ,则 $P\left( {A\bar{B}}\right)  =$ \blkda{$\dfrac{1}{9}$}
        \begin{jieda}
            因为事件 $A,B$ 是相互独立事件,所以 $A$ 与 $\bar{B}$ 相互独立， $P\left( \bar{B}\right)  = 1 - P\left( B\right)  = \dfrac{1}{3}$ ，则 $P\left( {A\bar{B}}\right)  = P\left( A\right) P\left( \bar{B}\right)  = \dfrac{1}{3} \times  \dfrac{1}{3} = \dfrac{1}{9}$ .
        \end{jieda}
        
        \item 如图，三个元件 $T_1$ ， $T_2$ ，$T_3$，正常工作的概率为$\dfrac{1}{3}$，且是相互独立的, 将它们接入电路中, 则电路不发生故障的概率是\blkda{$\dfrac{5}{27}$}

        \begin{center}
        \includegraphics[width=0.2\textwidth]{bodyxb/src/bo_d67fmev7aajc739ac51g_3_143_747_325_175_0.jpg}
        \end{center}
        \begin{jieda}
            $P = P(T_1) \times (1-P(\overline{T_2})P(\overline{T_3})) = \dfrac{1}{3} \times \left(1-\dfrac{2}{3}\times \dfrac{2}{3}\right) = \dfrac{5}{27}$
        \end{jieda}
        
         \item 如图,用 $A,B,C,D$ 四个不同的元件连接成一个工作系统, 当元件 ${A}$ 正常工作,且 $B,C,D$ 三个元件中至少有一个正常工作时,该系统正常工作. 已知元件 $A$ 正常工作的概率为 $\dfrac{1}{2}$ ，元件 $B,C,D$ 正常工作的概率均为 $\dfrac{1}{3}$ ，且这四个元件是否正常工作相互独立，则该系统正常工作的概率为\blkda{$\dfrac{19}{54}$}
        
        \begin{center}
        \includegraphics[width=0.3\textwidth]{bodyxb/src/bo_d67fmev7aajc739ac51g_4_142_190_378_227_0.jpg}
        \end{center}
        \begin{jieda}
            由题可知,元件 $B,C,D$ 均不正常工作的概率为 ${\left( 1 - \dfrac{1}{3}\right) }^{3} = \dfrac{8}{27}$，则元件 $B,C,D$ 中至少有一个正常工作的概率为 $1 - \dfrac{8}{27} = \dfrac{19}{27}$
            
            从而该系统正常工作的概率为 $\dfrac{1}{2} \times  \dfrac{19}{27} = \dfrac{19}{54}$
        \end{jieda}
    
        \item 设甲、乙两人每次投进篮球的概率分别为 $\dfrac{1}{3}$ 与 $\dfrac{2}{3}$ ，两人约定如下投篮:每次由一人投篮，若投进，下一次由另一人投篮；若没有投进，则继续投篮，则前 4 次中甲恰好投篮 3 次的概率为\blkda{$\dfrac{20}{27}$}
        \begin{jieda}
            若第一次甲先投篮,则前 4 次中甲恰好投篮 3 次的概率为: $\dfrac{1}{3} \times  \dfrac{2}{3} \times  \dfrac{2}{3} \times  1 + \dfrac{2}{3} \times  \dfrac{1}{3} \times  \dfrac{2}{3} \times  1 + \dfrac{2}{3} \times  \dfrac{2}{3} \times  \dfrac{1}{3} \times  1 = \dfrac{12}{27}$
            
            若第一次乙先投篮，则前 4 次中甲恰好投篮 3 次的概率为: $\dfrac{2}{3} \times  \dfrac{2}{3} \times  \dfrac{2}{3} \times  1 = \dfrac{8}{27}$
            
            故前 4 次中甲恰好投篮 3 次的概率为 $\dfrac{12}{27} + \dfrac{8}{27} = \dfrac{20}{27}$
        \end{jieda}
    \end{enumerate}
    \item 独立事件的实际应用
    \begin{enumerate}
        \item 已知甲、乙两人射击的命中率分别是 0.4 和 0.7 . 现二人同时向同一猎物射击, 发现猎物只中一枪, 则甲、乙分配猎物的比例应该是\blkda{$2:7$}
        \begin{jieda}
            因为甲、乙两人射击的命中率分别是 0.4 和 0.7，现二人同时向同一猎物射击，发现猎物只中一枪，只有甲打中猎物的概率为 ${0.4} \times  {0.3} = {0.12}$ ,只有乙打中猎物的概率为 ${0.6} \times  {0.7} = {0.42}$
            
            所以,甲、乙分配猎物的比例应该是 ${0.12} : {0.42} = 2 : 7$ .
        \end{jieda}
        
        \item 概率论起源于博弈游戏 17 世纪，曾有一个“赌金分配”的问题:博弈水平相当的甲、乙两人进行博弈游戏每局比赛都能分出胜负，没有平局. 双方约定，各出赌金 180 枚金币， 先赢 3 局者可获得全部赎金; 但比赛中途因故终止了, 此时甲赢了 2 局, 乙赢了 1 局.问这 360 枚金币的赌金该如何分配？数学家费马和帕斯卡都用了现在称之为“概率”的知识，合理地给出了赌金分配方案.该分配方案是\dotfill{(\kh{D})}
        \xx{甲 180 枚，乙 180 枚}{甲 288 枚，乙 72 枚}{甲 240 枚，乙 120 枚}{甲 270 枚，乙 90 枚}
        \begin{jieda}
            根据题意,甲、乙两人每局获胜的概率均为 $\dfrac{1}{2}$
            
            假设两人继续进行比赛，甲获取 360 枚金币有:第四局甲赢，或第四局甲输，第五局甲赢，故概率为 ${p}_{1} = \dfrac{1}{2} + \dfrac{1}{2} \times  \dfrac{1}{2} = \dfrac{3}{4}$
            
            乙获取 360 枚金币有:第四、五局乙都赢，故概率为 ${p}_{2} = \dfrac{7}{2} \times  \dfrac{1}{2} = \dfrac{1}{4}$ ,
            
            则甲应该获得枚金币 $\dfrac{3}{4} \times  {360} = {270}$ ，乙应该获得枚金币 $\dfrac{1}{4} \times  {360} = {90}$
        \end{jieda}
        
        \item 2020 年 1 月，教育部发布《关于在部分高校开展基础学科招生改革试点工作的意见》(简称“强基计划”)，明确从 2020 年起强基计划取代原有的高校自主招生方式. 某高校笔试环节要求考生参加三个科目考核，考生通过三个科目的笔试考核才能进入面试环节. 考生甲通过三个科目的笔试考核的概率分别为 $\dfrac{1}{2},\dfrac{2}{3},\dfrac{3}{4}$ ，且每个科目考核相互独立，则甲顺利进入面试环节的概率为\blkda{$\dfrac{1}{4}$}
        \begin{jieda}
            记甲通过三个科目的笔试考核分别为事件 $A,B,C$，显然 $A,B,C$ 为相互独立事件,则事件“甲通过三个科目的笔试考核”相当于事件 ${ABC}$ 
            
            $P\left( {ABC}\right)  = P\left( A\right) P\left( B\right) P\left( C\right)  = \dfrac{1}{2} \times  \dfrac{2}{3} \times  \dfrac{3}{4} = \dfrac{1}{4}$
        \end{jieda}
        
        \item 为庆祝我国第 39 个教师节，某校举办教师联谊会，甲、乙两名数学老师组成“几何队”参加“成语猜猜猜”比赛，每轮比赛由甲、乙两人各猜一个成语，已知甲每轮猜对的概率为 $\dfrac{4}{5}$ ，乙每轮猜对的概率为 $\dfrac{3}{4}$ 在每轮比赛中，甲和乙猜对与否互不影响，则 “几何队”在一轮比赛中至少猜对一个成语的概率为\blkda{$\dfrac{19}{20}$}
        \begin{jieda}
            设事件 $A =$ “甲猜对”， $B =$ “乙猜对”， $C =$ “几何队至少猜对一个成语”，所以 $P\left( A\right)  = \dfrac{4}{5},P\left( B\right)  = \dfrac{3}{4}$ 则 $P\left( \bar{A}\right)  = \dfrac{1}{5},P\left( \bar{B}\right)  = \dfrac{1}{4}$
            
            由题意知,事件 $A,B$ 相互独立,则 $\bar{A}$ 与 $B,A$ 与 $\bar{B},\bar{A}$ 与 $\bar{B}$ 也相互独立,
            
            法一: $C = \left( {\bar{A}B}\right)  \cup  \left( {A\bar{B}}\right)  \cup  \left( {AB}\right)$ ,且 $\bar{A}B,A\bar{B},{AB}$ 两两互互斥,
            
            则 $P\left( C\right)  = P\left( {\bar{A}B}\right)  + P\left( {A\bar{B}}\right)  + P\left( {AB}\right)  = P\left( \bar{A}\right) P\left( B\right)  + P\left( A\right) P\left( \bar{B}\right)  + P\left( A\right) P\left( B\right)$
            
            $= \dfrac{1}{5} \times  \dfrac{3}{4} + \dfrac{4}{5} \times  \dfrac{1}{4} + \dfrac{4}{5} \times  \dfrac{3}{4} = \dfrac{19}{20}$
            
            法二: 事件 $C$ 的对立事件 $\bar{C} =$ “几何队一个成语也没有猜对”,即 $\bar{C} = \overline{AB}$ ,
            
            则 $P\left( C\right)  = 1 - P\left( \bar{C}\right)  = 1 - P\left( {\bar{A}\bar{B}}\right)  = 1 - P\left( \bar{A}\right) P\left( \bar{B}\right)  = 1 - \dfrac{1}{5} \times  \dfrac{1}{4} = \dfrac{19}{20}$ .
        \end{jieda}
    \end{enumerate}
    
    \item 互斥事件、事件的相互独立性的综合应用
    \begin{enumerate}
        \item 甲、乙两位队员进行某种球类对抗赛，每局依次轮流发球， 连续赢 2 个球者获胜，通过分析甲、乙过去对抗赛的数据知，甲发球甲赢的概率为 $\dfrac{2}{3}$ ，乙发球甲赢的概率为 $\dfrac{1}{4}$ 不同球的结果互不影响,已知某局甲先发球.
        \begin{enumerate}
            \item 求该局打 4 个球甲赢的概率
            \item 求该局打 5 个球结束的概率
        \end{enumerate}
        \begin{jieda}
            \begin{enumerate}
                \item 设甲发球甲赢为事件 $A$ ，乙发球甲赢为事件 $B$ ，该局打 4 个球甲赢为事件 $C$
                
                由题意可知: $P\left( A\right)  = \dfrac{2}{3},P\left( B\right)  = \dfrac{1}{4},P\left( \bar{A}\right)  = \dfrac{1}{3},P\left( \bar{B}\right)  = \dfrac{3}{4}$ 且 $C = A\bar{B}{AB}$ ,
                
                可得 $P\left( C\right)  = P\left( {A\bar{B}{AB}}\right)  = P\left( A\right) P\left( \bar{B}\right) P\left( A\right) P\left( B\right)  = \dfrac{2}{3} \times  \dfrac{3}{4} \times  \dfrac{2}{3} \times  \dfrac{1}{4} = \dfrac{1}{12}$ ，
                
                所以该局打 4 个球甲赢的概率为 $\dfrac{1}{72}$

                \item 设该局打 5 个球结束时甲赢为事件 $D$ ，乙赢为事件 $E$ ，打 5 个球结束为事件 $F$

                可知事件 $D,E$ 为互斥事件,且 $D = \bar{A}B\bar{A}{BA},E = A\bar{B}A\bar{B}\bar{A},F = D \cup  E$ ,
                
                则 $P\left( D\right)  = P\left( {\bar{A}B\bar{A}{BA}}\right)  = P\left( \bar{A}\right) P\left( B\right) P\left( \bar{A}\right) P\left( B\right) P\left( A\right)  = \dfrac{1}{3} \times  \dfrac{1}{4} \times  \dfrac{1}{3} \times  \dfrac{1}{4} \times  \dfrac{2}{3} = \dfrac{1}{216}$ ,
                
                $P\left( E\right)  = P\left( {A\bar{B}A\bar{B}\bar{A}}\right)  = P\left( A\right) P\left( \bar{B}\right) P\left( A\right) P\left( \bar{B}\right) P\left( \bar{A}\right)  = \dfrac{2}{3} \times  \dfrac{3}{4} \times  \dfrac{2}{3} \times  \dfrac{3}{4} \times  \dfrac{1}{3} = \dfrac{1}{12},$
                
                可得 $P\left( F\right)  = P\left( {D \cup  E}\right)  = P\left( D\right)  + P\left( E\right)  = \dfrac{1}{216} + \dfrac{1}{12} = \dfrac{19}{216}$ ,
                
                所以该局打 5 个球结束的概率为 $\dfrac{19}{216}$
            \end{enumerate}
        \end{jieda}
        
        \item 某商场为了吸引顾客，规定购买一定价值的商品可以获得一次抽奖机会, 奖品价值分别为 10 元、20 元、30 元、40 元. 已知甲抽到价值为 10 元、20 元、30 元、40 元的奖品的概率分别为 $\dfrac{1}{6},\dfrac{1}{3},\dfrac{1}{3},\dfrac{1}{6}$ ，且每次抽奖结果相互独立.
        \begin{enumerate}
            \item 已知甲参与抽奖两次，求甲两次抽到的奖品价值不同的概率
            \item 求甲参与抽奖三次，抽到两种不同价值的奖品，且获得的奖品价值总和不低于 80 元的概率
        \end{enumerate}
        \begin{jieda}
            \begin{enumerate}
                \item 记甲两次抽到相同奖品为事件 $A$ ，
                
                记甲在一次抽奖中抽到值为 10 元、 20 元、 30 元、 40 元分别为事件 ${B}_{1},{B}_{2},{B}_{3},{B}_{4}$ ,
                
                则 $P\left( {B}_{1}\right)  = \dfrac{1}{6},P\left( {B}_{2}\right)  = \dfrac{1}{3},P\left( {B}_{3}\right)  = \dfrac{1}{3},P\left( {B}_{4}\right)  = \dfrac{1}{6}$ ,
                
                $P\left( A\right)  = P\left( {B}_{1}\right)  \times  P\left( {B}_{1}\right)  + P\left( {B}_{2}\right)  \times  P\left( {B}_{2}\right)  + P\left( {B}_{3}\right)  \times  P\left( {B}_{3}\right)  + P\left( {B}_{4}\right)  \times  P\left( {B}_{4}\right)$
                
                $= \dfrac{1}{6} \times  \dfrac{1}{6} + \dfrac{1}{3} \times  \dfrac{1}{3} + \dfrac{1}{3} \times  \dfrac{1}{3} + \dfrac{1}{6} \times  \dfrac{1}{6} = \dfrac{5}{18}$ ,
                
                所以甲两次抽到的奖品价值不同的概率为 $1 - P\left( A\right)  = 1 - \dfrac{5}{18} = \dfrac{13}{18}$

                \item 甲参与抽奖三次，抽到两种不同价值的奖品，所以其中一种奖品抽到两次，另一种抽到一次，又获得的奖品价值总和不低于 80 元,故可能两次抽到 40 元，一次抽到 30 元或两次抽到 40 元，一次抽到 20 元或两次抽到 40 元，一次抽到 10 元或两次抽到 30 元，一次抽到 40 元或两次抽到 30 元，一次抽到 20 元或两次抽到 20 元，一次抽到 40 元
                
                两次抽到 40 元，一次抽到 30 元的概率 ${P}_{1} = 3 \times  {\left( \dfrac{1}{6}\right) }^{2} \times  \dfrac{1}{3} = \dfrac{1}{36}$ ,
                
                两次抽到 40 元,一次抽到 20 元的概率 ${P}_{2} = 3 \times  {\left( \dfrac{1}{6}\right) }^{2} \times  \dfrac{1}{3} = \dfrac{1}{36}$ ,
                
                两次抽到 40 元，一次抽到 10 元的概率 ${P}_{3} = 3 \times  {\left( \dfrac{1}{6}\right) }^{2} \times  \dfrac{1}{6} = \dfrac{1}{72}$ ，
                
                两次抽到 30 元，一次抽到 40 元的概率 ${P}_{4} = 3 \times  {\left( \dfrac{1}{3}\right) }^{2} \times  \dfrac{1}{6} = \dfrac{1}{18}$ .
                
                两次抽到 30 元，一次抽到 20 元的概率 ${P}_{5} = 3 \times  {\left( \dfrac{1}{3}\right) }^{2} \times  \dfrac{1}{3} = \dfrac{1}{9}$ ，
                
                两次抽到 20 元，一次抽到 40 元的概率 ${P}_{6} = 3 \times  {\left( \dfrac{1}{3}\right) }^{2} \times  \dfrac{1}{6} = \dfrac{1}{18}$
                
                所以获得的奖品价值总和不低于 80 元的概率为:$P = \dfrac{1}{18} + \dfrac{1}{9} + \dfrac{1}{18} + \dfrac{1}{72} + \dfrac{1}{36} + \dfrac{1}{36} = \dfrac{7}{24}.$
            \end{enumerate}
        \end{jieda}
        
        \item 甲、乙两人轮流投篮，每人每次投一球. 约定甲先投且先投中者获胜，一直到有人获胜或每人都已投球 3 次时投篮结束. 设甲每次投篮投中的概率为 $\dfrac{1}{3}$ ，乙每次投篮投中的概率为 $\dfrac{1}{2}$ ，且各次投篮互不影响.
        \begin{enumerate}
            \item 求乙获胜的概率
            \item 求投篮结束时乙只投了 2 个球的概率
        \end{enumerate}
        \begin{jieda}
            \begin{enumerate}
                \item 设 ${A_k},{B_k}$ 分别表示甲、乙在第 $k$ 次投篮时投中，则 $P\left( {A}_{k}\right)  = \dfrac{1}{3},P\left( {B}_{k}\right)  = \dfrac{1}{2},k = 1,2,3$ .
                
                记“乙获胜”为事件 $C$
                
                $P\left( C\right)  = P\left( {{\bar{A}}_{1}{B}_{1}}\right)  + P\left( {{\bar{A}}_{1}{\bar{B}}_{1}{\bar{A}}_{2}{B}_{2}}\right)  + P\left( {{\bar{A}}_{1}{\bar{B}}_{1}{\bar{A}}_{2}{\bar{B}}_{2}{\bar{A}}_{3}{B}_{3}}\right) = P\left( {\bar{A}}_{1}\right) P\left( {B}_{1}\right)  + P\left( {\bar{A}}_{1}\right) P\left( {\bar{B}}_{1}\right) P\left( {\bar{A}}_{2}\right) P\left( {B}_{2}\right)  + P\left( {\bar{A}}_{1}\right) P\left( {\bar{B}}_{1}\right) P\left( {\bar{A}}_{2}\right) P\left( {\bar{B}}_{2}\right) P\left( {\bar{A}}_{3}\right) P\left( {B}_{3}\right) = \dfrac{2}{3} \times  \dfrac{1}{2} + {\left( \dfrac{2}{3}\right) }^{2} \times  {\left( \dfrac{1}{2}\right) }^{2} + {\left( \dfrac{2}{3}\right) }^{3} \times  {\left( \dfrac{1}{2}\right) }^{3} = \dfrac{13}{27}$

                \item 记“投篮结束时乙只投了 2 个球”为事件 $D$ ，
                
                $P\left( D\right)  = P\left( {{\bar{A}}_{1}{\bar{B}}_{1}{\bar{A}}_{2}{B}_{2}}\right)  + P\left( {{\bar{A}}_{1}{\bar{B}}_{1}{\bar{A}}_{2}{\bar{B}}_{2}{A}_{3}}\right) = P\left( {\bar{A}}_{1}\right)  \cdot  P\left( {\bar{B}}_{1}\right) P\left( {\bar{A}}_{2}\right) P\left( {B}_{2}\right)  + P\left( {\bar{A}}_{1}\right) P\left( {\bar{B}}_{1}\right) P\left( {\bar{A}}_{2}\right) P\left( {\bar{B}}_{2}\right) P\left( {A}_{3}\right) = {\left( \dfrac{2}{3}\right) }^{2} \times  {\left( \dfrac{1}{2}\right) }^{2} + {\left( \dfrac{2}{3}\right) }^{2} \times  {\left( \dfrac{1}{2}\right) }^{2} \times  \dfrac{1}{3} = \dfrac{4}{27}$
            \end{enumerate}
        \end{jieda}
        
        \item 临川二中两名优秀学子小明、小华同学独立地参加中国科技大学少科班的入学面试，入学面试时共有 3 道题目，答对 2 道题则通过面试(前 2 道题都答对或都答错，第 3 道题均不需要回答). 已知小明答对每道题目的概率均为 $\dfrac{3}{5}$ ，小华答对每道题目的概率依次为 $\dfrac{2}{3}\text{ 、 }\dfrac{2}{3}\text{ 、 }\dfrac{1}{2}$ ，且小明、小华两人对每道题能否答对相互独立. 记“小明只回答 2 道题就结束面试”为事件 $A$ ，记“小华 3 道题都回答且通过面试”为事件 $B$ .
        \begin{enumerate}
            \item 求事件 $A$ 发生的概率 $P\left( A\right)$
            \item 求事件 $A$ 和事件 $B$ 同时发生的概率 $P\left( {AB}\right)$
            \item 求小明、小华两人恰有一人通过面试的概率
        \end{enumerate}
        \begin{jieda}
            \begin{enumerate}
                \item 若事件 $A$ 发生，则小明前两题都答对或都答错，所以， $P\left( A\right)  = {\left( \dfrac{3}{5}\right) }^{2} + {\left( \dfrac{2}{5}\right) }^{2} = \dfrac{13}{25}$
                \item 若事件 $B$ 发生，则小华前两题一题，答错一题，第三题答对，根据题意则小华 3 道题都回答且通过面试的概率为 $P\left( B\right)  = \dfrac{2}{3} \times  \dfrac{1}{3} \times  \dfrac{1}{2} + \dfrac{1}{3} \times  \dfrac{2}{3} \times  \dfrac{1}{2} = \dfrac{2}{9}$ .
                
                由题意可知,事件 $A\text{ 、 }B$ 相互独立,则 $P\left( {AB}\right)  = P\left( A\right)  \cdot  P\left( B\right)  = \dfrac{13}{25} \times  \dfrac{2}{9} = \dfrac{26}{225}$

                \item 记小明没有通过面试为事件 $C$，分前两道回答对一道且最后一道错误或前两道均回答错误两种情况，则小明没有通过面试的概率为 $P\left( C\right)  = \dfrac{3}{5} \times  \dfrac{2}{5} \times  \dfrac{2}{5} + \dfrac{2}{5} \times  \dfrac{3}{5} \times  \dfrac{2}{5} + \dfrac{2}{5} \times  \dfrac{2}{5} = \dfrac{44}{125}$ ,
                
                可得小明通过面试的概率为 $P\left( \bar{C}\right)  = 1 - \dfrac{44}{125} = \dfrac{81}{125}$ ,
                
                而由 (1) 可得小华通过面试的事件记为 $D$ ,则概率为 $P\left( D\right)  = \dfrac{2}{9} + \dfrac{2}{3} \times  \dfrac{2}{3} = \dfrac{2}{3}$ ,
                
                由题意可知,事件 $C\text{ 、 }D$ 相互独立,
                
                则小明、小华两人恰有一人通过面试的事件记为 $E$ ，
                
                则概率为 $P\left( E\right)  = P\left( {D\bar{C} + \bar{D}C}\right)  = P\left( {D\bar{C}}\right)  + P\left( {\bar{D}C}\right)  = \dfrac{81}{125} \times  \dfrac{1}{3} + \dfrac{44}{125} \times  \dfrac{2}{3} = \dfrac{169}{375}$
            \end{enumerate}
        \end{jieda}
    \end{enumerate}
    \item 古典概型综合
    \begin{enumerate}
        \item 已知集合 $M = \{  - 1,1,2,3\} ,N = \{  - {20}, - 8,4,9\}$ ，若分别从集合 $M,N$ 中随机抽取一个数 $m$ 和 $n$ ，二次函数 $f\left( x\right)  = m{x}^{2} - {nx} - 1$ . 记事件 $A$ 为“ $\lbrack  - 4, + \infty )$ 是二次函数 $y = f\left( x\right)$ 的单调递增区间”,事件 $B$ 为“ $( - \infty ,0\rbrack$ 是二次函数 $y = f\left( x\right)$ 的单调递减区间”.
        \begin{enumerate}
            \item 求数对 $\left( {m,n}\right)$ 的样本空间中所含样本点的个数
            \item 分别求事件 $A$ 、事件 $B$ 的概率
            \item 求事件 $A$ 、事件 $B$ 至少一个发生的概率
        \end{enumerate}
        \begin{jieda}
            \begin{enumerate}
                \item 数对 $\left( {m,n}\right)$ 的样本空间中所含样本点的个数 $4 \times  4 = {16}$ 个.
                \item 函数 $f\left( x\right)  = m{x}^{2} - {nx} - 1$ 的对称轴为 $x =  - \dfrac{-n}{2m} = \dfrac{n}{2m}$
                
                对于事件 $A$ ,则 $\left\{  \begin{matrix} m > 0 \\  \dfrac{n}{2m} \leq   - 4 \end{matrix}\right.$ ,即 $n \leq   - {8m} < 0$ ,因 $m \in  \{  - 1,1,2,3\} ,n \in  \{  - {20}, - 8,4,9\}$
                
                则满足事件 ${A}$ 的数对有 $\left( {1, - {20}}\right) ,\left( {1, - 8}\right) ,\left( {2, - {20}}\right)$ 共 3 个,故 $P\left({A}\right)  = \dfrac{3}{16}$
                
                对于事件 $B$ ,则 $\left\{  \begin{matrix} m > 0 \\  \dfrac{n}{2m} \geq  0 \end{matrix}\right.$ ,则 $\left\{  \begin{matrix} m > 0 \\  n \geq  0 \end{matrix}\right.$ ,满足事件 $B$ 的数对 $\left( {m,n}\right)$ 有 $\left( {1,4}\right) ,\left( {1,9}\right) ,\left( {2,4}\right) ,\left( {2,9}\right) ,\left( {3,4}\right) ,\left( {3,9}\right)$ , 共 6 个，故 $P\left( B\right)  = \dfrac{6}{16} = \dfrac{3}{8}$

                \item 由 (2) 可知,事件 $A$ 发生时有 $n \leq   - {8m} < 0$ ,事件 $B$ 发生时有 $\left\{  \begin{array}{l} m > 0 \\  n \geq  0 \end{array}\right.$ ,则事件 $A$ 与事件 $B$ 互斥, 则事件 $A$ 、事件 $B$ 至少一个发生的概率 $P = P\left( A\right)  + P\left( B\right)  = \dfrac{9}{16}$
            \end{enumerate}
        \end{jieda}
        \item 袋中有 7 只大小形状相同颜色不全相同的小猫摆件，分别为黑猫、白猫、红猫,某同学从中任意取一只小猫摆件,得到黑猫或白猫的概率是 $\dfrac{5}{7}$ 得到白猫或红猫的概率是 $\dfrac{4}{7}$ . 试求:
        \begin{enumerate}
            \item 某同学从中任取一只小猫摆件，得到黑猫、白猫、红猫的概率各是多少？
            \item 某同学从中任取两只小猫摆件，得到的两只小猫颜色不相同的概率是多少？
        \end{enumerate}
        \begin{jieda}
            \begin{enumerate}
                \item 从中任取一只小猫摆件，分别记得到黑猫、白猫、红猫为事件 $A,B,C$ , 由于 $A,B,C$ 为互斥事件，则$\left\{\begin{array}{l}
                     P\left( {A \cup B}\right)  = P\left( A\right)  + P\left( B\right)  = \dfrac{5}{7}  \\
                     P\left( {B \cup C}\right)  = P\left( B\right)  + P\left( C\right)  = \dfrac{4}{7}  \\
                     P\left( {A \cup B \cup C}\right)  = P\left( A\right)  + P\left( B\right)  + P\left( C\right)  = 1
                \end{array}\right.$
                故$\left\{\begin{array}{l}
                     P\left( A\right)  = \dfrac{3}{7}  \\
                     P\left( B\right)  = \dfrac{2}{7} \\
                     P\left( C\right)  = \dfrac{2}{7} 
                \end{array}\right.$

                所以任取一只小猫摆件，得到黑猫、白猫、红猫的概率分别是 $\dfrac{3}{7}\dfrac{2}{7}\dfrac{2}{7}$

                \item 从7只小猫摆件中取出两只小猫摆件的基本事件有 $C_7^2 = 21$ 个，其中两只摆件是黑猫的基本事件有3个，两只白猫的基本事件 1 个，两只红猫的基本事件有1个
                
                于是两只小猫摆件同色的概率为 $\dfrac{3 + 1 + 1}{21} = \dfrac{5}{21}$ ,
                
                则两只小猫摆件颜色不相同的概率是 $1 - \dfrac{5}{21} = \dfrac{16}{21}$ .
            \end{enumerate}
        \end{jieda}
        \item 甲、乙、丙三人玩“剪刀、石头、布”游戏(剪刀赢布，布赢石头，石头赢剪刀)，规定每局中: \ding{172} 三人出现同一种手势，每人各得 1 分； \ding{173} 三人出现两种手势，赢者得 2 分，输者负 1 分； \ding{174} 三人出现三种手势均得 0 分.当有人累计得 3 分及以上时，游戏结束，得分最高者获胜，已知三人之间及每局游戏互不受影响.
        \begin{enumerate}
            \item 求甲在一局中得 2 分的概率 ${P}_{1}$
            \item 求游戏经过两局后甲恰得 3 分且为唯一获胜者的概率 ${P}_{2}$
        \end{enumerate}
        \begin{jieda}
            \begin{enumerate}
                \item 每局中共有 $3\times3\times3=27$ 种情况，其中甲在一局中得 2 分的情况有(按甲乙丙顺序):
                
                (剪刀、剪刀、布) 、 (剪刀、布、剪刀) 、 (剪刀、布、布) 、
                
                (石头、石头、剪刀)、(石头、剪刀、石头)、(石头、剪刀、剪刀)、
                
                (布、布、石头)、(布、石头、布)、(布、石头、石头)、
                
                一共有 9 种情况
                
                所以甲在一局中得 2 分的概率 ${P}_{1} = \dfrac{9}{27} = \dfrac{1}{3}$ .

                \item 游戏经过两局后甲恰得 3 分且为唯一获胜者的情况有 2 种:
                \begin{dingautolist}{172}
                    \item 第一局甲得 2 分，第二局甲得 1 分，此时有乙第一局得负 1 分，第二局得 1 分；丙第一局得负 1 分，第二局得 1 分

                    对应概率为$\dfrac{3}{27} \times \dfrac{3}{27} = \dfrac{1}{81}$.
                    
                    \item 第一局甲得 1 分，第二局甲得 2 分，此时有乙第一局得 1 分，第二局得负 1 分；丙第一局得 1 分，第二局得负 1 分

                    对应概率为$\dfrac{3}{27} \times \dfrac{3}{27} = \dfrac{1}{81}$.
                \end{dingautolist}

                综上所述: 游戏经过两局后甲恰得 3 分且为唯一获胜者的概率 ${P}_{2} = \dfrac{1}{81} + \dfrac{1}{81} = \dfrac{2}{81}$
            \end{enumerate}
        \end{jieda}
        \item 某班级在庆元旦联欢会时，主持人安排了跳双人舞、独唱、和独奏节目,指定 3 个男生和 2 个女生来参与,把五个人分别编号为1,2,3,4,5,其中1,2,3号是男生,4,5 号是女生. 将每个人的编号分别写在 5 张相同的卡片上, 放入一个不透明的箱子中, 并搅拌均匀, 每次从中随机取出一张卡片, 取出谁的编号谁就参与表演节目.
        \begin{enumerate}
            \item 为了选出 2 人来表演双人舞，不放回地抽取 2 张卡片，求选出的 2 人不全是男生的概率
            \item 为了确定表演独唱和独奏的人选，抽取并记录第一张卡片后，又放回箱子中，充分混合后再从中抽取第二张卡片，求：
            \begin{dingautolist}{172}
                \item 独唱和独奏由同一个人表演的概率
                \item 选出的不全是男生的概率
            \end{dingautolist}
        \end{enumerate}
        \begin{jieda}
            \begin{enumerate}
                \item 依题意, 不放回地抽取 2 张卡片, 抽取共有 $C_5^2 = 10$ 种可能的结果，全部是男生有$C_3^2 = 3$种可能的结果，则不全为男生由$10-3=7$种可能结果

                因此, $P = \dfrac{7}{10}$ ,即选出的 2 人不全是男生的概率为 $\dfrac{7}{10}$
                \item \begin{dingautolist}{172}
                    \item 依题意, 有放回地抽取 2 张卡片, 抽取共有 $5 \times 5 = 25$ 种可能的结果，其中抽取到同一个人的可能结果有5种，故独唱和独奏由同一个人表演的概率为 $\dfrac{1}{5}$
                    \item 设第一次选出的是男生为事件$A_1$，第二次选出的是男生为事件$A_2$，两者独立，故$P(A_1A_2) = P(A_1)P(A_2) = \dfrac{3}{5} \times \dfrac{3}{5} = \dfrac{9}{25}$，故选出的不全是男生的概率为 $1 - \dfrac{9}{25} = \dfrac{16}{25}$
                \end{dingautolist}
            \end{enumerate}
        \end{jieda}
        \item 某班元旦联欢会上开展趣味抽奖小游戏，在不透明的盒子里装有标号为 1,2 的两个红球和标号为 3,4,5 的三个白球, 五个小球除颜色外完全相同, 参与游戏的同学从中任取 1 个, 有放回的抽取 2 次, 根据抽到小球的情形分别设置一, 二, 三等奖.班委会讨论了以下两种规则:
        
        规则一:若抽到两个红球且标号和为偶数获一等奖，抽到两个白球且标号和为偶数获二等奖，抽到两个球标号和为奇数获三等奖, 其余不获奖;
        
        规则二:若抽到两个红球且标号和为奇数获一等奖，抽到两个球的标号和为 5 的倍数获二等奖，抽到两个球标号和为偶数获三等奖, 其余不获奖.

        \begin{enumerate}
            \item 请以标号$(x,y)$写出两次抽取小球的所有结果(其中 $x$ ， $y$ 分别为第一、第二次抽到的小球标号)
            \item 求两种规则下获得二等奖的概率
            \item 请问哪种规则获奖概率更大，并说明理由
        \end{enumerate}
        \begin{jieda}
            \begin{enumerate}
                \item 两次抽取小球的所有可能结果为:
                
                $\left( {1,1}\right) ,\left( {1,2}\right) ,\left( {1,3}\right) ,\left( {1,4}\right) ,\left( {1,5}\right)$ ,
                
                $\left( {2,1}\right) ,\left( {2,2}\right) ,\left( {2,3}\right) ,\left( {2,4}\right) ,\left( {2,5}\right)$ ,
                
                $\left( {3,1}\right) ,\left( {3,2}\right) ,\left( {3,3}\right) ,\left( {3,4}\right) ,\left( {3,5}\right)$ ,
                
                $\left( {4,1}\right) ,\left( {4,2}\right) ,\left( {4,3}\right) ,\left( {4,4}\right) ,\left( {4,5}\right)$ ,
                
                $\left( {5,1}\right) ,\left( {5,2}\right) ,\left( {5,3}\right) ,\left( {5,4}\right) ,\left( {5,5}\right)$ 
                \item 记规则一中获得二等奖为事件 ${A}_{2}$ ，记规则二中获得二等奖为事件 ${B}_{2}$
                
                事件 ${A}_{2}$ 包含 $\left( {3,3}\right) ,\left( {3,5}\right) ,\left( {4,4}\right) ,\left( {5,3}\right) ,\left( {5,5}\right)$ 五个样本点,
                
                故 $P\left( {A}_{2}\right)  = \dfrac{5}{25} = \dfrac{1}{5}$
                
                事件 ${B}_{2}$ 包含 $\left( {1,4}\right) ,\left( {2,3}\right) ,\left( {3,2}\right) ,\left( {4,1}\right) ,\left( {5,5}\right)$ 五个样本点,
                
                故 $P\left( {B}_{2}\right)  = \dfrac{5}{25} = \dfrac{1}{5}$
                \item 两种规则获奖的概率一样大

                理由如下: 记规则一获得一,二,三等奖分别为事件 ${A}_{1},{A}_{2},{A}_{3}$ ,
                
                规则二获得一，二，三等奖分别为事件 ${B}_{1},{B}_{2},{B}_{3}$ ,
                
                ${A}_{1}$ 包含 $\left( {1,1}\right) ,\left( {2,2}\right)$ 两个样本点, $\therefore P\left( {A}_{1}\right)  = \dfrac{2}{25}$
                
                ${A}_{3}$ 包含 $\left( {1,2}\right) ,\left( {1,4}\right) ,\left( {2,1}\right) ,\left( {2,3}\right) ,\left( {2,5}\right) ,\left( {3,2}\right) ,\left( {3,4}\right) ,\left( {4,1}\right) ,\left( {4,3}\right) ,\left( {4,5}\right) ,\left( {5,2}\right) ,\left( {5,4}\right)$ 12个样本点，$\therefore P\left( {A}_{3}\right)  = \dfrac{12}{25}$
                
                所以规则一获奖的概率
                
                $P\left( {{A}_{1} \cup {A}_{2} \cup {A}_{3}}\right)  = P\left( {A}_{1}\right)  + P\left( {A}_{2}\right)  + P\left( {A}_{3}\right)  = \dfrac{2}{25} + \dfrac{1}{5} + \dfrac{12}{25} = \dfrac{19}{25}$
                
                ${B}_{1}$ 包含 $\left( {1,2}\right) ,\left( {2,1}\right)$ 两个样本点, $\therefore P\left( {B}_{1}\right)  = \dfrac{2}{25}$ ;
                
                ${B}_{3}$ 包含 $\left( {1,1}\right) ,\left( {1,3}\right) ,\left( {1,5}\right) ,\left( {2,2}\right) ,\left( {2,4}\right) ,\left( {3,1}\right) ,\left( {3,3}\right) ,\left( {3,5}\right) ,\left( {4,2}\right) ,\left( {4,4}\right) ,\left( {5,1}\right) ,\left( {5,3}\right)$ 12个样本点，($\left( {5,5}\right)$在 ${B}_{2}$ 中已经记录,不再计算)，$\therefore P\left( {B}_{3}\right)  = \dfrac{12}{25}$
                
                所以规则二获奖的概率
                
                $P\left( {{B}_{1} \cup {B}_{2} \cup {B}_{3}}\right)  = P\left( {B}_{1}\right)  + P\left( {B}_{2}\right)  + P\left( {B}_{3}\right)  = \dfrac{2}{25} + \dfrac{1}{5} + \dfrac{12}{25} = \dfrac{19}{25}$
            \end{enumerate}
        \end{jieda}
        
        \item 某商场有奖销售中，购满 100 元商品得一张奖券，多购多得，每 1000 张奖券为一个开奖单位.设特等奖 1 个, 一等奖 10 个, 二等奖 50 个.设 1 张奖券中特等奖、一等奖、二等奖的事件分别为 $A,B,C$ ,求:
        \begin{enumerate}
            \item $P\left( A\right) ,P\left( B\right) ,P\left( C\right)$
            \item 抽取 1 张奖券中奖概率
            \item 抽取 1 张奖券不中特等奖或一等奖的概率
        \end{enumerate}
        \begin{jieda}
            \begin{enumerate}
                \item 因为每 1000 张奖券中设特等奖 1 个，一等奖 10 个，二等奖 50 个，所以 $P\left( A\right)  = \dfrac{1}{1000},P\left( B\right)  = \dfrac{10}{1000} = \dfrac{1}{100},P\left( C\right)  = \dfrac{50}{1000} = \dfrac{1}{20}$ .

                \item 设“抽取 1 张奖券中奖”为事件 $D$ ，则 $P\left( D\right)  = P\left( A\right)  + P\left( B\right)  + P\left( C\right)  = \dfrac{1}{1000} + \dfrac{1}{100} + \dfrac{1}{20} = \dfrac{61}{1000}$ .

                \item 设“抽取 1 张奖券不中特等奖或一等奖”为事件 $E$ ，则 $P\left( E\right)  = 1 - P\left( A\right)  - P\left( B\right)  = 1 - \dfrac{1}{1000} - \dfrac{1}{100} = \dfrac{989}{1000}$ .
            \end{enumerate}
        \end{jieda}
    \end{enumerate}
    \item 古典概型和统计结合
    \begin{enumerate}
        \item 某地文化和旅游局统计了春节期间 100 个家庭的旅游支出情况， 统计得到这 100 个家庭的旅游支出(单位:千元)数据，按 [3,5), [5,7), [7,9), [9,11), [11,13] 分成 5 组，并绘制成频率分布直方图, 如图所示.
        
        \begin{center}
        \includegraphics[width=0.4\textwidth]{bodyxb/src/bo_d67fmaf7aajc739ac4u0_19_138_203_556_330_0.jpg}
        \end{center}

        \begin{enumerate}
            \item 估计这 100 个家庭的旅游支出的平均值(同一组中的数据用该组区间的中点值作为代表)

            \item 估计这 100 个家庭的旅游支出的第 70 百分位数(结果保留一位小数)

            \item 在这 100 个家庭中，旅游支出在 $[3,5)$, $[5,7)$ (千元)的家庭中，按分层抽样的方法抽取 5 个家庭，再从这 5 个家庭中抽取 2 个家庭，求至少有 1 个家庭的旅游支出在 $[3,5)$ 千元内的概率.
        \end{enumerate}
        \begin{jieda}
            \begin{enumerate}
                \item 估计这 100 个家庭的旅游支出的平均值为: $2\left( {4 \times  {0.050} + 6 \times  {0.075} + 8 \times  {0.175} + {10} \times  {0.150} + {12} \times  {0.050}}\right)  = {8.3}$ (千元).

                \item 由频率分布直方图知，旅游支出在 $\lbrack 3,9)$ 千元的频率为 $2 \times  {0.050} + 2 \times  {0.075} + 2 \times  {0.175} = {0.6}$
                
                在 $\lbrack 3,{11})$ 千元的频率为 ${0.6} + 2 \times  {0.15} = {0.9}$ ,则这 100 个家庭的旅游支出的第 70 百分位数 $a \in  \lbrack 9,{11})$，则 ${0.6} + \left( {a - 9}\right)  \times  {0.150} = {0.7}$ ，解得 $a \approx  {9.7}$
                
                所以估计这 100 个家庭的旅游支出的第 70 百分位数为 9.7 .

                \item 按分层抽样的方法抽取 5 个家庭, $\lbrack 3,5),\lbrack 5,7)$ 千元内抽取家庭数之比为$0.050:0.075$，即$2:3$,所以 $\lbrack 3,5)$ 千元内抽取 2 个家庭, $[5,7)$千元内抽取 3 个家庭
                
                设旅游支出在 $[3,5)$ 千元的 2 个家庭记为 $A\text{ 、 }B$ ,在 $[5,7)$ 千元的 3 个家庭记为 $C\text{ 、 }D\text{ 、 }E$ . 从这 5 个家庭中抽取 2 个家庭的所有可能情况有$C_5^2 = 10$ 种.
                
                2个家庭旅游支出均不在 $\lbrack 3,5)$ 千元内的情况有$C_3^2 = 3$种，故至少有 1 个家庭指出在$\lbrack 3,5)$ 千元内的情况有 7 种. 所以至少有 1 个家庭的旅游支出在 $[3,5)$ 千元内的概率 $P = \dfrac{7}{10}$
            \end{enumerate}
        \end{jieda}
        \item 某企业以“庆祝春节，迎接新年”为主题的职工歌手大赛决赛如期举行， 满分 100 分，共有 100 人参赛，将参赛歌手的成绩分成如下五组:第一组[50,60)，第二组[60,70)，第三组[70,80), 第四组[80,90), 第五组[90,100], 得到如图所示的频率分布直方图.
        
        \begin{center}
        \includegraphics[width=0.3\textwidth]{bodyxb/src/bo_d67fmaf7aajc739ac4u0_20_138_261_452_395_0.jpg}
        \end{center}

        \begin{enumerate}
            \item 根据频率分布直方图，求 $a$ 的值及参赛歌手的平均成绩(同一组数据用该组区间的中点值作代表)
            \item 根据频率分布直方图，求参赛歌手成绩的 40\% 分位数
            \item 从参赛成绩在 $\left\lbrack  {{50},{60}}\right\rbrack$ 和 $\left\lbrack  {{90},{100}}\right\rbrack$ 的歌手中，采用分层随机抽样方法抽取 6 名歌手，再从抽取的这 6 名歌手中随机抽取 2 名歌手，求这 2 名歌手比赛成绩在 $\lbrack {50},{60})$ 和 $\left\lbrack  {{90},{100}}\right\rbrack$ 内各 1 人的概率
        \end{enumerate}
        \begin{jieda}
            \begin{enumerate}
                \item 第一至第五组对应的频率分别为 ${0.010} \times  {10} = {0.1};{0.015} \times  {10} = {0.15}$ ;
                
                ${0.040} \times  {10} = {0.4};a \times  {10} = {10a};{0.005} \times  {10} = {0.05}$，所以 ${0.1} + {0.15} + {0.4} + {10a} + {0.05} = 1$ ,解得 $a = {0.030}$，所以参赛歌手的平均成绩为 ${0.1} \times  {55} + {0.15} \times  {65} + {0.4} \times  {75} + {0.3} \times  {85} + {0.05} \times  {95} = {75.5}$ 分.

                \item 由 ${0.1} + {0.15} = {0.25} < {0.4},{0.1} + {0.15} + {0.4} = {0.65} > {0.4}$，得参赛歌手成绩的 40\% 分位数为 ${70} + \dfrac{{0.4} - {0.25}}{0.4} \times  {10} = {73.75}$ 分.

                \item 由 ${0.1} : {0.05} = 2 : 1$ ，得这 6 人中参赛成绩在 $\lbrack {50},{60})$ 的人数为 $6 \times  \dfrac{2}{3} = 4$ 人，分别记为 $a,b,c,d$ ； 在 $\left\lbrack  {{90},{100}}\right\rbrack$ 的人数为 $6 - 4 = 2$ 人，分别记为 $x$ ， $y$ .
                
                在这 6 个人中抽取 2 个人，共 $\left( {a,b}\right) ,\left( {a,c}\right) ,\left( {a,d}\right) ,\left( {a,x}\right) ,\left( {a,y}\right) ,\left( {b,c}\right) ,\left( {b,d}\right) ,\left( {b,x}\right) ,\left( {b,y}\right) ,\left( {c,d}\right)$ , $\left( {c,x}\right) ,\left( {c,y}\right) ,\left( {d,x}\right) ,\left( {d,y}\right) ,\left( {x,y}\right) ,{15}$ 个基本事件,
                
                这 2 名歌手比赛成绩在 $\lbrack {50},{60})$ 和 $\left\lbrack  {{90},{100}}\right\rbrack$ 内各 1 人，共 $\left( {a,x}\right) ,\left( {a,y}\right) ,\left( {b,x}\right) ,\left( {b,y}\right) ,\left( {c,x}\right) ,\left( {c,y}\right) ,\left( {d,x}\right)$ ， $\left( {d,y}\right) ,8$ 个基本事件,
                
                故这 2 名歌手比赛成绩在 $\lbrack {50},{60})$ 和 $\left\lbrack  {{90},{100}}\right\rbrack$ 内各 1 人的概率为 $\dfrac{8}{15}$
            \end{enumerate}
        \end{jieda}
        
        \item 其校为了解学生的综合素养情况，从该校学生中随机地抽取了 40 名学生作为样本，进行综合素养测评，将他们的得分(满分:100 分)分成$\lbrack {40},{50}),\lbrack {50},{60}),\lbrack {60},{70}),\lbrack {70},{80}),\lbrack {80},{90}),\left\lbrack  {{90},{100}}\right\rbrack$ ,共六组. 根据他们的得分绘制了如图所示的频率分布直方图.

        \begin{center}
        \includegraphics[width=0.4\textwidth]{bodyxb/src/bo_d67fmaf7aajc739ac4u0_21_137_468_516_382_0.jpg}
        \end{center}

        \begin{enumerate}
            \item 从得分低于 60 分的样本中随机地选取 2 个样本，求这 2 个样本的得分在同一组的概率
            \item 若在 $\lbrack {80},{90})$ 内的样本得分的平均数为 86 分,方差为 10,在 $\left\lbrack  {{90},{100}}\right\rbrack$ 内的样本得分的平均数为 92 分,方差为 6 , 求在 $\left\lbrack  {{80},{100}}\right\rbrack$ 内的样本得分的平均数和方差
        \end{enumerate}
        \begin{jieda}
            \begin{enumerate}
                \item 由图可知， ${{10} \times  \left( {{0.005} + {0.020} + {0.025} + {0.030} + {2a}}\right) } = 1$ ，解得 $a = {0.010}$ ，则在 $\lbrack {40},{50})$ 内的样本容量为 ${40} \times  {10} \times  {0.005} = 2$ ,将这 2 个样本分别记为 ${{A}}_{1},{{A}}_{2}$
                
                在 $\lbrack {50},{60})$ 内的样本容量为 ${40} \times  {10} \times  {0.010} = 4$ ,将这 4 个样本分别记为 ${{B}}_{1},{{B}}_{2},{{B}}_{3},{{B}}_{4}$ .
                
                从中随机地选取 2 个，可知样本空间
                
                $\Omega  = \left\{  {{{A}}_{1}{{A}}_{2},{{A}}_{1}{{B}}_{1},{{A}}_{1}{{B}}_{2},{{A}}_{1}{{B}}_{3},{{A}}_{1}{{B}}_{4},{{A}}_{2}{{B}}_{1},{{A}}_{2}{{B}}_{2},{{A}}_{2}{{B}}_{3},{{A}}_{2}{{B}}_{4},{{B}}_{1}{{B}}_{2},}\right.$
                
                $\left. {{{B}}_{1}{{B}}_{3},{{B}}_{1}{{B}}_{4},{{B}}_{2}{{B}}_{3},{{B}}_{2}{{B}}_{4},{{B}}_{3}{{B}}_{4}}\right\}$ ,
                
                共有 15 个样本点.
                
                用事件 $X$ 表示“这 2 个样本的得分在同一组”,
                
                则 $X = \left\{  {{A}_{1}{A}_{2},{B}_{1}{B}_{2},{B}_{1}{B}_{3},{B}_{1}{B}_{4},{B}_{2}{B}_{3},{B}_{2}{B}_{4},{B}_{3}{B}_{4}}\right\}$ ,有 7 个样本点,
                
                则 $P\left( X\right)  = \dfrac{7}{15}$ ，即这 2 个样本得分在同一组的概率为 $\dfrac{7}{15}$

                \item 由图可知，在 $\lbrack {80},{90})$ 内的样本数与在 $\left\lbrack  {{90},{100}}\right\rbrack$ 内的样本数之比为 $2 : 1$
                
                所以在 $\left\lbrack  {{80},{100}}\right\rbrack$ 内的样本得分的平均数 $\bar{x} = \dfrac{2}{3} \times  {86} + \dfrac{1}{3} \times  {92} = {88}$ 分,
                
                方差 ${s}^{2} = \dfrac{2}{3}\left\lbrack  {{10} + {\left( {88} - {86}\right) }^{2}}\right\rbrack   + \dfrac{1}{3}\left\lbrack  {6 + {\left( {88} - {92}\right) }^{2}}\right\rbrack   = \dfrac{50}{3}$
            \end{enumerate}
        \end{jieda}
    \end{enumerate}
    \item 条件概率的计算与性质
    \begin{enumerate}
        \item 已知事件 $A,B$ ，若 $P\left( {B \mid  A}\right)  = \dfrac{3}{4},P\left( A\right)  = \dfrac{1}{3}$ ，则 $P\left( {AB}\right)  =$ \blkda{$\dfrac{1}{4}$}
        \begin{jieda}
            由题可知, $P\left( {AB}\right)  = P\left( A\right) P\left( {B \mid  A}\right)  = \dfrac{1}{3} \times  \dfrac{3}{4} = \dfrac{1}{4}$
        \end{jieda}
        
        \item 已知 $P\left( {A|B}\right)  = {0.6},P\left( {B|A}\right)  = {0.3}$ ，且 $A,B$ 相互独立，则 $P\left( {AB}\right)  =$ \blkda{$0.18$}
        \begin{jieda}
            $A,B$ 相互独立, $\therefore P\left( {A|B}\right)  = P\left( A\right)  = {0.6},P\left( {B|A}\right)  = P\left( B\right)  = {0.3}$
            
            $P\left( {AB}\right)  = P\left( A\right) P\left( B\right)  = {0.6} \times  {0.3} = {0.18}$ .
        \end{jieda}

        \item 假设 $A,B$ 是两个事件，且 $P\left( A\right)  > 0,P\left( B\right)  > 0$ ，则下列结论一定正确的是\dotfill{(\kh{A})}
        \xx{$P\left( {A \mid  B}\right) P\left( B\right)  = P\left( {AB}\right)$}{$P\left( {A \mid  B}\right)  = P\left( {B \mid  A}\right)$}{$P\left( {A \mid  B}\right)  \leq  P\left( A\right) P\left( B\right)$}{$P\left( {A \mid  B}\right)  \leq  P\left( A\right)$}

        \item 下列说法正确的是\dotfill{(\kh{B})}
        \xx{$P\left( {B \mid  A}\right)  < P\left( {AB}\right)$}{$P\left( {B \mid  A}\right)  = \dfrac{P\left( B\right) }{P\left( A\right) }$ 是可能的}{$0 < P\left( {B \mid  A}\right)  < 1$}{$P\left( {A \mid  A}\right)  = 0$}

        \item 某医学院校计划从 5 名男生和 3 名女生中选派 2 人参加义诊活动，则在派出的 2 人中有男生的条件下，另一人恰好是女生的概率是\blkda{$\dfrac{3}{5}$}
        \begin{jieda}
            记“派出的 2 人中有男生”为事件 $A$ ，“另一人恰好是女生”为事件 $B$
            
            $P\left( {B|A}\right)  = \dfrac{P\left( {AB}\right) }{P\left( A\right) } = \dfrac{n\left( {AB}\right) }{n\left( A\right) } = \dfrac{{{C}}_{5}^{1}{{C}}_{3}^{1}}{{{C}}_{5}^{2} + {{C}}_{5}^{1}{{C}}_{3}^{1}} = \dfrac{3}{5}$
        \end{jieda}
        
        \item 投掷 3 枚质地均匀的骰子，设事件 $A{ = }$ “这 3 枚骰子朝上的点数之和为奇数”,事件 $B{ = }$ “恰有 1 枚骰子朝上的点数为奇数”,则 $P\left( {B \mid  A}\right)  =$\blkda{$\dfrac{3}{4}$}
        \begin{jieda}
            因为每枚骰子朝上的点数有奇数 1,3,5 三个,偶数有 2,4,6 三个，所以 3 枚骰子朝上的点数之和为奇数的情况有奇数+奇数+奇数, 偶数+偶数+奇数,
            
            共两种情况,可得 $P\left( A\right)  = \dfrac{{{C}}_{3}^{1}{{C}}_{3}^{1}{{C}}_{3}^1 + 3{{C}}_{3}^1{{C}}_{3}^1{{C}}_{3}^1}{{{C}}_{6}^1{{C}}_{6}^1{{C}}_{6}^1} = \dfrac{1}{2}$ ,
            
            恰有 1 枚骰子朝上的点数为奇数的情况有偶数+偶数+奇数,
            
            可得 $P\left( {AB}\right)  = \dfrac{3{{C}}_{3}^{1}{{C}}_{3}^{1}{{C}}_{3}^{1}}{{{C}}_{6}^{1}{{C}}_{6}^{1}{{C}}_{6}^{1}} = \dfrac{3}{8}$，则 $P\left( {B \mid  A}\right)  = \dfrac{P\left( {AB}\right) }{P\left( A\right) } = \dfrac{3}{4}$
        \end{jieda}
        
        \item 甲、乙分别用弓箭对准同一个弓箭靶，两人同时射箭.已知甲、乙中靶的概率分别为 0.5 和 0.4 ，且两人是否中靶互不影响，若弓箭靶被射中，则只被乙射中的概率为\blkda{$\dfrac{2}{7}$}
        \begin{jieda}
            设事件 ${A}$ : 甲中靶，事件 ${B}$ :乙中靶，事件 ${C}$ :弓箭靶被射中
            
            $P\left( A\right)  = {0.5},P\left( B\right)  = {0.4}$，所以 $P\left( C\right)  = 1 - P\left( \overline{A} \cap  \overline{B}\right)  = 1 - \left( {1 - {0.5}}\right) \left( {1 - {0.4}}\right)  = {0.7}$
            
            $P\left( {\bar{A}BC}\right) = P\left( {\bar{A}B}\right)  = \left( {1 - {0.5}}\right)  \times  {0.4} = {0.2}$，$P\left( {\left( {\bar{A}B}\right) |C}\right)  = \dfrac{0.2}{0.7} = \dfrac{2}{7}$
        \end{jieda}
    \end{enumerate}
    \item 利用全概率公式求概率
    \begin{enumerate}
        \item 某人参加抽奖游戏，现有三叠外形、大小、图案均相同的卡片， 分别有 10 张、15 张、20 张, 若每叠中有 2 张中奖卡片, 则随机选择一叠卡片抽取, 中奖的概率是\blkda{$\dfrac{13}{90}$}
        \begin{jieda}
            记事件 ${A}_{i} = \left\{  \begin{array}{ll} \text{ 在第 }i\text{ 叠卡片中抽奖, }i = 1,2,3\} , & \text{ 事件 }B = \{ \text{ 中奖 }\} , \end{array}\right.$
            
            则 $P\left( {B|{A}_{1}}\right)  = \dfrac{2}{10} = \dfrac{1}{5},P\left( {B|{A}_{2}}\right)  = \dfrac{2}{15},P\left( {B|{A}_{3}}\right)  = \dfrac{2}{20} = \dfrac{1}{10}$ .
            
            由全概率公式可得 $P\left( B\right)  = P\left( {A}_{1}\right) P\left( {B|{A}_{1}}\right)  + P\left( {A}_{2}\right) P\left( {B|{A}_{2}}\right)  + P\left( {A}_{3}\right) P\left( {B|{A}_{3}}\right) = \dfrac{1}{3} \times  \dfrac{1}{5} + \dfrac{1}{3} \times  \dfrac{2}{15} + \dfrac{1}{3} \times  \dfrac{1}{10} = \dfrac{13}{90}$
        \end{jieda}
        \item 无人酒店是利用人工智能与物联网技术，为客人提供自助人住等服务的新型酒店，胜在科技感与新奇感. 去某地旅游的游客有无人酒店和常规酒店两种选择，某游客去该地旅游，第一天随机选择一种酒店入住，如果第一天入住无人酒店，那么第二天还入住无人酒店的概率为 0.8 ，如果第一天入住常规酒店，那么第二天入住无人酒店的概率为 0.6 ，则该游客第二天入住无人酒店的概率为\blkda{$0.7$}
        \begin{jieda}
            记事件 ${A}_{1} =$ \{第一天入住无人酒店\}, ${A}_{2} =$ \{第二天入住无人酒店\}, ${B}_{1} =$ \{第一天入住常规酒店\},
            
            根据题意可知 $P\left( {A}_{1}\right)  = P\left( {B}_{1}\right)  = {0.5},P\left( {{A}_{2}|{A}_{1}}\right)  = {0.8},P\left( {{A}_{2}|{B}_{1}}\right)  = {0.6}$ ,
            
            则由全概率公式可得 $P\left( {A}_{2}\right)  = P\left( {A}_{1}\right) P\left( {{A}_{2}|{A}_{1}}\right)  + P\left( {B}_{1}\right) P\left( {{A}_{2}|{B}_{1}}\right)  = {0.7}$ .
        \end{jieda}
        
        \item 已知一道解答题共有两小问,第一问 5 分,第二问 8 分,现每 10 个人有 6 个人能够解答出第一问, 在第一问解答不出的情况下, 解答出第二问的概率为 $0.1$ , 第一问解答出来的情况下，第二问解答不出来的概率为 0.7 ，则解答出第二问的概率为\blkda{$0.22$}
        \begin{jieda}
            设“解出第一问”为事件 $A$ ，“解出第二问”为事件 $B$
            
            $P\left( A\right)  = \dfrac{6}{10} = {0.6},P\left( {B|\bar{A}}\right)  = {0.1},P\left( {\bar{B}|A}\right)  = {0.7}$，则 $P\left( \bar{A}\right)  = {0.4},P\left( {B|A}\right)  = {0.3}$
            
            所以 $P\left( {AB}\right)  = P\left( A\right) P\left( {B|A}\right)  = {0.18}$，$P\left( {\bar{A}B}\right)  = P\left( \bar{A}\right) P\left( {B|\bar{A}}\right)  = {0.04}$
            
            所以 $P\left( B\right)  = P\left( {AB}\right)  + P\left( {\bar{A}B}\right)  = {0.22}$ .
        \end{jieda}
        
        \item 已知某羽毛球小组共有 20 名运动员，其中一级运动员 4 人，二级运动员 6 人，三级运动员 10 人. 现在举行一场羽毛球选拔赛，若一级、二级、三级运动员能够晋级的概率分别为 0.9, 0.6, 0.2, 则这 20 名运动员中任选一名运动员能够晋级的概率为\blkda{$0.46$}
        \begin{jieda}
            设事件 $B$ 为“选出的运动员能晋级”,
            
            ${A}_{1}$ 为“选出的运动员是一级运动员”,
            
            ${A}_{2}$ 为“选出的运动员是二级运动员”,
            
            ${A}_{3}$ 为“选出的运动员是三级运动员”,
            
            则 $P\left( {B \mid  {A}_{1}}\right)  = {0.9},P\left( {B \mid  {A}_{2}}\right)  = {0.6},P\left( {B \mid  {A}_{3}}\right)  = {0.2}$ ,
            
            又根据题意可得 $P\left( {A}_{1}\right)  = \dfrac{4}{20} = \dfrac{1}{5},P\left( {A}_{2}\right)  = \dfrac{6}{20} = \dfrac{3}{10},P\left( {A}_{3}\right)  = \dfrac{10}{20} = \dfrac{1}{2}$ ,
            
            $\therefore$ 由全概率公式可得:
            
            $P\left( B\right)  = P\left( {A}_{1}\right) P\left( {B \mid  {A}_{1}}\right)  + P\left( {A}_{2}\right) P\left( {B \mid  {A}_{2}}\right)  + P\left( {A}_{3}\right) P\left( {B \mid  {A}_{3}}\right)  = \dfrac{4}{20} \times  {0.9} + \dfrac{6}{20} \times  {0.6} + \dfrac{10}{20} \times  {0.2} = {0.46}$ , $\therefore$ 任选一名运动员能够晋级的概率为 0.46 .
        \end{jieda}
        
    \end{enumerate}
    \item 利用贝叶斯公式求概率
    \begin{enumerate}
        \item 已知事件 $A$ 发生的概率为 0.4，事件 $B$ 发生的概率为 0.5，若在事件 $B$ 发生的条件下,事件 $A$ 发生的概率为 0.6,则在事件 $A$ 发生的条件下,事件 $B$ 发生的概率为\blkda{$0.75$}
        \begin{jieda}
            因为 $P\left( A\right)  = {0.4},P\left( B\right)  = {0.5},P\left( {A|B}\right) = {0.6}$
            
            所以 $P\left( {B|A}\right)  = \dfrac{P\left( {A|B}\right) P(B) }{P\left( A\right) } = \dfrac{0.6 \times 0.5}{0.4} = {0.75}$ .
        \end{jieda}
        \item 有 3 台车床加工同一型号的零件, 第 1 台加工的次品率为 5\%，第 2，3 台加工的次品率均为 3\%，加工出来的零件混放在一起. 已知第 1,2,3 台车床加工的零件数分别占总数的 20\%，30\%，50\%. 如果取到的零件是次品，则它是第 3 台车床加工的概率是\blkda{$\dfrac{15}{34}$}
        \begin{jieda}
            记取到“第1,2,3台车床加工的零件”分别为事件 ${A}_{1},{A}_{2},{A}_{3}$，“取到次品”为事件 $B$
            
            故 $P\left( {A}_{1}\right)  = \dfrac{1}{5},P\left( {A}_{2}\right)  = \dfrac{3}{10},P\left( {A}_{3}\right)  = \dfrac{1}{2}$，$P\left( {B|{A}_{1}}\right)  = \dfrac{1}{20},P\left( {B|{A}_{2}}\right)  = \dfrac{3}{100},P\left( {B|{A}_{3}}\right)  = \dfrac{3}{100},$
            
            由全概率公式可得: $P\left( B\right)  = P\left( {A}_{1}\right)  \times  P\left( {B|{A}_{1}}\right)  + P\left( {A}_{2}\right)  \times  P\left( {B|{A}_{2}}\right)  + P\left( {A}_{3}\right)  \times  P\left( {B|{A}_{3}}\right)  = \dfrac{17}{500}$
            
            由贝叶斯公式: $P\left( {{A}_{3}|B}\right)  = \dfrac{P\left( {A}_{3}\right) P\left( {B|{A}_{3}}\right) }{P\left( B\right) } = \dfrac{15}{34}$
        \end{jieda}
        
        \item 某工厂有 $A,B$ 两个生产车间，所生产的同一批产品合格率分别是 ${98}\%$ 和 ${99}\%$ ，已知某批产品的 60\%和 40\%分别是 $A,B$ 两个车间生产，质量跟踪小组从中随机抽取一件，发现不合格，则该产品是由 $A$ 车间生产的概率为\blkda{$\dfrac{3}{4}$}
        \begin{jieda}
            依题意,该产品是由 $A$ 车间生产的概率为: $\dfrac{{0.6} \times  \left( {1 - {0.98}}\right) }{{0.6} \times  \left( {1 - {0.98}}\right)  + {0.4} \times  \left( {1 - {0.99}}\right) } = \dfrac{{0.6} \times  {0.02}}{{0.6} \times  {0.02} + {0.4} \times  {0.01}} = \dfrac{3}{4}.$
        \end{jieda}
        
        \item 托马斯·贝叶斯(ThomasBayes)在研究“逆向概率”的问题中得到了一个公式: $P\left( {A|B}\right)  = \dfrac{P\left( {B|A}\right)  \cdot  P\left( A\right) }{P\left( {B|A}\right)  \cdot  P\left( A\right)  + P\left( {B|\bar{A}}\right)  \cdot  P\left( \bar{A}\right) }$ . 这个定理在实际生活中有着重要的应用价值. 假设某种疾病在所有人群中的感染率是 0.1\%，医院现有的技术对于该疾病检测召回率为 99\%，即已知患病情况下， 99\% 的可能性可以检查出阳性. 如果从人群中随机抽一个人去检测，经计算检测结果为阳性的全概率为 0.01098 ，请你用这个公式估计医院技术的准确率，即在医院给出的检测结果为阳性的条件下这个人得病的概率为\dotfill{(\kh{C})}
        \xx{$0.1\%$}{$8\%$}{$9\%$}{$99\%$}
        \begin{jieda}
            记一个人得病为事件 $A$ ,检测结果为阳性为事件 $B$
            
            则 $P\left( A\right)  = {0.1}\% ,P\left( {B|A}\right)  = {99}\% ,P(B) = P\left( {B|A}\right)  \cdot  P\left( A\right)  + P\left( {B|\bar{A}}\right)  \cdot  P\left( \bar{A}\right)  = {0.01098}$ ,
            
            所以 $P\left( {A|B}\right)  = \dfrac{P\left( {B|A}\right)  \cdot  P\left( A\right) }{P(B)} = \dfrac{{99}\%  \times  {0.1}\% }{0.01098} \approx  9\%$ 
        \end{jieda}
        
        \item 假设甲袋中有 3 个白球和 3 个红球，乙袋中有 2 个白球和 2 个红球. 现从甲袋中任取 2 个球放入乙袋，再从乙袋中任取 2 个球. 已知从乙袋中取出的是 2 个红球，则从甲袋中取出的也是 2 个红球的概率为\blkda{$\dfrac{3}{8}$}
        \begin{jieda}
            设从甲中取出 2 个球,其中红球的个数为 $i$ 个的事件为 ${A}_{i}$ ,事件 ${A}_{i}$ 的概率为 $P\left( {A}_{i}\right)$，从乙中取出 2 个球,其中红球的个数为 2 个的事件为 $B$ ,事件 $B$ 的概率为 $P\left( B\right)$ ,由题意:
            \begin{dingautolist}{172}
                \item $P\left( {A}_{0}\right)  = \dfrac{{{C}}_{5}^{2}{{C}}_{5}^{0}}{{{C}}_{6}^{2}} = \dfrac{1}{5},P\left( {B|{A}_{0}}\right)  = \dfrac{{{C}}_{5}^{2}{{C}}_{4}^{0}}{{{C}}_{6}^{2}} = \dfrac{1}{15}$
                \item $P\left( {A}_{1}\right)  = \dfrac{{{C}}_{3}^{1}{{C}}_{3}^{1}}{{{C}}_{6}^{2}} = \dfrac{3}{5}P\left( {B|{A}_{1}}\right)  = \dfrac{{{C}}_{3}^{2}{{C}}_{3}^{0}}{{{C}}_{6}^{2}} = \dfrac{1}{5}$
                \item $P\left( {A}_{2}\right)  = \dfrac{{{C}}_{3}^{0}{{C}}_{3}^{2}}{{{C}}_{6}^{2}} = \dfrac{1}{5}P\left( {B|{A}_{2}}\right)  = \dfrac{{{C}}_{2}^{2}{{C}}_{2}^{0}}{{{C}}_{6}^{2}} = \dfrac{2}{5}$
            \end{dingautolist}

            $P\left( B\right)  = P\left( {A}_{0}\right) P\left( {B|{A}_{0}}\right)  + P\left( {A}_{1}\right) P\left( {B|{A}_{1}}\right)  + P\left( {A}_{2}\right) P\left( {B|{A}_{2}}\right) = \dfrac{1}{5} \times  \dfrac{1}{15} + \dfrac{3}{5} \times  \dfrac{1}{5} + \dfrac{1}{5} \times  \dfrac{2}{5} = \dfrac{16}{75}$
            
            $P\left( {{A}_{2}|B}\right)  = \dfrac{P\left( {B{A}_{2}}\right) }{P\left( B\right) } = \dfrac{P\left( {A}_{2}\right) P\left( {B|{A}_{2}}\right) }{P\left( B\right) } = \dfrac{3}{8}$
            
            即已知从乙袋中取出的是 2 个红球,则从甲袋中取出的也是 2 个红球的概率为 $\dfrac{3}{8}$
        \end{jieda}
    \end{enumerate}
    \item 条件概率与全概率公式的综合应用
    \begin{enumerate}
        \item 某运动队为评估短跑运动员在接力赛中的作用，对运动员进行数据分析，运动员甲在接力赛中跑第一棒、第二棒、第三棒、第四棒四个位置，统计以往多场比赛，其出场率与出场时比赛获胜率如下表所示.
        
        \begin{center}
        \begin{tabular}{|c|c|c|c|c|}
        \hline
        比赛位置 & 第一棒 & 第二棒 & 第三棒 & 第四棒 \\
        \cline{1-5}
        出场率 & 0.2 & 0.3 & 0.3 & 0.2 \\
        \cline{1-5}
        比赛胜率 & 0.6 & 0.8 & 0.7 & 0.7 \\
        \cline{1-5}
        \hline
        \end{tabular}
        \end{center}
        \begin{enumerate}
            \item 当甲出场比赛时，求该运动队获胜的概率
            \item 当甲出场比赛时，在该运动队获胜的条件下，求甲跑第四棒的概率
        \end{enumerate}
        \begin{jieda}
            \begin{enumerate}
                \item 记“甲跑第一棒”为事件 ${A}_{1}$ ，“甲跑第二棒”为事件 ${A}_{2}$ ，“甲跑第三棒”为事件 ${A}_{3}$ ，“甲跑第四棒”为事件 ${A}_{4}$ ,“运动队获胜”为事件 $B$ .
                
                $P\left( B\right)  = P\left( {A}_{1}\right) P\left( {B \mid  {A}_{1}}\right)  + P\left( {A}_{2}\right) P\left( {B \mid  {A}_{2}}\right)  + P\left( {A}_{3}\right) P\left( {B \mid  {A}_{3}}\right)  + P\left( {A}_{4}\right) P\left( {B \mid  {A}_{4}}\right) = {0.2} \times  {0.6} + {0.3} \times  {0.8} + {0.3} \times  {0.7} + {0.2} \times  {0.7} = {0.71}.$
                
                \item $P\left( {{A}_{4} \mid  B}\right)  = \dfrac{P\left( {{A}_{4}B}\right) }{P\left( B\right) } = \dfrac{P\left( {A}_{4}\right) P\left( {B \mid  {A}_{4}}\right) }{P\left( B\right) } = \dfrac{{0.2} \times  {0.7}}{0.71} = \dfrac{14}{71}$
            \end{enumerate}
        \end{jieda}
        \item 设甲袋中有 4 个白球和 2 个红球, 乙袋中有 2 个白球和 2 个红球.
        \begin{enumerate}
            \item 现从甲袋中任取 2 个球放入乙袋，再从乙袋中任取 2 个球. 求从乙袋中取出的是 2 个红球的概率
            \item 先随机取一只袋，在再从该袋中先后随机取 2 个球，求第一次取出的是红球的前提下，第二次取出的球是白球的概率
        \end{enumerate}
        \begin{jieda}
            \begin{enumerate}
                \item 记事件 ${A}_1$ :从甲袋中取出 2 个红球， ${A}_{2}$ :从甲袋中取出 2 个白球， ${A}_{3}$ :从甲袋中取出 1 个白球和 1 个红球, $B$ : 从乙袋中取出 2 个红球.
                
                显然， ${A}_{1}$ ， ${A}_{2}$ ， ${A}_{3}$ 两两互斥，且 ${A}_{1} + {A}_{2} + {A}_{3}$ 正好为“从甲袋中任取 2 个球”的样本空间 $\Omega$ .
                
                由全概率公式,得 $P\left( B\right)  = P\left( {A}_{1}\right) P\left( {B \mid  {A}_{1}}\right)  + P\left( {A}_{2}\right) P\left( {B \mid  {A}_{2}}\right)  + P\left( {A}_{3}\right) P\left( {B \mid  {A}_{3}}\right)$
                
                $= \dfrac{{{C}}_{2}^{2}}{{{C}}_{6}^{2}} \cdot  \dfrac{{{C}}_{4}^{2}}{{{C}}_{6}^{2}} + \dfrac{{{C}}_{4}^{2}{{C}}_{2}^{2}}{{{C}}_{6}^{2}} \cdot  \dfrac{{{C}}_{3}^{2}}{{{C}}_{6}^{2}} + \dfrac{{{C}}_{4}^{2}}{{{C}}_{6}^{2}} \cdot  \dfrac{{{C}}_{2}^{2}}{{{C}}_{6}^{2}} = \dfrac{4}{25}$

                \item 设“取出的是甲袋”为事件 ${A}_{1}$ ，“取出的是乙袋”为事件 ${A}_{2}$ ，“第一次取出的球是红球”为事件 $B$ ， “第二次取出的球是白球”为事件 $C$ ，则 $P\left( {A}_{1}\right)  = P\left( {A}_{2}\right)  = \dfrac{1}{2}$
                
                $P\left( {{BC} \mid  {A}_{1}}\right)  = \dfrac{2 \times  4}{6 \times  5} = \dfrac{4}{15},P\left( {{BC} \mid  {A}_{2}}\right)  = \dfrac{2 \times  2}{4 \times  3} = \dfrac{1}{3}.$
                
                故 $P\left( {BC}\right)  = P\left( {A}_{7}\right) P\left( {{BC} \mid  {A}_{1}}\right)  + P\left( {A}_{2}\right) P\left( {{BC} \mid  {A}_{2}}\right)  = \dfrac{1}{2} \times  \dfrac{4}{15} + \dfrac{1}{2} \times  \dfrac{1}{3} = \dfrac{3}{10}$ ,
                
                $P\left( B\right)  = P\left( {A}_{1}\right) P\left( {B \mid  {A}_{1}}\right)  + P\left( {A}_{2}\right) P\left( {B \mid  {A}_{2}}\right)  = \dfrac{1}{2} \times  \dfrac{1}{3} + \dfrac{1}{2} \times  \dfrac{1}{2} = \dfrac{5}{12}$
                
                所以 $P\left( {C \mid  B}\right)  = \dfrac{P\left( {BC}\right) }{P\left( B\right) } = \dfrac{18}{25}$
            \end{enumerate}
        \end{jieda}
        \item 两批同种规格的产品，第一批占 40\%，次品率 3\%，第二批占 60\%，次品率为 $p$ .将两批产品混合，从混合产品中任取 1 件是合格品的概率为 97.6\%.
        \begin{enumerate}
            \item 求 $p$
            \item 已知取到的是次品，求它取自第二批产品的概率
        \end{enumerate}
        \begin{jieda}
            \begin{enumerate}
                \item 记事件 $A$ :任取一件产品是次品，记事件 ${B}_{i}$ :第 $i$ 批的产品， $i = {1.2}$ .

                则 $P\left( {A \mid  {B}_{1}}\right)  = 3\% ,P\left( {A \mid  {B}_{2}}\right)  = p,P\left( {B}_{1}\right)  = {0.40},P\left( {B}_{2}\right)  = {0.60}$ ,
                
                由 $P\left( A\right)  = P\left( {A{B}_{1}}\right)  + P\left( {A{B}_{2}}\right)  = {0.4} \times  {0.03} + {0.6} \times  p = 1 - {0.976}$ ,解得 $p = {0.02}$

                \item $P\left( {{B}_{2} \mid  A}\right)  = \dfrac{P\left( {A{B}_{2}}\right) }{P\left( A\right) } = \dfrac{P\left( {A \mid  {B}_{2}}\right) P\left( {B}_{2}\right) }{P\left( A\right) } = \dfrac{{0.6} \times  {0.02}}{1 - {0.976}} = \dfrac{1}{2}$
                
                已知取到的是次品，则它取自第二批产品的概率 $\dfrac{1}{2}$
            \end{enumerate}
        \end{jieda}
        
        \item 甲袋中有 3 个白球和 2 个红球, 乙袋中有 2 个白球和 3 个红球. 先随机取一只袋, 再从该袋中先后随机取 2 个球.
        \begin{enumerate}
            \item 求随机取到的是甲袋且从中取出的两球均为白球的概率
            \item 求第一次取出的是白球的概率
            \item 求第一次取出的是白球的前提下，第二次取出的依然是白球的概率
        \end{enumerate}
        \begin{jieda}
            \begin{enumerate}
                \item 记“随机取到甲袋”为事件 ${A}_{1}$ ,“随机取到乙袋”为事件 ${A}_{2}$ ,“第一次取出的是白球”为事件 $B$ ，“第二次取出的是白球”为事件 $C$ .
                
                $P\left( {{A}_{1}{BC}}\right)  = P\left( {A}_{1}\right) P\left( {{BC} \mid  {A}_{1}}\right)  = \dfrac{1}{2} \times  \dfrac{{{A}}_{3}^{2}}{{{A}}_{5}^{2}} = \dfrac{3}{20}.$
                
                所以取到甲袋且从中取出的两球均为白球的概率为 $\dfrac{3}{20}$ .

                \item $P\left( B\right)  = P\left( {A}_{1}\right) P\left( {B \mid  {A}_{1}}\right)  + P\left( {A}_{2}\right) P\left( {B \mid  {A}_{2}}\right)  = \dfrac{1}{2} \times  \dfrac{3}{5} + \dfrac{1}{2} \times  \dfrac{2}{5} = \dfrac{1}{2}$

                所以第一次取到白球的概率为 $\dfrac{1}{2}$

                \item $P\left( {BC}\right)  = P\left( {A}_{1}\right) P\left( {{BC} \mid  {A}_{1}}\right)  + P\left( {A}_{2}\right) P\left( {{BC} \mid  {A}_{2}}\right)  = \dfrac{1}{2} \times  \dfrac{{A}_{3}^{2}}{{A}_{5}^{2}} + \dfrac{1}{2} \times  \dfrac{{A}_{2}^{2}}{{A}_{5}^{2}} = \dfrac{1}{5}$
                
                所以 $P\left( {C \mid  B}\right)  = \dfrac{P\left( {BC}\right) }{P\left( B\right) } = \dfrac{2}{5}$
                
                所以第一次取出的是白球的前提下，第二次取出的依然是白球的概率为 $\dfrac{2}{5}$
            \end{enumerate}
        \end{jieda}
\end{enumerate}

\item 条件概率与其他知识综合
\begin{enumerate}
    \item 我们平时常用的视力表叫做对数视力表, 视力呈现为 4.8, 4.9, 5.0, 5.1. 视力 5.0 为正常视力. 否则就是近视. 某校进行一次对学生视力与学习成绩的相关调查，随机抽查了 100 名近视学生的成绩 (按照各科占一定权重计算而得的满分 100 分的综合成绩), 得到频率分布直方图如下:
    
    \begin{center}
    \includegraphics[width=0.3\textwidth]{bodyxb/src/bo_d67fmsv7aajc739ac5b0_9_140_961_449_267_0.jpg}
    \end{center}
    \begin{enumerate}
        \item 估计该校近视学生学习成绩的第 85 百分位数；(精确到 0.1)
        
        \item 已知该校学生的近视率为 54\%，学生成绩的优秀率为 36\% (成绩≥ 85 分视作优秀)，从该校学生中任选一人，若此人的成绩为优秀，求此人近视的概率. (以样本中的频率作为相应的概率)
    \end{enumerate}
    \begin{jieda}
        \begin{enumerate}
            \item 由频率分布直方图可知，成绩 90 分以下所占比例为 $7\%  + {13}\%  + {20}\%  + {24}\%  = {64}\%$ , 因此第 85 百分位数一定位于 $\left\lbrack  {{90},{100}}\right\rbrack$ 内,由 ${90} + \dfrac{{0.85} - {0.64}}{0.036} \approx  {95.8}$，可以估计该地区近视学生的学习成绩的第 85 百分位数约为 95.8

            \item 设 $A =$ “该地区近视学生”， $B =$ “该地区优秀学生”，
            
            由频率分布直方图可得 $P\left( {B \mid  A}\right)  = \left( {\dfrac{0.024}{2} + {0.036}}\right)  \times  {10} = {0.48}$ ,
            
            又 $P\left( A\right)  = {0.54},P\left( B\right)  = {0.36}$ ,
            
            所以 $P\left( {A \mid  B}\right)  = \dfrac{P\left( {AB}\right) }{P\left( B\right) } = \dfrac{P\left( {B \mid  A}\right) P\left( A\right) }{P\left( B\right) } = \dfrac{{0.48} \times  {0.54}}{0.36} = {0.72}$ ,
            
            即若此人的成绩为优秀，则此人近视的概率为 0.72 .
        \end{enumerate}
    \end{jieda}

    \item 将一颗质地均匀的正四面体骰子 (四个面的点数分别为 1, 2,3,4) 先后抛掷两次，记第一次出现的点数为 $x$ ，第二次出现的点数为 $y$
    \begin{enumerate}
        \item 记事件 $A$ 为“ $x + y \leq  5$ ”，求 $P\left( A\right)$
        
        \item 记事件 $B$ 为“ $\left| {x - y}\right|  = 3$ ” ，求 $P\left( {B \mid  A}\right)$
    \end{enumerate}
    \begin{jieda}
        \begin{enumerate}
            \item 投掷骰子 2 次得到的所有结果共 $4\times4 = 16$ 种
            
            事件 $A$ 包含的结果有: $\left( {1,1}\right) ,\left( {1,2}\right) ,\left( {1,3}\right) ,\left( {1,4}\right) ,\left( {2,1}\right) ,\left( {2,2}\right) ,\left( {2,3}\right) ,\left( {3,1}\right) ,\left( {3,2}\right) ,\left( {4,1}\right)$ 共 10 种
            
            则 $P\left( A\right)  = \dfrac{10}{16} = \dfrac{5}{8}$

            \item 事件 ${AB}$ 包含的结果有:$( 1,4)$，$(4,1)$ 共 2 种，则 $P\left( {AB}\right)  = \dfrac{2}{16} = \dfrac{1}{8}$ ，
            
            $P\left( {B \mid  A}\right)  = \dfrac{P\left( {AB}\right) }{P\left( A\right) } = \dfrac{\dfrac{1}{8}}{\dfrac{5}{8}} = \dfrac{1}{5}$
        \end{enumerate}
    \end{jieda}
    \item 我们平时常用的视力表叫做对数视力表，视力呈现为 4.8，4.9，5.0,5.1. 视力 $\geq$ 5.0 为正常视力. 否则就是近视. 某地区对学生视力与学习成绩进行调查，随机抽查了 100 名近视学生的成绩，得到频率分布直方图:
    
    \begin{center}
    \includegraphics[width=0.3\textwidth]{bodyxb/src/bo_d67fmsv7aajc739ac5b0_10_139_1487_448_371_0.jpg}
    \end{center}
    \begin{enumerate}
        \item 能否据此判断学生的学习成绩与视力状况相关；(不需说明理由)
        \item 估计该地区近视学生学习成绩的第 85 百分位数；(精确到 0.1)
        \item 已知该地区学生的近视率为 54\%，学生成绩的优秀率为 36\% (成绩≥ 85 分为优秀)，从该地区学生中任选一人，若此人的成绩为优秀，求此人近视的概率. (以样本中的频率作为相应的概率)
    \end{enumerate}
    \begin{jieda}
        \begin{enumerate}
            \item 因从题干频率分布直方图不可以看到，不同成绩层次的同学近视率的情况，故不能据此判断学生的学习成绩与视力状况相关
            
            \item 由频率分布直方图可知，成绩 90 分以下所占比例为 $7\%  + {13}\%  + {20}\%  + {24}\%  = {64}\%$ ,因此第 85 百分位数一定位于 $\left\lbrack  {{90},{100}}\right\rbrack$ 内，由 ${90} + {10} \times  \dfrac{{85} - {64}}{{100} - {64}} = {90} + \dfrac{35}{6} \approx  {95.8}$ ,可以估计该地区近视学生的学习成绩的第 85 百分位数约为 95.8
            
            \item 设 $A =$ “该地区近视学生”， $B =$ “该地区优秀学生”，由频率分布直方图可得 $P\left( {B \mid  A}\right)  = \left( {\dfrac{0.024}{2} + {0.036}}\right)  \times  {10} = {0.48}$ , $P\left( A\right)  = {0.54},P\left( B\right)  = {0.36}$ ,所以 $P\left( {A \mid  B}\right)  = \dfrac{P\left( {AB}\right) }{P\left( B\right) } = \dfrac{P\left( {B \mid  A}\right) P\left( A\right) }{P\left( B\right) } = \dfrac{{0.48} \times  {0.54}}{0.36} = {0.72}$ .
            
            即若此人的成绩为优秀，则此人近视的概率为 0.72 .
        \end{enumerate}
    \end{jieda}
    \item 某学校高二 1 班有五名学生报名参加社团活动，社团活动共有 “记者在线”、“机器人行动”、“音乐之声”三个项目，每人都要报名且限报其中一项.
    \begin{enumerate}
        \item 求“每个项目都有人报名”的报名情况种数
        \item 已知其中一项目恰只有三名学生报名，求只有甲同学一人报“记者在线”的概率
    \end{enumerate}
    \begin{jieda}
        \begin{enumerate}
            \item “每个项目都有人报名”，则 5 名学生分三组，即人数分为 3,1，1 或 2,2，1； 故此时报名情况有 ${{C}}_{5}^{3}{{A}}_{3}^{3} + \dfrac{{{C}}_{5}^{2}{{C}}_{3}^{2}}{{{A}}_{2}^{2}}{{A}}_{3}^{3} = {150}$ 种.
            \item 记事件A为“其中一项目恰只有三名学生报名”，事件B为“只有甲同学一人报记者在线”，
            
            事件 $A$ 为“其中一项目恰只有三名学生报名”，报名情况有 ${{C}}_{5}^{3}{{A}}_{3}^{3} + {{C}}_{5}^{3}{{C}}_{3}^{2}{{A}}_{2}^{2} = {120}$ 种，
            
            所以 $P\left( A\right)  = \dfrac{120}{{3}^{5}}$ ,
            
            若 $A,B$ 同时发生,即其中一项目恰只有三名学生报名,且只有甲同学一人报“记者在线”,则有 ${{C}}_{4}^{3}{{A}}_{2}^{2} = 8$ 种, 所以 $P\left( {AB}\right)  = \dfrac{8}{{3}^{5}}$ ,所以 $P\left( {B|A}\right)  = \dfrac{P\left( {AB}\right) }{P\left( A\right) } = \dfrac{1}{15}$
        \end{enumerate}
    \end{jieda}
    \end{enumerate}
    \item 离散型随机变量的判断
    \begin{enumerate}
        \item 下面给出的四个随机变量中是离散型随机变量的是\blkda{\ding{172}\ding{175}}
        \begin{dingautolist}{172}
            \item 某食堂在中午半小时内进的人数 ${Z}_{1}$
            \item 某元件的测量误差 ${Z}_{2}$
            \item 小明在一天中浏览网页的时间 ${Z}_{3}$
            \item 高一 2 班参加运动会的人数 ${Z}_{4}$
        \end{dingautolist}
        \begin{jieda}
            \begin{dingautolist}{172}
                \item 某食堂在中午半小时内进的人数 ${Z}_{1}$ 可以一一列举出来，故 ${Z}_{1}$ 是离散型随机变量

                \item 某元件的测量误差 ${Z}_{2}$ 不能一一列举出来，故 ${Z}_{2}$ 不是离散型随机变量

                \item 小明在一天中浏览网页的时间 ${Z}_{3}$ 不能一一列举出来，故 ${Z}_{3}$ 不是离散型随机变量

                \item 高一 2 班参加运动会的人数 ${Z}_{4}$ 可以一一列举出来，故 ${Z}_{4}$ 是离散型随机变量
            \end{dingautolist}
        \end{jieda}
        
        \item 袋中有 2 个黑球、 5 个红球，从中任取 2 个，可以作为随机变量的是\dotfill{(\kh{B})}
        \xx{取到的球的个数}{取到红球的个数}{至少取到一个红球}{至少取到一个红球的概率}
        
        \item 下列叙述中，是离散型随机变量的为\dotfill{(\kh{C})}
        \xx{将一枚质地均匀的硬币掷五次，出现正面和反面向上的次数之和}{某人早晨在车站等出租车的时间}{连续不断地射击，首次命中目标所需要的次数}{袋中有 2 个黑球 6 个红球，任取 2 个，取得一个红球的可能性}
        
        \item 一串钥匙有 6 把，只有一把能打开锁，依次试验，打不开的扔掉，直到找到能开锁的钥匙为止，则试验次数 $X$ 的可能取值为\dotfill{(\kh{D})}
        \xx{$1,2,3,\cdots ,6$}{$0,1,2,\cdots ,6$}{$0,1,2,\cdots ,5$}{$1,2,3,\cdots ,5$}
        \begin{jieda}
            由试验次数 $X$ 的含义可知,至少试验一次才可能刚好打开，如果第五次依然没有打开，此时不管开锁是否成功，都能确定能开锁的钥匙.
            
            所以 $X$ 的所有可能取值为: 1,2,3,4,5.
        \end{jieda}
    \end{enumerate}
    \item 求离散型随机变量的分布列、随机变量分布列的性质
    \begin{enumerate}
        \item 投掷两枚质地均匀的骰子，记偶数点朝上的骰子的个数为 $X$ ，则 $X$ 的分布列为\dotfill{(\kh{C})}
        \xx{\begin{tabular}{|c|c|c|}
        \hline
        $X$ & 1 & 2 \\
        \cline{1-3}
        $P$ & $\dfrac{1}{2}$ & $\dfrac{1}{2}$ \\
        \cline{1-3}
        \hline
        \end{tabular}}{\begin{tabular}{|c|c|c|}
        \hline
        $X$ & 0 & 1 \\
        \cline{1-3}
        $P$ & $\dfrac{1}{2}$ & $\dfrac{1}{2}$ \\
        \cline{1-3}
        \hline
        \end{tabular}}{\begin{tabular}{|c|c|c|c|}
        \hline
        $X$ & 0 & 1 & 2 \\
        \cline{1-4}
        $P$ & $\dfrac{1}{4}$ & $\dfrac{1}{2}$ & $\dfrac{1}{4}$ \\
        \cline{1-4}
        \hline
        \end{tabular}}{\begin{tabular}{|c|c|c|c|}
        \hline
        $X$ & 0 & 1 & 2 \\
        \cline{1-4}
        $P$ & $\dfrac{1}{2}$ & $\dfrac{1}{4}$ & $\dfrac{1}{4}$ \\
        \cline{1-4}
        \hline
        \end{tabular}}
        \item 甲、乙两名篮球运动员每次投篮的命中率分别为 0.8, 0.7，他们各自投篮 1 次，设两人命中总次数为 $X$ ，则 $X$ 的分布列为\dotfill{(\kh{D})}
        \xx{\begin{tabular}{|c|c|c|c|}
        \hline
        $X$ & 0 & 1 & 2 \\
        \cline{1-4}
        $P$ & 0.08 & 0.14 & 0.78 \\
        \cline{1-4}
        \hline
        \end{tabular}}{\begin{tabular}{|c|c|c|c|}
        \hline
        $X$ & 0 & 1 & 2 \\
        \cline{1-4}
        $P$ & 0.06 & 0.24 & 0.70 \\
        \cline{1-4}
        \hline
        \end{tabular}}{\begin{tabular}{|c|c|c|c|}
        \hline
        $X$ & 0 & 1 & 2 \\
        \cline{1-4}
        $P$ & 0.06 & 0.56 & 0.38 \\
        \cline{1-4}
        \hline
        \end{tabular}}{\begin{tabular}{|c|c|c|c|}
        \hline
        $X$ & 0 & 1 & 2 \\
        \cline{1-4}
        $P$ & 0.06 & 0.38 & 0.56 \\
        \cline{1-4}
        \hline
        \end{tabular}}
        \item 盒中有四张卡片，分别标有数字 1, 2, 3, 4, 现从盒中任取两张卡片，记取到偶数的个数为 $X$ .
        \begin{enumerate}
            \item 求 $P\left( {X = 1}\right)$
            \item 求 $X$ 的分布列
        \end{enumerate}
        \begin{jieda}
            \begin{enumerate}
                \item $X = 1$ 表示的随机事件是“取到的两张卡片上的数字是一个偶数、一个奇数”，所以 $P\left( {X = 1}\right)  = \dfrac{{{C}}_{2}^{1} \times  {{C}}_{2}^{1}}{{{C}}_{4}^{2}} = \dfrac{2}{3}$ .
                \item 依题意 $X$ 的可能取值为 0，1，2，则 $P\left( {X = 0}\right)  = \dfrac{{{C}}_{2}^{2}}{{{C}}_{4}^{2}} = \dfrac{1}{6},P\left( {X = 1}\right)  = \dfrac{{{C}}_{2}^{1} \times  {{C}}_{2}^{1}}{{{C}}_{4}^{2}} = \dfrac{2}{3},P\left( {X = 2}\right)  = \dfrac{{{C}}_{2}^{2}}{{{C}}_{4}^{2}} = \dfrac{1}{6}$ ，所以 $X$ 的分布列如下所示:
                
                \begin{center}
                
                \begin{tabular}{|c|c|c|c|}
                \hline
                $X$ & 0 & 1 & 2 \\
                \cline{1-4}
                $P$ & 1 & 2 & 1 6 \\
                \cline{1-4}
                \hline
                \end{tabular}
                
                \end{center}
            \end{enumerate}
        \end{jieda}

        
        \item 下表是离散型随机变量 $X$ 的概率分布，则常数 $a$ 的值是\blkda{$\dfrac{1}{9}$}
        
        \begin{center}
        
        \begin{tabular}{|c|c|c|c|c|}
        \hline
        $X$ & 3 & 4 & 5 & 6 \\
        \cline{1-5}
        $P$ & $\dfrac{a}{2}$ & $\dfrac{1}{6}+a$ & $\dfrac{1}{2}$ & $\dfrac{1}{6}$ \\
        \cline{1-5}
        \hline
        \end{tabular}
        
        \end{center}
        \begin{jieda}
            由 $\dfrac{a}{2} + \dfrac{1}{6} + a + \dfrac{1}{2} + \dfrac{1}{6} = 1$ ,解得 $a = \dfrac{1}{9}$
        \end{jieda}

        \item 下表是离散型随机变量 $X$ 的分布列，则常数 $q$ 的值是\blkda{$0.2$}
        
        \begin{center}
        
        \begin{tabular}{|c|c|c|c|}
        \hline
        $X$ & 0 & 1 & 2 \\
        \cline{1-4}
        $P$ & $0.36$ & $1-2q$ & ${q}^{2}$ \\
        \cline{1-4}
        \hline
        \end{tabular}
        \end{center}
        \begin{jieda}
            由已知得 ${0.36} + 1 - {2q} + {q}^{2} = 1$ ,解得 $q = {0.2}$ 或 $q = {1.8}$ (舍去).
        \end{jieda}
        
        \item 已知离散型随机变量 $X$ 的分布列为 $P\left( {X = n}\right)  = \dfrac{a}{n\left( {n + 1}\right) }(n = 1$ , 2,3),则 $a =$\blkda{$\dfrac{4}{3}$}
        \begin{jieda}
            因为 $P\left( {X = 1}\right)  + P\left( {X = 2}\right)  + P\left( {X = 3}\right)  = 1$
            
            所以 $\dfrac{a}{1 \times  2} + \dfrac{a}{2 \times  3} + \dfrac{a}{3 \times  4} = a\left( {1 - \dfrac{1}{2} + \dfrac{1}{2} - \dfrac{1}{3} + \dfrac{1}{3} - \dfrac{1}{4}}\right)  = a\left( {1 - \dfrac{1}{4}}\right)  = 1$ ,
            
            所以 $a = \dfrac{4}{3}$ .
        \end{jieda}

        
        \item 已知随机变量 $X$ 的分布列为$P\left( {X = n}\right)  = \dfrac{a}{\left( {n + 1}\right) \left( {n + 2}\right) },\left( {n = 0,1,2}\right)$ ,其中 $a$ 是常数,则下列说法不正确的是\dotfill{(\kh{D})}
        \xx{$P\left( {X = 0}\right)  + P\left( {X = 1}\right)  + P\left( {X = 2}\right)  = 1$}{$a = \dfrac{4}{3}$}{$P\left( {0 \leq  X < 2}\right)  = \dfrac{8}{9}$}{$P\left( {X \geq  1}\right)  = \dfrac{4}{9}$}
        \begin{jieda}
            由 $P\left( {X = n}\right)  = \dfrac{a}{\left( {n + 1}\right) \left( {n + 2}\right) },\left( {n = 0,1,2}\right)$ ,
            
            得 $P\left( {X = 0}\right)  + P\left( {X = 1}\right)  + P\left( {X = 2}\right)  = 1$ ,
            
            即 $\dfrac{a}{2} + \dfrac{a}{6} + \dfrac{a}{12} = 1$ ，解得 $a = \dfrac{4}{3}$ ，故 ${{AB}}$ 正确；
            
            $P\left( {0 \leq  X < 2}\right)  = P\left( {X = 0}\right)  + P\left( {X = 1}\right)  = \dfrac{2}{3} + \dfrac{2}{9} = \dfrac{8}{9}$ 故 C 正确；
            
            $P\left( {X \geq  1}\right)  = P\left( {X = 1}\right)  + P\left( {X = 2}\right)  = \dfrac{2}{9} + \dfrac{1}{9} = \dfrac{1}{3}$ ，故 D 错误.
        \end{jieda}
        
    \end{enumerate}
    \item 由随机变量的分布列求概率、两个相关的随机变量的分布列
    \begin{enumerate}
         \item 公园的某个位置摆放了 10 盆牡丹花,编号分别为 $0,1,2,3,\cdots$ , 9，若从中任取 1 盆，则编号“大于 5”的概率是\blkda{$\dfrac{2}{5}$}
         \begin{jieda}
             设任取 1 盆的编号为随机变量 $X$，则 $X$ 的可能取值为 $0,1,2,\cdots ,9$，且 $P\left( {X = 0}\right)  = P\left( {X = 1}\right)  = P\left( {X = 2}\right)  = \cdots  = P\left( {X = 9}\right)  = \dfrac{1}{10}$ .
            
            $\therefore P\left( {X > 5}\right)  = P\left( {X = 6}\right)  + P\left( {X = 7}\right)  + P\left( {X = 8}\right)  + P\left( {X = 9}\right)  = \dfrac{4}{10} = \dfrac{2}{5}$
         \end{jieda}

         
        \item 设随机变量 $X$ 的分布列 $P\left( {X = i}\right)  = \dfrac{k}{{2}^{i}}\left( {i = 1,2,3}\right)$ ,则 $P\left( {X \geq  2}\right)$ 的值为\blkda{$\dfrac{3}{7}$}
        \begin{jieda}
            因为随机变量 $X$ 的分布列 $P\left( {X = i}\right)  = \dfrac{k}{{2}^{i}}\left( {i = 1,2,3}\right)$，所以 $\dfrac{k}{{2}^1} + \dfrac{k}{{2}^{2}} + \dfrac{k}{{2}^{3}} = 1$ ,解得: $k = \dfrac{8}{7}$
            
            $P\left( {X \geq  2}\right)  = P\left( {X = 2}\right)  + P\left( {X = 3}\right)  = \dfrac{2}{7} + \dfrac{1}{7} = \dfrac{3}{7}$
        \end{jieda}

        \item 设离散型随机变量 $X$ 的分布列为
        
        \begin{center}
        
        \begin{tabular}{|c|c|c|c|c|c|}
        \hline
        $X$ & 0 & 1 & 2 & 3 & 4 \\
        \cline{1-6}
        $P$ & 0.1 & 0.1 & $m$ & 0.3 & 0.2 \\
        \cline{1-6}
        \hline
        \end{tabular}
        
        \end{center}
        
        \item 若随机变量 $Y = {2X} - 2$ ，则 $P\left( {Y = 2}\right)$ 等于\blkda{$0.3$}
        \begin{jieda}
            由离散型随机变量的分布列的性质得 ${0.1} + {0.1} + m + {0.2} + {0.3} = 1$ ,
            
            解得 $m = {0.3}$ ,
            
            $\because$ 随机变量 $Y = {2X} - 2$ ,
            
            $\therefore P\left( {Y = 2}\right)  = P\left( {X = 2}\right)  = {0.3}$ .
        \end{jieda}

        \item 已知随机变量 $X$ 的分布列如表所示.
        
        \begin{center}
        
        \begin{tabular}{|c|c|c|c|c|c|c|}
        \hline
        $X$ & -2 & -1 & 0 & 1 & 2 & 3 \\
        \cline{1-7}
        $P$ & $\dfrac{1}{2}$ & $\dfrac{1}{4}$ & $\dfrac{1}{3}$ & $\dfrac{1}{2}$ & $\dfrac{1}{6}$ & $\dfrac{1}{12}$ \\
        \cline{1-7}
        \hline
        \end{tabular}
        
        \end{center}
        \begin{enumerate}
            \item 求随机变量 $Y = {X}^{2}$ 的分布列
            \item 若 $P\left( {Y < x}\right)  = \dfrac{11}{12}$ ，求实数 $x$ 的取值范围.
        \end{enumerate}
        \begin{jieda}
            \begin{enumerate}
                \item 由随机变量 $X$ 的分布列知， $Y$ 的可能取值为0,1,4,9,
                
                则 $P\left( {Y = 0}\right)  = P\left( {X = 0}\right)  = \dfrac{1}{3}$
                
                $P\left( {Y = 1}\right)  = P\left( {X =  - 1\text{ 或 }X = 1}\right)  = \dfrac{1}{4} + \dfrac{1}{12} = \dfrac{4}{12} = \dfrac{1}{3}$ ,
                
                $P\left( {Y = 4}\right)  = P\left( {X =  - 2\text{ 或 }X = 2}\right)  = \dfrac{1}{12} + \dfrac{1}{6} = \dfrac{3}{12} = \dfrac{1}{4}$
                
                $P\left( {Y = 9}\right)  = P\left( {X = 3}\right)  = \dfrac{1}{12}$
                
                可得随机变量 $Y$ 的分布列如表所示.
                
                \begin{center}
                
                \begin{tabular}{|c|c|c|c|c|}
                \hline
                $Y$ & 0 & 1 & 4 & 9 \\
                \cline{1-5}
                $P$ & $\dfrac{1}{3}$ & $\dfrac{1}{3}$ & $\dfrac{1}{4}$ & $\dfrac{1}{2}$ \\
                \cline{1-5}
                \hline
                \end{tabular}
                
                \end{center}
                
                \item 因为 $\dfrac{1}{3} + \dfrac{1}{3} + \dfrac{1}{4} = \dfrac{11}{12},\dfrac{1}{3} + \dfrac{1}{3} + \dfrac{1}{4} + \dfrac{1}{12} = 1$ ，
                
                又因为 $P\left( {Y < x}\right)  = \dfrac{11}{12}$ ,所以 $4 < x \leq  9$ .
                
                $\therefore$ 实数 $x$ 的取值范围是 $(4,9\rbrack$ .
            \end{enumerate}
        \end{jieda}
        
        \item 设随机变量 $\xi$ 的分布列如下表，则 $P\left( {\dfrac{{2\xi } - 5}{\xi  - 4} < 1}\right)  =$ \blkda{$\dfrac{7}{12}$}
        
        \begin{center}
        
        \begin{tabular}{|c|c|c|c|c|}
        \hline
        $\xi$ & 1 & 2 & 3 & 4 \\
        \cline{1-5}
        $P$ & $\dfrac{1}{2}$ & $a$ & $\dfrac{1}{3}$ & $\dfrac{1}{3}$ \\
        \cline{1-5}
        \hline
        \end{tabular}
        
        \end{center}
        \begin{jieda}
            根据题意, $\dfrac{{2\xi } - 5}{\xi  - 4} < 1$ ,解得 $1 < \xi  < 4$ ,则 $P\left( {\dfrac{{2\xi } - 5}{\xi  - 4} < 1}\right)  = P\left( {1 < \xi  < 4}\right)$，结合分布列:
            
            $P\left( {1 < \xi  < 4}\right)  = P\left( {\xi  = 2}\right)  + P\left( {\xi  = 3}\right)  = 1 - P\left( {\xi  = 1}\right)  - P\left( {\xi  = 4}\right)  = 1 - \dfrac{1}{12} - \dfrac{1}{3} = \dfrac{7}{12}$ .
        \end{jieda}
    
        \item 设 $X$ 是一个离散型随机变量，其分布列如下，则 $P\left( |X| = 1\right)$ 等于\blkda{$\dfrac{2}{3}$}

        \begin{center}
        
        \begin{tabular}{|c|c|c|c|}
        \hline
        $X$ & $-1$ & 0 & 1 \\
        \cline{1-4}
        $P$ & $\dfrac{1}{3}$ & $1-2q$ & $3{q}^{2} - q  + \dfrac{1}{3}$ \\
        \cline{1-4}
        \hline
        \end{tabular}
        
        \end{center}
        \begin{jieda}
            由离散型随机变量的性质可得 $\dfrac{1}{3} + 1 - {2q} + 3{q}^{2} - q + \dfrac{1}{3} = 1$，即 $\left( {{3q} - 1}\right) \left( {{3q} - 2}\right)  = 0$ ，解得 $q = \dfrac{1}{3}$ 或 $q = \dfrac{2}{3}$
            
            $q = \dfrac{2}{3}$ 时 $1 - {2q} < 0$ ，不合题意， $\therefore q = \dfrac{1}{3}$
            
            $\therefore P\left( |X| = 1\right)  = P\left( {X =  - 1}\right)  + P\left( {X = 1}\right)  = 1 - P\left( {X = 0}\right)  = \dfrac{2}{3}$
        \end{jieda}

        
        \item 设离散型随机变量 $X$ 的分布列为
        
        \begin{center}
        
        \begin{tabular}{|c|c|c|c|c|}
        \hline
        $X$ & 0 & 1 & 2 & 3 \\
        \cline{1-5}
        $P$ & 0.2 & 0.1 & 0.1 & 0.3 \\
        \cline{1-5}
        \hline
        \end{tabular}
        
        \end{center}
        
        若随机变量 $Y = \left| {X - 1}\right|$ ，则 $P\left( {Y = 1}\right)$ 等于\blkda{$0.3$}
        \begin{jieda}
            因为 $Y = \left| {X - 1}\right|$，所以 $P\left( {Y = 1}\right)  = P\left( {X = 0}\right)  + P\left( {X = 2}\right)  = {0.2} + {0.1} = {0.3}$.
        \end{jieda}
        
        \item 设离散型随机变量 $X$ 的分布列为

        \begin{center}
        
        \begin{tabular}{|c|c|c|c|c|c|}
        \hline
        $X$ & 0 & 1 & 2 & 3 & 4 \\
        \cline{1-6}
        $P$ & 0.2 & 0.1 & 0.1 & 0.3 & $m$ \\
        \cline{1-6}
        \hline
        \end{tabular}
        
        \end{center}
        \begin{enumerate}
            \item 求 $\eta  = \left| {X - 1}\right|$ 的分布列
            \item 求 $P\left( {1 < {2X} + 1 < 9}\right)$
        \end{enumerate}
        \begin{jieda}
            \begin{enumerate}
                \item $0.2 + 0.1 + 0.1 + 0.3 + m = 1$ , 故 $m = 0.3$
                
                $\eta$ 的可能取值为 $0\text{ 、 }1\text{ 、 }2\text{ 、 }3$ ,
                
                $P\left( {\eta  = 0}\right)  = P\left( {X = 1}\right)  = {0.1},P\left( {\eta  = 1}\right)  = P\left( {X = 0}\right)  + P\left( {X = 2}\right)  = {0.2} + {0.1} = {0.3}$ ,
                
                $P\left( {\eta  = 2}\right)  = P\left( {X = 3}\right)  = {0.3},P\left( {\eta  = 3}\right)  = P\left( {X = 4}\right)  = {0.3},$
                
                故其分布列为:
                
                \begin{center}
                \begin{tabular}{|c|c|c|c|c|}
                \hline
                $\eta$ & 0 & 1 & 2 & 3 \\
                \cline{1-5}
                $P$ & 0.1 & 0.3 & 0.3 & 0.3 \\
                \cline{1-5}
                \hline
                \end{tabular}
                \end{center}
                \item 由 $1 < {2X} + 1 < 9$ ，可得 $0 < X < 4$，$P\left( {1 < {2X} + 1 < 9}\right)  = P\left( {X = 1}\right)  + P\left( {X = 2}\right)  + P\left( {X = 3}\right)  = {0.1} + {0.1} + {0.3} = {0.5}$ .
            \end{enumerate}
        \end{jieda}
        \item 已知随机变量 $X$ 的分布列为

        \begin{center}
        
        \begin{tabular}{|c|c|c|c|c|c|c|}
        \hline
        $X$ & 1 & 2 & 3 & 4 & 5 & 6 \\
        \cline{1-7}
        $P$ & $\dfrac{1}{2}$ & $\dfrac{1}{4}$ & $\dfrac{1}{8}$ & $\dfrac{1}{16}$ & $\dfrac{1}{32}$ & $\dfrac{1}{32}$ \\
        \cline{1-7}
        \hline
        \end{tabular}
        
        \end{center}

        求随机变量 $Y = \sin \dfrac{\pi }{2}X$ 的分布列.
        \begin{jieda}
            由 $Y = \sin \dfrac{\pi }{2}X$ ,得
            
            当 $X = 1,5$ 时, $Y = 1$ ;
            
            当 $X = 2,4,6$ 时, $Y = 0$ ;
            
            当 $X = 3$ 时, $Y =  - 1$ .
            
            则 $P\left( {Y = 1}\right)  = P\left( {X = 1}\right)  + P\left( {X = 5}\right)  = \dfrac{1}{2} + \dfrac{1}{32} = \dfrac{17}{32}$ .
            
            $P\left( {Y = 0}\right)  = P\left( {X = 2}\right)  + P\left( {X = 4}\right)  + P\left( {X = 6}\right)  = \dfrac{1}{4} + \dfrac{1}{16} + \dfrac{1}{32} = \dfrac{11}{32}$
            
            $P\left( {Y =  - 1}\right)  = P\left( {X = 3}\right) = \dfrac{1}{8}$
            
            所以随机变量 $Y$ 的分布列为
            
            \begin{center}
            
            \begin{tabular}{|c|c|c|c|}
            \hline
            $Y$ & 1 & 0 & -1 \\
            \cline{1-4}
            $P$ & $\dfrac{17}{32}$ & $\dfrac{17}{32}$ & $\dfrac{1}{8}$ \\
            \cline{1-4}
            \hline
            \end{tabular}
            
            \end{center}
        \end{jieda}
    \end{enumerate}

    \item 两点分布
    \begin{enumerate}
        \item 若随机变量 $X$ 服从两点分布， $P\left( {X = 1}\right)  = {0.33}$ ，则 $P\left( {X = 0}\right)  =$ \blkda{$0.67$}
        \begin{jieda}
            $P\left( {X = 0}\right)  = 1 - P\left( {X = 1}\right)  = {0.67}$ .
        \end{jieda}
        \item 若 $X$ 服从两点分布, $P\left( {X = 1}\right)  - P\left( {X = 0}\right)  = {0.34}$ ,则 $P(X = 1) = $\blkda{$0.67$}
        \begin{jieda}
            依题意可得 $P\left( {X = 1}\right)  + P\left( {X = 0}\right)  = 1$， $P\left( {X = 1}\right)  - P\left( {X = 0}\right)  = {0.34}$
            
            所以 $P\left( {X = 1}\right)  = \dfrac{1 + {0.34}}{2} = {0.67}$ .
        \end{jieda}
        
        \item 已知随机变量 $X$ 服从两点分布,且 $P\left( {X = 0}\right)  = 2{a}^{2},P\left( {X = 1}\right) = a$ ,那么 $a = $\blkda{$\dfrac{1}{2}$}.
        \begin{jieda}
            由题意可知 $P\left( {X = 0}\right)  + P\left( {X = 1}\right)  = a + 2{a}^{2} = 1 \Longrightarrow  a = \dfrac{1}{2}$ 或 $a =  - 1$ ,
            
            由于 $a > 0$ ，所以 $a = \dfrac{1}{2}$
        \end{jieda}
        
        \item 已知随机变量 $X$ 服从两点分布,且 $P\left( {X = 1}\right)  = \dfrac{3}{2}P\left( {X = 0}\right)$ , 则 $P\left( {X = 1}\right)  = $\blkda{$ \dfrac{3}{5}$} .
        \begin{jieda}
            由随机变量 $X$ 服从两点分布,得 $P\left( {X = 1}\right)  + P\left( {X = 0}\right)  = 1$，又因为 $P\left( {X = 1}\right)  = \dfrac{3}{2}P\left( {X = 0}\right)$，所以 $P\left( {X = 1}\right)  = \dfrac{3}{5}$
        \end{jieda}
    \end{enumerate}

    \item 分布列与其他知识综合
    \begin{enumerate}
        \item 甲、乙两人各进行 3 次射击，甲每次击中目标的概率时 $\dfrac{3}{4}$ ，乙每次击中目标的概率 $\dfrac{2}{3}$ ，假设两人射击是否击中目标. 相互之间没有影响；每次射击是否击中目标，相互之间没有影响.
        \begin{enumerate}
            \item 求甲至少有 1 次未击中目标的概率
            \item 记甲击中目标的次数为 $\xi$ ，求 $\xi$ 的概率分布列
            \item 求甲恰好比乙多击中目标 2 次的概率
        \end{enumerate}
        \begin{jieda}
            \begin{enumerate}
                \item 记“甲连续射击 3 次，至少 1 次未击中目标”为事件 ${A}_{1}$ ，由题意知两人射击是否击中目标，相互之间没有影响，射击 3 次，相当于 3 次独立重复试验，故 $P\left( {A}_{1}\right)  = 1 - P\left( \overline{{A}_{1}}\right)  = 1 - {\left( \dfrac{3}{4}\right) }^{3} = \dfrac{37}{64}$
                
                \item 依题可知 $\xi$ 的可能取值为0,1,2,3，并且 $\xi  \sim  B\left( {3,\dfrac{3}{4}}\right) ,P\left( {\xi  = k}\right)  = {{C}}_{3}^{k}{\left( \dfrac{3}{4}\right) }^{k}{\left( \dfrac{1}{4}\right) }^{3 - k}\left( {k = 0,1,2,3}\right)$，即 $P\left( {\xi  = 0}\right)  = {\left( \dfrac{1}{4}\right) }^{3} = \dfrac{1}{64},P\left( {\xi  = 1}\right)  = {{C}}_{3}^{1}\left( \dfrac{3}{4}\right) {\left( \dfrac{1}{4}\right) }^{2} = \dfrac{9}{64}$，$P\left( {\xi  = 2}\right)  = {{C}}_{3}^{2}{\left( \dfrac{3}{4}\right) }^{2}\dfrac{1}{4} = \dfrac{27}{64},P\left( {\xi  = 3}\right)  = {\left( \dfrac{3}{4}\right) }^{3} = \dfrac{27}{64}$
                
                $\xi$ 的概率分布列为:
                
                \begin{center}
                
                \begin{tabular}{|c|c|c|c|c|}
                \hline
                $\xi$ & 0 & 1 & 2 & 3 \\
                \cline{1-5}
                $P$ & $\dfrac{1}{6}$ & $\dfrac{9}{64}$ & $\dfrac{27}{64}$ & $\dfrac{27}{64}$ \\
                \cline{1-5}
                \hline
                \end{tabular}
                
                \end{center}

                \item 设甲恰好比乙多击中目标 2 次为事件 $A$ ，甲恰击中目标 2 次且乙恰击中目标 0 次为事件 ${B}_{1}$ ，甲恰击中目标 3 次且乙恰击中目标 1 次为事件 ${B}_{2}$，则 $A = {B}_{1} + {B}_{2},{B}_{1}\text{ 、 }{B}_{2}$ 为互斥事件, $P\left( A\right)  = P\left( {B}_{1}\right)  + P\left( {B}_{2}\right)  = \dfrac{27}{64} \times  \dfrac{1}{27} + \dfrac{27}{64} \times  \dfrac{6}{27} = \dfrac{7}{64}$
                
                $\therefore$ 甲恰好比乙多击中目标 2 次的概率为 $\dfrac{7}{64}$
            \end{enumerate}
        \end{jieda}
        
        \item 某校高三年级拟派出甲、乙、丙三人去参加校运动会 100m 跑项目. 比赛分为初赛和决赛,其中初赛有两轮,只有两轮都获胜才能进入决赛. 已知甲在每轮比赛中获胜的概率均为 $\dfrac{3}{4}$ ; 乙在第一轮和第二轮比赛中获胜的概率分别为 $\dfrac{2}{3}$ 和 $\dfrac{3}{4}$ ; 丙在第一轮和第二轮获胜的概率分别为 $p$ 和 $\dfrac{3}{2} - p$ ,其中 $\dfrac{1}{2} < p < \dfrac{3}{4}$
        \begin{enumerate}
            \item 甲、乙、丙三人中，谁进入决赛的可能性最大
            \item 若甲、乙、丙三人中恰有两人进入决赛的概率为 $\dfrac{29}{72}$ ，求 $p$ 的值
            \item 在(2)的条件下，设进入决赛的人数为 $\xi$ ，求 $\xi$ 的分布列
        \end{enumerate}
        \begin{jieda}
            \begin{enumerate}
                \item 甲进入决赛的概率为 ${p}_{1} = {\left( \dfrac{3}{4}\right) }^{2} = \dfrac{9}{16}$ ，乙进入决赛的概率为 ${p}_{2} = \dfrac{2}{3} \times  \dfrac{3}{4} = \dfrac{1}{2}$ ，丙进入决赛的概率为 ${p}_{3} = p\left( {\dfrac{3}{2} - p}\right)  =  - {\left( p - \dfrac{3}{4}\right) }^{2} + \dfrac{9}{16}$ ,而 $\dfrac{1}{2} < p < \dfrac{3}{4}$ ,则 $\dfrac{1}{2} < {p}_{3} < \dfrac{3}{16}$，所以甲进入决赛的可能性最大.
                
                \item 甲、乙、丙三人中恰有两人进入决赛的概率为$\dfrac{9}{16} \cdot  \dfrac{1}{2} \cdot  \left\lbrack  {1 - p\left( {\dfrac{3}{2} - p}\right) }\right\rbrack   + \dfrac{9}{16} \cdot  \left( {1 - \dfrac{1}{2}}\right)  \cdot  p\left( {\dfrac{3}{2} - p}\right)  + \left( {1 - \dfrac{9}{16}}\right)  \cdot  \dfrac{1}{2} \cdot  p\left( {\dfrac{3}{2} - p}\right)  = \dfrac{29}{72}$
                
                整理可得 ${18}{p}^{2} - {27p} + {10} = 0$ ，而 $\dfrac{1}{2} < p < \dfrac{3}{4}$ ，所以 $p = \dfrac{2}{3}$
                
                \item 依题意，甲、乙、丙进入决赛的概率分别为 $\dfrac{9}{16},\dfrac{1}{2},\dfrac{5}{9}$ .
                
                随机变量 $\xi$ 的可能取值有 0,1,2,3,
                
                $P\left( {\xi  = 0}\right)  = \dfrac{7}{16} \times  \dfrac{1}{2} \times  \dfrac{4}{9} = \dfrac{7}{72},P\left( {\xi  = 2}\right)  = \dfrac{29}{72},P\left( {\xi  = 3}\right)  = \dfrac{9}{16} \times  \dfrac{1}{2} \times  \dfrac{5}{9} = \dfrac{5}{32}$ ,
                
                $P\left( {\xi  = 1}\right)  = \dfrac{9}{16} \cdot  \left( {1 - \dfrac{1}{2}}\right)  \cdot  \left( {1 - \dfrac{5}{9}}\right)  + \left( {1 - \dfrac{9}{16}}\right)  \cdot  \dfrac{1}{2} \cdot  \left( {1 - \dfrac{5}{9}}\right)  + \left( {1 - \dfrac{9}{16}}\right)  \cdot  \left( {1 - \dfrac{1}{2}}\right)  \cdot  \dfrac{5}{9} = \dfrac{11}{32}$
                
                所以随机变量 $\xi$ 的分布列为:
                
                \begin{center}
                
                \begin{tabular}{|c|c|c|c|c|}
                \hline
                $\xi$ & 0 & 1 & 2 & 3 \\
                \cline{1-5}
                $P$ & $\dfrac{7}{72}$ & $\dfrac{11}{32}$ & $\dfrac{29}{72}$ & $\dfrac{5}{32}$ \\
                \cline{1-5}
                \hline
                \end{tabular}
                
                \end{center}
            \end{enumerate}
        \end{jieda}
        
        \item 综合素质评价是高考招生制度改革的内容之一. 某高中采用多维评分的方式进行综合素质评价. 下图是该校高三学生“运动与健康”评价结果的频率直方图，评分在区间 $[90,100]$, $\lbrack {70},{90})$ , $\lbrack {60},{70})$ , $\lbrack {50},{60})$ 上，分别对应为 $A,B,C,D$ 四个等级. 为了进一步引导学生对运动与健康的重视，初评获 $A$ 等级的学生不参加复评，等级不变，对其余学生学校将进行一次复评. 复评中， 原获 $B$ 等级的学生有 $\dfrac{1}{4}$ 的概率提升为 $A$ 等级；原获 $C$ 等级的学生有 $\dfrac{1}{5}$ 的概率提升为 $B$ 等级；原获 $D$ 等级的学生有 $\dfrac{1}{6}$ 的概率提升为 $C$ 等级. 用频率估计概率,每名学生复评结果相互独立.
        
        \begin{center}
        \includegraphics[width=0.3\textwidth]{bodyxb/src/bo_d67fmqn7aajc739ac5a0_13_138_1340_450_360_0.jpg}
        \end{center}

        \begin{enumerate}
            \item 若初评中甲获得 $C$ 等级，乙、丙获得 $D$ 等级，记甲、乙、丙三人复评后等级为 $C$ 等级的人数为 $\xi$ ，求 $\xi$ 的分布列
            \item 从全体高三学生中任选 1 人，在已知该学生是复评晋级的条件下，求他初评是 $B$ 等级的概率
        \end{enumerate}
        \begin{jieda}
            \begin{enumerate}
                \item 依题意 $\xi$ 的所有可能取值为0,1,2,3，
                
                所以 $P\left( {\xi  = 0}\right)  = \dfrac{1}{5} \times  \dfrac{5}{6} \times  \dfrac{5}{6} = \dfrac{5}{36},P\left( {\xi  = 1}\right)  = \dfrac{4}{5} \times  \dfrac{5}{6} \times  \dfrac{5}{6} + \dfrac{1}{5} \times  {{C}}_{2}^{1} \times  \dfrac{1}{6} \times  \dfrac{5}{6} = \dfrac{11}{18}$ ，
                
                $P\left( {\xi  = 2}\right)  = \dfrac{4}{5} \times  {{C}}_{2}^{1} \times  \dfrac{1}{6} \times  \dfrac{5}{6} + \dfrac{1}{5} \times  \dfrac{1}{6} \times  \dfrac{1}{6} = \dfrac{41}{180},P\left( {\xi  = 3}\right)  = \dfrac{4}{5} \times  \dfrac{1}{6} \times  \dfrac{1}{6} = \dfrac{1}{45},$
                
                $\therefore \xi$ 的分布列如下:
                
                \begin{center}
                
                \begin{tabular}{|c|c|c|c|c|}
                \hline
                $\xi$ & 0 & 1 & 2 & 3 \\
                \cline{1-5}
                $P$ & $\dfrac{5}{36}$ & $\dfrac{11}{18}$ & $\dfrac{41}{180}$ & $\dfrac{1}{45}$ \\
                \cline{1-5}
                \hline
                \end{tabular}
                
                \end{center}
                
                \item 记事件 $M$ 为“该学生复评晋级”，事件 $B$ 为“该学生初评是 $B$ ”.
                
                所以 $P\left( {B \mid  M}\right)  = \dfrac{P\left( {BM}\right) }{P\left( M\right) } = \dfrac{{0.6} \times  \dfrac{1}{4}}{{0.6} \times  \dfrac{1}{4} + {0.15} \times  \dfrac{1}{5} + {0.05} \times  \dfrac{1}{6}} = \dfrac{90}{113}$ .
            \end{enumerate}
        \end{jieda}
        \item 从甲、乙、丙、丁 4 人中随机抽取 3 个人去做传球训练.训练规则是确定一人第一次将球传出，每次传球时，传球者都等可能地将球传给另外两个人中的任何一人，每次必须将球传出.
        \begin{enumerate}
            \item 记甲乙丙三人中被抽到的人数为随机变量 $X$ ，求 $X$ 的分布列
            \item 若刚好抽到甲乙丙三个人相互做传球训练，且第 1 次由甲将球传出，记 $n$ 次传球后球在甲手中的概率为 ${p}_{n}$ ， $n = 1,2,3,\cdots$ 
            \begin{dingautolist}{172}
                \item 直接写出 ${p}_{1},{p}_{2},{p}_{3}$ 的值
                \item 求 ${p}_{n + 1}$ 与 ${p}_{n}$ 的关系式 $\left( {n \in  {{N}}^{ * }}\right)$ ，并求 ${p}_{n}\left( {n \in  {{N}}^{ * }}\right)$
            \end{dingautolist}
        \end{enumerate}
        \begin{jieda}
            \begin{enumerate}
                \item $X$ 的可能取值为 2 和 3，则 $P\left( {X = 2}\right)  = \dfrac{C_{3}^{2}}{C_{4}^{3}} = \dfrac{3}{4},P\left( {X = 3}\right)  = \dfrac{C_{3}^{3}}{C_{4}^{3}} = \dfrac{1}{4}$
                
                所以随机变量 $X$ 的分布列为:
                
                \begin{center}
                
                \begin{tabular}{|c|c|c|}
                \hline
                $X$ & 2 & 3 \\
                $P$ & $\dfrac{3}{4}$ & $\dfrac{1}{4}$ \\
                \cline{1-3}
                \hline
                \end{tabular}
                
                \end{center}
                \item \begin{dingautolist}{172}
                    \item 若刚好抽到甲乙丙三个人相互做传球训练，且第 1 次由甲将球传出， $n$ 次传球后球在甲手中的概率为 ${p}_{n},n = 1,2,3,\cdots$，则有 ${p}_{1} = 0,{p}_{2} = 2 \times  \dfrac{1}{2} \times  \dfrac{1}{2} = \dfrac{2}{{2}^{2}} = \dfrac{1}{2},{p}_{3} = 2 \times  \dfrac{1}{2} \times  \dfrac{1}{2} \times  \dfrac{1}{2} = \dfrac{2}{{2}^{3}} = \dfrac{1}{4}$

                    \item 记 ${A}_{n}$ 表示事件“经过 $n$ 次传球后，球在甲手中”，
                    
                    ${A}_{n + 1} = \overline{{A}_{n}} \cdot  {A}_{n + 1} + {A}_{n} \cdot  {A}_{n + 1}$
                    
                    所以 ${p}_{n + 1} = P\left( {\overline{{A}_{n}} \cdot  {A}_{n + 1} + {A}_{n} \cdot  {A}_{n + 1}}\right)  = P\left( {\overline{{A}_{n}} \cdot  {A}_{n + 1}}\right)  + P\left( {{A}_{n} \cdot  {A}_{n + 1}}\right)$
                    
                    $= P\left( \overline{{A}_{n}}\right)  \cdot  P\left( {{A}_{n + 1} \mid  \overline{{A}_{n}}}\right)  + P\left( {A}_{n}\right)  \cdot  P\left( {{A}_{n + 1} \mid  {A}_{n}}\right)  = \left( {1 - {p}_{n}}\right)  \cdot  \dfrac{1}{2} + {p}_{n} \cdot  0 = \dfrac{1}{2}\left( {1 - {p}_{n}}\right)$
                    
                    即 ${p}_{n + 1} =  - \dfrac{1}{2}{p}_{n} + \dfrac{1}{2},n = 1,2,3$ ,
                    
                    所以 ${p}_{n + 1} - \dfrac{1}{3} =  - \dfrac{1}{2}\left( {{p}_{n} - \dfrac{1}{3}}\right)$ ,且 ${p}_{1} - \dfrac{1}{3} =  - \dfrac{1}{3} \neq  0$ .
                    
                    所以数列 $\left\{  {{p}_{n} - \dfrac{1}{3}}\right\}$ 表示以 $- \dfrac{1}{3}$ 为首项， $= \dfrac{1}{2}$ 为公比的等比数列，
                    
                    所以 ${p}_{n} - \dfrac{1}{3} =  - \dfrac{1}{3} \times  {\left( -\dfrac{1}{2}\right) }^{n - 1}$ ，所以 ${p}_{n} =  - \dfrac{1}{3} \times  {\left( -\dfrac{1}{2}\right) }^{n - 1} + \dfrac{1}{3} = \dfrac{1}{3}\left\lbrack  {1 - {\left( -\dfrac{1}{2}\right) }^{n - 1}}\right\rbrack$
                    
                    即 $n$ 次传球后球在甲手中的概率是 $\dfrac{1}{3}\left\lbrack  {1 + \dfrac{{\left( -1\right) }^{n}}{{2}^{n - 1}}}\right\rbrack$ .
                \end{dingautolist}
            \end{enumerate}
        \end{jieda}
    \end{enumerate}
    \item 离散型随机变量的均值
    \begin{enumerate}
        \item 若随机变量 $X$ 的分布列如表,则 $E\left( X\right)  =$ \blkda{$\dfrac{37}{12}$}

        \begin{center}
        
        \begin{tabular}{|c|c|c|c|c|}
        \hline
        $X$ & 1 & 2 & 3 & 4 \\
        \cline{1-5}
        $P$ & $\dfrac{1}{4}$ & $\dfrac{1}{4}$ & $a$ & $\dfrac{1}{3}$ \\
        \cline{1-5}
        \hline
        \end{tabular}
        
        \end{center}
        \begin{jieda}
            根据题意可得 $a = 1 - \dfrac{1}{4} - \dfrac{1}{4} - \dfrac{1}{3} = \dfrac{1}{6}$，则 $E\left( X\right)  = 1 \times  \dfrac{1}{4} + 2 \times  \dfrac{1}{4} + 3 \times  \dfrac{1}{6} + 4 \times  \dfrac{1}{3} = \dfrac{31}{12}$
        \end{jieda}

        
        \item 近年来中国人工智能产业爆发式的增长，推动了 AI 电商行业的快速发展, 已知 2020-2023 年中国 AI 解决方案提供商企业数量分别为 1617, 2106, 2329, 2896, 从这 4 个数字中任取 2 个数字,当所取两个数字差的绝对值小于 500 时,随机变量 $X = \dfrac{1}{2}$ ; 当所取两个数字差的绝对值不小于 500 时,随机变量 $X = 1$ ,则 $E\left( X\right)  =$\blkda{$\dfrac{5}{6}$}
        \begin{jieda}
            从这 4 个数字中任取 2 个数字, 结果有 6 种，所取两个数字差的绝对值小于 500 的结果有 2 种,故 $P\left( {X = \dfrac{1}{2}}\right)  = \dfrac{1}{3}$，不小于 500 的结果有 4 种,故 $P\left( {X = 1}\right)  = \dfrac{2}{3}$
            
            所以 $E\left( X\right)  = \dfrac{1}{2} \times  \dfrac{1}{3} + 1 \times  \dfrac{2}{3} = \dfrac{5}{6}$
        \end{jieda}

        
        
        \item 若随机变量 $X$ 服从两点分布,其中 $P\left( {X = 0}\right)  = \dfrac{1}{3}$ ,则以下正确的是\dotfill{(\kh{D})}
        \xx{$E\left( X\right)  = \dfrac{1}{3}$}{$E\left( {{2X} + 3}\right)  = \dfrac{4}{3}$}{$E\left( {{2X} + 2}\right)  = \dfrac{8}{3}$}{$E\left( {{2X} + 1}\right)  = \dfrac{7}{3}$}
        \begin{jieda}
            因为随机变量 $X$ 服从两点分布,且 $P\left( {X = 0}\right)  = \dfrac{1}{3}$ 则 $P\left( {X = 1}\right)  = \dfrac{2}{3}$
            
            故 $E\left( X\right)  = 0 \times  \dfrac{1}{3} + 1 \times  \dfrac{2}{3} = \dfrac{2}{3}$ 故 A 错误;
            
            $E\left( {{2X} + 3}\right)  = {2E}\left( X\right)  + 3 = \dfrac{13}{3}$ ，故 B 错误；
            
            $E\left( {{2X} + 2}\right)  = {2E}\left( X\right)  + 2 = \dfrac{10}{3}$ ，故 C 错误；
            
            $E\left( {{2X} + 1}\right)  = {2E}\left( X\right)  + 1 = \dfrac{7}{3}$ 故 D 正确.
        \end{jieda}

        \item 已知 $\xi$ 的分布列为

        \begin{center}
        
        \begin{tabular}{|c|c|c|c|c|}
        \hline
        $\xi$ & 1 & 2 & 3 & 4 \\
        \cline{1-5}
        $P$ & $\dfrac{1}{6}$ & $\dfrac{1}{6}$ & $\dfrac{1}{3}$ & $m$ \\
        \cline{1-5}
        \hline
        \end{tabular}
        
        \end{center}
        
        设 $\eta  = {2\xi } - 5$ ,则 $E\left( \eta \right)  =$\blkda{$\dfrac{2}{3}$}
        \begin{jieda}
            由题意得 $\dfrac{1}{6} + \dfrac{1}{6} + \dfrac{1}{3} + m = 1$ ,解得 $m = \dfrac{1}{3}$，故 $E\left( \xi \right)  = 1 \times  \dfrac{1}{6} + 2 \times  \dfrac{1}{6} + 3 \times  \dfrac{1}{3} + 4 \times  \dfrac{1}{3} = \dfrac{17}{6}$ ,而 $\eta  = {2\xi } - 5$，则 $E\left( \eta \right)  = {2E}\left( \xi \right)  - 5 = 2 \times  \dfrac{17}{6} - 5 = \dfrac{2}{3}$
        \end{jieda}

        \item 已知随机变量 $X$ 的概率分布列为$P\left( {X = n}\right)  = a\sin \dfrac{\pi }{{4n} - 2}\left( {n = 1,2}\right)$ ，其中 $a$ 是常数，则 $E\left( {\dfrac{1}{a}X}\right)  =$\blkda{2}
        \begin{jieda}
            由 $P\left( {X = n}\right)  = a\sin \dfrac{\pi }{{4n} - 2}\left( {n = 1,2}\right)$ ,得 $P\left( {X = 1}\right)  = a,P\left( {X = 2}\right)  = \dfrac{1}{2}a$
            
            由 $P\left( {X = 1}\right)  + P\left( {X = 2}\right)  = 1$ ,得 $a = \dfrac{2}{3}$ ,
            
            于是 $E\left( X\right)  = 1 \times  P\left( {X = 1}\right)  + 2 \times  P\left( {X = 2}\right)  = {2a} = \dfrac{4}{3}$ .
            
            所以 $E\left( {\dfrac{1}{a}X}\right)  = \dfrac{3}{2}E\left( X\right)  = 2$ .
        \end{jieda}

        \item 从一批含有 8 件正品，2 件次品的产品中不放回地抽 3 次，每次抽取 1 件，设抽取的次品数为 $\xi$ ，则 $E\left( {{5\xi } + 1}\right)  =$\blkda{4}
        \begin{jieda}
            由题意,随机变量 $\xi$ 的取值分别为0,1,2,
            
            可得 $P\left( {\xi  = 0}\right)  =
            \dfrac{C_8^3}{C_{10}^3}= \dfrac{7}{15}$；$P\left( {\xi  = 1}\right)  = \dfrac{C_8^2C_2^1}{C_{10}^3} = \dfrac{7}{15}$；$P\left( {\xi  = 2}\right)  = \dfrac{C_8^1C_2^2}{C_{10}^3} = \dfrac{1}{15}$
            
            所以 $E\left( \xi \right)  = 0 \times  \dfrac{7}{15} + 1 \times  \dfrac{7}{15} + 2 \times  \dfrac{1}{15} = \dfrac{3}{5}$ 可得 $E\left( {{5\xi } + 1}\right)  = {5E}\left( \xi \right)  + 1 = 5 \times  \dfrac{3}{5} + 1 = 4$ .
        \end{jieda}
        \item 已知随机变量 $X$ 的分布列:
        
        \begin{center}
        
        \begin{tabular}{|c|c|c|c|}
        \hline
        $x$ & -1 & 0 & 1 \\
        \cline{1-4}
        $P$ & 1 2 & 1 3 & 1 \\
        \cline{1-4}
        \hline
        \end{tabular}
        
        \end{center}
        
        满足 $Y = {aX} + 3,E\left( Y\right)  = \dfrac{5}{3}$ ，则 $a$ 的值为\blkda{4}
        \begin{jieda}
            由表可得 $E\left( X\right)  =  - \dfrac{1}{2} + 0 + \dfrac{1}{6} =  - \dfrac{1}{3}$，又 $Y = {aX} + 3,E\left( Y\right)  = \dfrac{5}{3}$ 可得 $E\left( Y\right)  = {aE}\left( X\right)  + 3 =  - \dfrac{1}{3}a + 3 = \dfrac{5}{3}$ 解得 $a = 4$
        \end{jieda}

        \item 随机变量 $Y$ 的分布列为下表所示,若 $Y$ 的期望值为 1，则\dotfill{(\kh{A})}
        
        \begin{center}
        
        \begin{tabular}{|c|c|c|c|}
        \hline
        $Y$ & -2 & 0 & 2 \\
        \cline{1-4}
        $P\left( Y\right)$ & $a$ & $b$ & $c$ \\
        \cline{1-4}
        \hline
        \end{tabular}
        
        \end{center}
        \xx{$a + b + c = 1$}{${2a} + {2c} = 4$}{$- {2a} + {2c} = 2$}{$a = c$}
        \begin{jieda}
            由分布列的性质可知, $a + b + c = 1$ ,故 ${A}$ 正确;
            
            因为 $Y$ 的期望值为 1,所以 $- {2a} + {2c} = 1$ ,所以 $C$ 错.
            
            若 ${2a} + {2c} = 4 \Rightarrow  a + c = 2 > 1$ ,不满足分布列性质, ${B}$ 错,
            
            由上,有 $c - a = \dfrac{1}{2}$ 显然 ${D}$ 错.
        \end{jieda}
        
        \item 已知随机变量 $X$ 的分布列如下所示，且 $E\left( X\right)  = \dfrac{13}{6}$ ，则 $m =$\blkda{$\dfrac{1}{6}$}
        
        \begin{center}
        
        \begin{tabular}{|c|c|c|c|}
        \hline
        $X$ & 1 & 2 & 3 \\
        \cline{1-4}
        $P$ & $n$ & $m$ & $\dfrac{1}{2}$ \\
        \cline{1-4}
        \hline
        \end{tabular}
        
        \end{center}

        \begin{jieda}
            由分布列的性质可得, $n + m + \dfrac{1}{2} = 1$ ,所以 $n + m = \dfrac{1}{2}$，又因为 $E\left( X\right)  = \dfrac{13}{6}$，所以 $E\left( X\right)  = n + {2m} + 3 \times  \dfrac{1}{2} = \dfrac{13}{6}$ ,即 $n + {2m} = \dfrac{2}{3}$ ,
            
            联立方程 $\left\{  \begin{array}{l} n + m = \dfrac{1}{2} \\  n + {2m} = \dfrac{2}{3} \end{array}\right.$ ，解得 $\left\{  \begin{array}{l} m = \dfrac{1}{6} \\  n = \dfrac{1}{3} \end{array}\right.$ ，所以 $m = \dfrac{1}{6}$
        \end{jieda}
        
        \item 下表是离散型随机变量 $X$ 的分布列,且满足 $E\left( X\right)  = \dfrac{28}{5}$ , 则 $p,q$ 的值分别是\blkda{$\dfrac{4}{15},\dfrac{1}{5}$}
        
        \begin{center}
        
        \begin{tabular}{|c|c|c|c|c|}
        \hline
        $X$ & 3 & 4 & 5 & 9 \\
        \cline{1-5}
        $P$ & $p$ & $q$ & $\dfrac{1}{5}$ & $\dfrac{1}{3}$ \\
        \cline{1-5}
        \hline
        \end{tabular}
        
        \end{center}
        \begin{jieda}
            由题意,得 $p + q + \dfrac{1}{5} + \dfrac{1}{3} = 1$ ,所以 $p + q = \dfrac{7}{15}$.
            
            因为 $E\left( X\right)  = 3 \cdot  p + 4 \cdot  q + 5 \times  \dfrac{1}{5} + 9 \times  \dfrac{1}{3} = \dfrac{28}{5}$ ,所以 ${3p} + {4q} = \dfrac{8}{5}$.
            
            解得: $q = \dfrac{1}{5},p = \dfrac{4}{15}$ .
        \end{jieda} 
        \item 甲、乙、丙三位同学报名参加学校的社团活动，每个同学彼此独立地从足球、篮球、围棋、合唱四个社团中随机选报两个社团.
        \begin{enumerate}
            \item 求恰有两个同学选报的社团完全相同的概率.
            \item 求同学甲选报足球社的概率.
            \item 若甲已经报名参加了合唱社团，只需在其余三个社团中再选报一个，乙、丙从四个社团中随机选两个，设报名足球社的同学人数为 $X$ ,求随机变量 $X$ 的分布列及数学期望.
        \end{enumerate}
        \begin{jieda}
            \begin{enumerate}
                \item 先从 3 个同学中选出 2 个同学，有 $C_{3}^{2} = 3$ ，从 4 个社团中选 2 个，有 $C_{4}^{2} = 6$ 种方法，因此每位同学选报社团都有 6 种方法，因此恰好两个同学选报的社团一样的概率为 $P = 3 \times  \dfrac{1}{6} \times  \dfrac{5}{6} = \dfrac{5}{12}$

                \item 同学甲选报足球社的概率为 $\dfrac{C_{3}^1}{C_{4}^{2}} = \dfrac{1}{2}$ .

                \item 甲报足球的概率为 $\dfrac{1}{3}$ 不报的概率为 $\dfrac{2}{3}$ .
                
                乙丙报足球的概率均为 $\dfrac{1}{2}$ ，不报的概率为 $\dfrac{1}{2}$ ，
                
                故 $X$ 可取0,1,2,3,
                
                $P\left( {X = 3}\right)  = \dfrac{1}{3} \times  \dfrac{1}{2} \times  \dfrac{1}{2} = \dfrac{1}{12},P\left( {X = 2}\right)  = \dfrac{1}{3} \times  \dfrac{1}{2} \times  \left( {1 - \dfrac{1}{2}}\right)  \times  2 + \dfrac{2}{3} \times  \dfrac{1}{2} \times  \dfrac{1}{2} = \dfrac{1}{3}$ ,
                
                $P\left( {X = 1}\right)  = \dfrac{1}{3} \times  \left( {1 - \dfrac{1}{2}}\right)  \times  \left( {1 - \dfrac{1}{2}}\right)  + \dfrac{2}{3} \times  \dfrac{1}{2} \times  \left( {1 - \dfrac{1}{2}}\right)  \times  2 = \dfrac{5}{12},P\left( {X = 0}\right)  = \dfrac{2}{3} \times  \left( {1 - \dfrac{1}{2}}\right)  \times  \left( {1 - \dfrac{1}{2}}\right)  = \dfrac{1}{6},$ 故 $X$ 的分布列为:
                
                \begin{center}
                
                \begin{tabular}{|c|c|c|c|c|}
                \hline
                $X$ & 0 & 1 & 2 & 3 \\
                \cline{1-5}
                $P$ & $\dfrac{2}{12}$ & $\dfrac{5}{12}$ & $\dfrac{4}{12}$ & $\dfrac{1}{12}$ \\
                \cline{1-5}
                \hline
                \end{tabular}
                
                \end{center}
                
                故 $E\left( X\right)  = 0 + \dfrac{5}{12} + \dfrac{8}{12} + \dfrac{3}{12} = \dfrac{16}{12} = \dfrac{4}{3}$ .
            \end{enumerate}
        \end{jieda}
        \item 已知某单位甲、乙、丙三个部门的员工人数分别为 18,36， 9. 现采用分层抽样的方法从中抽取 7 人，进行睡眠时间的调查.
        \begin{enumerate}
            \item 应从甲、乙、丙三个部门的员工中分别抽取多少人？
            \item 若抽出的 7 人中有 3 人睡眠不足，4 人睡眠充足，现从这 7 人中随机抽取 3 人做进一步的身体检查. 用 $X$ 表示抽取的 3 人中睡眠充足的员工人数，求随机变量 $X$ 的分布列与数学期望.
        \end{enumerate}
        \begin{jieda}
            \begin{enumerate}
                \item 某单位甲乙丙三个部门的员工人数分别为 18,36,9，现采用分层抽样的方法, 从中抽取 7 人，进行睡眠时间的调查，则从甲部门的员工中抽取 $7 \times  \dfrac{18}{{18} + {36} + 9} = 2$ 人，从乙部门的员工中抽取 $7 \times  \dfrac{36}{{18} + {36} + 9} = 4$ 人，从丙部门的员工中抽取 $7 \times  \dfrac{9}{{18} + {36} + 9} = 1$ 人.
                \item 若抽取的 7 人中有 3 人睡眠不足，4 人睡眠充足，现从这 7 人中随机抽取 3 人做进一步的身体检查,用 $X$ 表示抽取的 3 人中睡眠充足的员工人数,则 $X$ 的可能取值为 0,1,2,3 ,
                
                $P\left( {X = 0}\right)  = \dfrac{C_{3}^{3}}{C_7^{3}} = \dfrac{1}{35},P\left( {X = 1}\right)  = \dfrac{C_{4}^{1}C_{3}^{2}}{C_7^{3}} = \dfrac{12}{35}$
                
                $P\left( {X = 2}\right)  = \dfrac{C_{4}^{2}C_{3}^{1}}{C_7^{3}} = \dfrac{18}{35},P\left( {X = 3}\right)  = \dfrac{C_{4}^{3}}{C_7^{3}} = \dfrac{4}{35}$
                
                所以随机变量 $X$ 的分布列为:
                
                \begin{center}
                
                \begin{tabular}{|c|c|c|c|c|}
                \hline
                $X$ & 0 & 1 & 2 & 4 \\
                \cline{1-5}
                $P$ & $\dfrac{1}{35}$ & $\dfrac{12}{35}$ & $\dfrac{18}{35}$ & $\dfrac{4}{35}$ \\
                \cline{1-5}
                \hline
                \end{tabular}
                
                \end{center}
                
                则数学期望为 $E\left( X\right)  = 0 \times  \dfrac{1}{35} + 1 \times  \dfrac{12}{35} + 2 \times  \dfrac{18}{35} + 3 \times  \dfrac{4}{35} = \dfrac{12}{7}$ .
            \end{enumerate}
        \end{jieda}
    \end{enumerate}
    \item 离散型随机变量的方差、标准差
    \begin{enumerate}
        
        \item 甲、乙两人进行定点投篮游戏，投篮者若投中，则继续投篮， 否则由对方投篮，第一次由甲投篮. 已知每次投篮甲、乙命中的概率分别为 $\dfrac{1}{3},\dfrac{3}{4}$ ，在前 3 次投篮中，乙投篮的次数为 $X$ ,求 $X$ 的分布列、方差及标准差.
        \begin{jieda}
            由题意得, $X$ 的可能取值为 0,1,2 .
            
            $P\left( {X = 0}\right)  = \dfrac{1}{3} \times  \dfrac{1}{3} = \dfrac{1}{9}$
            
            $P\left( {X = 1}\right)  = \dfrac{1}{3} \times  \dfrac{2}{3} + \dfrac{2}{3} \times  \dfrac{1}{4} = \dfrac{7}{18}$
            
            $P\left( {X = 2}\right)  = \dfrac{2}{3} \times  \dfrac{3}{4} = \dfrac{1}{2}$
            
            故 $X$ 的分布列为
            
            \begin{center}
            
            \begin{tabular}{|c|c|c|c|}
            \hline
            $X$ & 0 & 1 & 2 \\
            \cline{1-4}
            $P$ & $\dfrac{1}{9}$ & $\dfrac{7}{18}$ & $\dfrac{1}{2}$ \\
            \cline{1-4}
            \hline
            \end{tabular}
            
            \end{center}
            
            $E\left( X\right)  = 0 \times  \dfrac{1}{9} + 1 \times  \dfrac{7}{18} + 2 \times  \dfrac{1}{2} = \dfrac{25}{18}$
            
            $D\left( X\right)  = {\left( 0 - \dfrac{25}{18}\right) }^{2} \times  \dfrac{1}{9} + {\left( 1 - \dfrac{25}{18}\right) }^{2} \times  \dfrac{7}{18} + {\left( 2 - \dfrac{25}{18}\right) }^{2} \times  \dfrac{1}{2} = \dfrac{149}{324}$
            
            $\therefore \sqrt{D\left( X\right) } = \dfrac{\sqrt{149}}{18}$ .
        \end{jieda}
        \item 不透明口袋中有 $n$ 个相同的黑色小球和红色、白色、蓝色的小球各 1 个,从中任取 4 个小球, $\xi$ 表示当 $n = 2$ 时取出黑球的数目, $\eta$ 表示当 $n = 3$ 时取出黑球的数目,则下列结论中成立的是\dotfill{(\kh{A})}
        \xx{$E\left( \xi \right)  < E\left( \eta \right) ,D\left( \xi \right)  < D\left( \eta \right)$}{$E\left( \xi \right)  > E\left( \eta \right) ,D\left( \xi \right)  < D\left( \eta \right)$}{$E\left( \xi \right)  < E\left( \eta \right) ,D\left( \xi \right)  > D\left( \eta \right)$}{$E\left( \xi \right)  > E\left( \eta \right) ,D\left( \xi \right)  > D\left( \eta \right)$}
        \begin{jieda}
            当 $n = 2$ 时, $\xi$ 的可能取值为 1,2,
            
            $P\left( {\xi  = 1}\right)  = \dfrac{C_{2}^{1}C_{3}^{3}}{C_{5}^{4}} = \dfrac{2}{5},P\left( {\xi  = 2}\right)  = \dfrac{C_{2}^{2}C_{3}^{2}}{C_{5}^{4}} = \dfrac{3}{5}.$
            
            因此 $E\left( \xi \right)  = 1 \times  \dfrac{2}{5} + 2 \times  \dfrac{3}{5} = \dfrac{8}{5},D\left( \xi \right)  = \dfrac{9}{25} \times  \dfrac{2}{5} + \dfrac{4}{25} \times  \dfrac{3}{5} = \dfrac{6}{25}$ ;
            
            当 $n = 3$ 时, $\eta$ 的可能取值为1,2,3,
            
            $P\left( {\eta  = 1}\right)  = \dfrac{C_{3}^{1}C_{3}^{3}}{C_{6}^{4}} = \dfrac{1}{5},P\left( {\eta  = 2}\right)  = \dfrac{C_{3}^{2}C_{3}^{2}}{C_{6}^{4}} = \dfrac{3}{5}P\left( {\eta  = 3}\right)  = \dfrac{C_{3}^{3}C_{3}^{1}}{C_{6}^{4}} = \dfrac{1}{5},$
            
            因此 $E\left( \eta \right)  = 1 \times  \dfrac{1}{5} + 2 \times  \dfrac{3}{5} + 3 \times  \dfrac{1}{5} = 2,D\left( \eta \right)  = 1 \times  \dfrac{1}{5} + 0 \times  \dfrac{3}{5} + 1 \times  \dfrac{1}{5} = \dfrac{2}{5}$ ,
            
            所以 $E\left( \xi \right)  < E\left( \eta \right) ,D\left( \xi \right)  < D\left( \eta \right)$ .
        \end{jieda}
        \item 设离散型随机变量 $X$ 的分布列为
        
        \begin{center}
        
        \begin{tabular}{|c|c|c|c|c|c|}
        \hline
        $X$ & 0 & 1 & 2 & 3 & 4 \\
        \cline{1-6}
        $P$ & 0.1 & 0.4 & $q$ & 0.2 & 0.2 \\
        \cline{1-6}
        \hline
        \end{tabular}
        
        \end{center}

        若离散型随机变量 $Y$ 满足 $Y = {2X} - 1$ ，则下列结论错误的是\dotfill{(\kh{B})}
        \xx{$q = {0.1}$}{$E\left( X\right)  = 2,D\left( X\right)  = {1.4}$}{$E\left( X\right)  = 2,D\left( X\right)  = {1.8}$}{$E\left( Y\right)  = 3,D\left( Y\right)  = {7.2}$}
        \begin{jieda}
            对于选项 ${A}$ ,因为 ${0.1} + {0.4} + q + {0.2} + {0.2} = 1$ ,解得 $q = {0.1}$ ,所以选项 ${A}$ 正确

            又 $E\left( X\right)  = 0 \times  {0.1} + 1 \times  {0.4} + 2 \times  {0.1} + 3 \times  {0.2} + 4 \times  {0.2} = 2,D\left( X\right)  = {\left( 0 - 2\right) }^{2} \times  {0.1} + {\left( 1 - 2\right) }^{2} \times  {0.4} +  {\left( 2 - 2\right) }^{2} \times  {0.1} + {\left( 3 - 2\right) }^{2} \times  {0.2} + {\left( 4 - 2\right) }^{2} \times  {0.2} = {1.8}$ ,
            
            所以选项 B 错误,选项 C 正确,
            
            对于选项 D,因为 $Y = {2X} - 1$ ,所以 $E\left( Y\right)  = E\left( {{2X} - 1}\right)  = {2E}\left( X\right)  - 1 = 2 \times  2 - 1 = 3,D\left( Y\right)  = D({2X} -  1) = {4D}\left( X\right)  = 4 \times  {1.8} = {7.2}$ ,所以选项 D 正确
        \end{jieda}

        \item 随机变量 $X$ 的分布列如下,且 $E\left( X\right)  = \dfrac{7}{5}$ 则 $D\left( {\dfrac{1}{2}X}\right)  =$\blkda{$0.16$}
        
        \begin{center}
        
        \begin{tabular}{|c|c|c|c|}
        \hline
        $X$ & 0 & 1 & 2 \\
        \cline{1-4}
        $P$ & 0.2 & ${p}_{1}$ & ${p}_{2}$ \\
        \cline{1-4}
        \hline
        \end{tabular}
        
        \end{center}
        
        \begin{jieda}
            根据题意可得 $\left\{  \begin{matrix} {p}_{1} + {p}_{2} = {0.8}, \\  0 \times  {0.2} + 1 \times  {p}_{1} + 2 \times  {p}_{2} = \dfrac{7}{5}, \end{matrix}\right.$ 解得 $\left\{  \begin{array}{l} {p}_{1} = {0.2}, \\  {p}_{2} = {0.6}, \end{array}\right.$
            
            $D\left( {\dfrac{1}{2}X}\right)  = \dfrac{1}{4}D\left( X\right)  = \dfrac{1}{4} \times  \left\lbrack  {{0.2} \times  {\left( {1.4} - 0\right) }^{2} + {0.2} \times  {\left( 1 - {1.4}\right) }^{2} + {0.6} \times  {\left( 2 - {1.4}\right) }^{2}}\right\rbrack = \dfrac{1}{4} \times  \left( {{0.392} + {0.032} + {0.216}}\right)  = {0.16}.$
        \end{jieda}

        \item 已知随机变量 $X$ 满足 $E\left( {{2X} + 3}\right)  = 7,D\left( {{2X} + 3}\right)  = {16}$ , 则下列选项正确的是\dotfill{(\kh{B})}
        \xx{$E\left( X\right)  = \dfrac{7}{2},D\left( X\right)  = \dfrac{13}{2}$}{$E\left( X\right)  = 2,D\left( X\right)  = 4$}{$E\left( X\right)  = 2,D\left( X\right)  = 8$}{$E\left( X\right)  = \dfrac{7}{2},D\left( X\right)  = 8$}
        \begin{jieda}
            因为 $E\left( {{2X} + 3}\right)  = 7,D\left( {{2X} + 3}\right)  = {16}$，所以 ${2E}\left( X\right)  + 3 = 7,{2}^{2}D\left( X\right)  = {16}$，解得 $E\left( X\right)  = 2,D\left( X\right)  = 4$ .
        \end{jieda}

        \item 已知离散型随机变量 $\xi$ 的分布列为
        
        \begin{center}
        
        \begin{tabular}{|c|c|c|c|c|}
        \hline
        $\xi$ & 0 & 1 & 2 & 3 \\
        \cline{1-5}
        $P$ & $m$ & $\dfrac{4}{9}$ & $\dfrac{2}{9}$ & $n$ \\
        \cline{1-5}
        \hline
        \end{tabular}
        
        \end{center}
        
        若 $E\left( \xi \right)  = 1$ ,则 $D\left( {{3\xi } + 1}\right)  =$\blkda{6}
        \begin{jieda}
            由题意知 $m + n = 1 - \dfrac{4}{9} - \dfrac{2}{9} = \dfrac{3}{9}$ ,
            
            由 $E\left( \xi \right)  = 1$ 得 $0 \times  m + 1 \times  \dfrac{4}{9} + 2 \times  \dfrac{2}{9} + {3n} = 1$ ,解得 $n = \dfrac{1}{27},m = \dfrac{8}{27}$
            
            故 $D\left( \xi \right)  = {\left( 0 - 1\right) }^{2} \times  \dfrac{8}{27} + {\left( 1 - 1\right) }^{2} \times  \dfrac{4}{9} + {\left( 2 - 1\right) }^{2} \times  \dfrac{2}{9} + {\left( 3 - 1\right) }^{2} \times  \dfrac{1}{27} = \dfrac{2}{3}$ ,
            
            故 $D\left( {{3\xi } + 1}\right)  = {9D}\left( \xi \right)  = 6$
        \end{jieda}

        \item 设 $a > 0$ ，若随机变量 $\zeta$ 的分布列如下表:

        \begin{center}
        
        \begin{tabular}{|c|c|c|c|}
        \hline
        $\zeta$ & $-1$ & $0$ & $2$ \\
        \cline{1-4}
        $P$ & $a$ & ${2a}$ & ${3a}$ \\
        \cline{1-4}
        \hline
        \end{tabular}
        
        \end{center}
        
        则下列方差中最大的是\dotfill{(\kh{C})}
        \xx{$D\left( \zeta \right)$}{$D\left( \left| \zeta \right| \right)$}{$D\left( {{2\zeta } - 1}\right)$}{$D\left( {2\left| \zeta \right|  - 1}\right)$}
        \begin{jieda}
            由题意,得 $a + {2a} + {3a} = 1$ ,则 $a = \dfrac{1}{6}$，所以 $E\left( \zeta \right)  =  - 1 \times  \dfrac{1}{6} + 0 \times  \dfrac{1}{3} + 2 \times  \dfrac{1}{2} = \dfrac{5}{6},E\left( \left| \zeta \right| \right)  = 1 \times  \dfrac{1}{6} + 0 \times  \dfrac{1}{3} + 2 \times  \dfrac{1}{2} = \dfrac{7}{6}$
            
            所以 $D\left( \zeta \right)  = \dfrac{1}{6} \times  {\left( -1 - \dfrac{5}{6}\right) }^{2} + \dfrac{1}{3} \times  {\left( 0 - \dfrac{5}{6}\right) }^{2} + \dfrac{1}{2} \times  {\left( 2 - \dfrac{5}{6}\right) }^{2} = \dfrac{53}{36},D\left( \left| \zeta \right| \right)  = \dfrac{1}{6} \times  {\left( 1 - \dfrac{7}{6}\right) }^{2} + \dfrac{1}{3} \times  {\left( 0 - \dfrac{7}{6}\right) }^{2} + \dfrac{1}{2} \times {\left( 2 - \dfrac{7}{6}\right) }^{2} = \dfrac{29}{36}$
            
            所以 $D\left( {{2\zeta } - 1}\right)  = {4D}\left( \zeta \right)  = 4 \times  \dfrac{53}{36} = \dfrac{53}{9},D\left( {2\left| \zeta \right|  - 1}\right)  = {4D}\left( \left| \zeta \right| \right)  = \dfrac{29}{9}$，即 $D\left( {{2\zeta } - 1}\right)$ 最大
        \end{jieda}
        
        \item 已知离散型随机变量 $X$ 的分布列 $P\left( {X = \dfrac{k}{5}}\right)  = {ak}\left( {k = 1,2,3,4,5}\right)$ .
        \begin{enumerate}
            \item 求常数 $a$ 的值
            \item 求 $P\left( {X \geq  \dfrac{3}{5}}\right)$
            \item 求随机变量 $\eta  = \left| {{5X} - 3}\right|$ 的分布列及方差
        \end{enumerate}
        \begin{jieda}
            \begin{enumerate}
                \item 由题意得随机变量 $X$ 的分布列如下表所示.
                
                \begin{center}
                
                \begin{tabular}{|c|c|c|c|c|c|}
                \hline
                $X$ & $\dfrac{1}{5}$ & $\dfrac{2}{5}$ & $\dfrac{3}{5}$ & $\dfrac{4}{5}$ & $1$ \\
                \cline{1-6}
                $P$ & $a$ & $2a$ & $3a$ & $4a$ & $5a$ \\
                \cline{1-6}
                \hline
                \end{tabular}
                
                \end{center}
                
                由分布列的性质得 $a + {2a} + {3a} + {4a} + {5a} = 1$ ,解得 $a = \dfrac{1}{75}$
                
                \item $P\left( {X \geq  \dfrac{3}{5}}\right)  = P\left( {X = \dfrac{3}{5}}\right)  + P\left( {X = \dfrac{4}{5}}\right)  + P\left( {X = 1}\right)  = \dfrac{3}{15} + \dfrac{4}{15} + \dfrac{5}{15} = \dfrac{4}{5}$
                
                \item $\eta$ 的所有可能值为 0,1,2,
                
                $\therefore P\left( {\eta  = 0}\right)  = \dfrac{1}{5},P\left( {\eta  = 1}\right)  = \dfrac{2}{5},P\left( {\eta  = 2}\right)  = \dfrac{2}{5}$
                
                所以 $\eta$ 的分布列为:
                
                \begin{center}
                
                \begin{tabular}{|c|c|c|c|}
                \hline
                $\eta$ & 0 & 1 & 2 \\
                \cline{1-4}
                $P$ & $\dfrac{1}{5}$ & $\dfrac{1}{5}$ & $\dfrac{1}{5}$ \\
                \cline{1-4}
                \hline
                \end{tabular}
                
                \end{center}
                
                所以 $E\left( \eta \right)  = 0 \times  \dfrac{1}{5} + 1 \times  \dfrac{2}{5} + 2 \times  \dfrac{2}{5} = \dfrac{6}{5}$ ,
                
                $D\left( \eta \right)  = {\left( 0 - \dfrac{6}{5}\right) }^{2} \times  \dfrac{1}{5} + {\left( 1 - \dfrac{6}{5}\right) }^{2} \times  \dfrac{2}{5} + {\left( 2 - \dfrac{6}{5}\right) }^{2} \times  \dfrac{2}{5} = \dfrac{14}{25}$
            \end{enumerate}
        \end{jieda}
    \end{enumerate}

    \item 两点分布的均值与方差
    \begin{enumerate}
        \item 设随机变量 $X$ 服从两点分布,若 $P\left( {X = 1}\right)  - P\left( {X = 0}\right)  = {0.4}$ , 则 $E\left( X\right)  =$\blkda{$0.7$}
        \begin{jieda}
            由题意得 $P\left( {X = 1}\right)  + P\left( {X = 0}\right)  = 1$
            
            因为 $P\left( {X = 1}\right)  - P\left( {X = 0}\right)  = {0.4}$，所以解得 $P\left( {X = 1}\right)  = {0.7},P\left( {X = 0}\right)  = {0.3}$
            
            所以 $E\left( X\right)  = 1 \times  {0.7} + 0 \times  {0.3} = {0.7}$
        \end{jieda}

        \item 已知离散型随机变量 $X$ 的分布列服从两点分布,满足 $P(X =  0) = \dfrac{2}{{9P}\left( {X = 1}\right) }$ ,且 $P\left( {X = 0}\right)  < P\left( {X = 1}\right)$ ,则 $E\left( X\right)  =$\blkda{$\dfrac{2}{3}$}
        \begin{jieda}
            因为随机变量 $X$ 的分布列服从两点分布，所以 $P\left( {X = 0}\right)  + P\left( {X = 1}\right)  = 1$，则 $P\left( {X = 1}\right)  + \dfrac{2}{{9P}\left( {X = 1}\right) } = 1$ ，解得 $P\left( {X = 1}\right)  = \dfrac{1}{3}$ 或 $\dfrac{2}{3}$
            
            又因 $P\left( {X = 0}\right)  < P\left( {X = 1}\right)$，所以 $P\left( {X = 1}\right)  = \dfrac{2}{3}$ ,则 $P\left( {X = 0}\right)  = \dfrac{1}{3}$，所以 $E\left( X\right)  = \dfrac{2}{3}$
        \end{jieda}
        
        \item 已知随机变量 ${\xi }_{i}$ 服从两点分布，且 $P\left( {{\xi }_{i} = 1}\right)  = {p}_{i}\left( {i = 1,2}\right)$ ， 若 $\dfrac{1}{2} < {p}_{1} < {p}_{2} < 1$ ，则下列判断正确的是\dotfill{(\kh{D})}
        \xx{$E\left( {\xi }_{2}\right)  < D\left( {\xi }_{2}\right)$}{$E\left( {\xi }_{1}\right)  > E\left( {\xi }_{2}\right)$}{$E\left( {\xi }_{1}\right)  < D\left( {\xi }_{1}\right)$}{$D\left( {\xi }_{1}\right)  > D\left( {\xi }_{2}\right)$}
        \begin{jieda}
            $\because E\left( {\xi }_{1}\right)  = {p}_{1},E\left( {\xi }_{2}\right)  = {p}_{2},\therefore E\left( {\xi }_{1}\right)  < E\left( {\xi }_{2}\right)$

            $\because D\left( {\xi }_{1}\right)  = {p}_{1}\left( {1 - {p}_{1}}\right) ,D\left( {\xi }_{2}\right)  = {p}_{2}\left( {1 - {p}_{2}}\right)$
            
            二次函数 $y = x\left( {1 - x}\right)$ 在区间 $\left( {\dfrac{1}{2},1}\right)$ 上单调递减
            
            $\therefore E\left( {\xi }_{2}\right)  > D\left( {\xi }_{2}\right) ,E\left( {\xi }_{1}\right)  > D\left( {\xi }_{1}\right)$ ,且 $D\left( {\xi }_{1}\right)  > D\left( {\xi }_{2}\right)$ .
        \end{jieda}
        
        \item 已知随机变量 ${\xi }_{i}$ 满足 $P\left( {{\xi }_{i} = 0}\right)  = {p}_{i},P\left( {{\xi }_{i} = 1}\right)  = 1 - {p}_{i}$ ，且 $0 < {p}_{i} < \dfrac{1}{2},i = 1,2$ . 若 $E\left( {\xi }_{1}\right)  < E\left( {\xi }_{2}\right)$ ,则\dotfill{(\kh{B})}
        \xx{${p}_{1} < {p}_{2}$ ,且 $D\left( {\xi }_{1}\right)  < D\left( {\xi }_{2}\right)$}{${p}_{1} > {p}_{2}$ ,且 $D\left( {\xi }_{1}\right)  > D\left( {\xi }_{2}\right)$}{${p}_{1} < {p}_{2}$ ,且 $D\left( {\xi }_{1}\right)  > D\left( {\xi }_{2}\right)$}{${p}_{1} > {p}_{2}$ ,且 $D\left( {\xi }_{1}\right)  < D\left( {\xi }_{2}\right)$}
        \begin{jieda}
            由题知变量 ${\xi }_{1},{\xi }_{2}$ 的分布列均为两点分布. 变量 ${\xi }_{1},{\xi }_{2}$ 的分布列如下:
            
            \begin{center}
            
            \begin{tabular}{|c|c|c|c|c|c|c|}
            \hline
            ${\xi }_{1}$ & 0 & 1 & \phantom{X} & ${\xi }_{2}$ & 0 & 1 \\
            \cline{1-7}
            $P$ & ${p}_{1}$ & $1 - {p}_{1}$ & \phantom{X} & $P$ & ${p}_{2}$ & $1 - {p}_{2}$ \\
            \cline{1-7}
            \hline
            \end{tabular}
            
            \end{center}
            
            则 $E\left( {\xi }_{1}\right)  = 1 - {p}_{1},E\left( {\xi }_{2}\right)  = 1 - {p}_{2},D\left( {\xi }_{1}\right)  = {p}_{1}\left( {1 - {p}_{1}}\right) ,D\left( {\xi }_{2}\right)  = {p}_{2}\left( {1 - {p}_{2}}\right)$ ,
            
            由 $E\left( {\xi }_{1}\right)  < E\left( {\xi }_{2}\right)  \Rightarrow  1 - {p}_{1} < 1 - {p}_{2} \Rightarrow  {p}_{1} > {p}_{2}$ ,因为 $0 < {p}_{i} < \dfrac{1}{2},i = 1,2$ ,
            
            函数 $y = x\left( {1 - x}\right)$ 在 $\left( {0,\dfrac{1}{2}}\right)$ 上单调递增,所以 $D\left( {\xi }_{1}\right)  > D\left( {\xi }_{2}\right)$ .
        \end{jieda}
    \end{enumerate}

    \item 决策问题
    \begin{enumerate}
        \item 为更好利用“学习强国”平台开展学习，推动学习型单位建设， 某单位组织开展“学习强国”知识竞赛. 竞赛设置 6 个不同的题目, 参赛人员从中随机抽取 3 个题目进行作答, 若所抽取的 3 个题目全部作答正确，则进入下一轮比赛，否则退出比赛. 对这 6 个题目，某科室的甲能正确作答其中的 4 个题目，乙能正确作答每个题目的概率均为 $\dfrac{2}{3}$ ，且甲乙对每个题目的作答都是相互独立的.
        \begin{enumerate}
            \item 已知甲乙两人总共正确作答 3 个题目，求甲答对 1 道乙答对 2 道的概率
            \item 如果该科室要在甲乙两人中选择一人去参加竞赛，你认为派谁去较为合适？说明理由.
        \end{enumerate}
        \begin{jieda}
            \begin{enumerate}
                \item 甲乙两人总共正确作答 3 个题目，包括:甲笞对 3 道乙笞对 0 道、甲笞对 2 道乙笞对 1 道、甲笞对 1 道乙笞对 2 道.
                
                甲乙两人总共正确作答 3 个题目的概率为 ${P}_{1} = \dfrac{C_{3}^{2}}{C_{6}^{2}} \cdot  {\left( \dfrac{1}{3}\right) }^{3} + \dfrac{C_{2}^{2}C_{2}^{1}}{C_{6}^{2}} \cdot  C_{3}^{1}\dfrac{2}{3}{\left( 1 - \dfrac{2}{3}\right) }^{2} + \dfrac{C_{2}^{1}C_{2}^{2}}{C_{6}^{2}} \cdot  C_{3}^{2}{\left( \dfrac{3}{3}\right) }^{2}\left( {1 - \dfrac{2}{3}}\right)  = \dfrac{31}{135}$
                
                甲答对 1 道题目乙答对 2 道题的概率为 ${P}_{2} = \dfrac{C_{4}^{2}C_{2}^{2}}{C_{6}^{3}} \cdot  C_{3}^{2}{\left( \dfrac{2}{3}\right) }^{2}\left( {1 - \dfrac{2}{3}}\right)  = \dfrac{4}{45}$
                
                所以甲乙两人总共正确作答 3 个题目,甲答对 1 道乙答对 2 道的概率为 $P = \dfrac{{P}_{2}}{{P}_{1}} = \dfrac{12}{37}$
                \item 设甲正确作答题目个数为 $X$ ，则 $X$ 可以取值为 1，2，3.
                
                则 $P\left( {X = 1}\right)  = \dfrac{C_{4}^{1}C_{2}^{2}}{C_{6}^{3}} = \dfrac{1}{5},P\left( {X = 2}\right)  = \dfrac{C_{4}^{2}C_{2}^{1}}{C_{6}^{3}} = \dfrac{3}{5},P\left( {X = 3}\right)  = \dfrac{C_{4}^{3}C_{2}^{0}}{C_{6}^{3}} = \dfrac{1}{5}$ ，$X$ 的分布列为
                
                \begin{center}
                
                \begin{tabular}{|c|c|c|c|}
                \hline
                $X$ & 1 & 2 & 3 \\
                \cline{1-4}
                $P$ & $\dfrac{1}{5}$ & $\dfrac{3}{5}$ & $\dfrac{1}{5}$ \\
                \cline{1-4}
                \hline
                \end{tabular}
                
                \end{center}
                
                $E\left( X\right)  = 1 \times  \dfrac{1}{5} + 2 \times  \dfrac{3}{5} + 3 \times  \dfrac{1}{5} = 2,$
                
                $D\left( X\right)  = {\left( 1 - 2\right) }^{2} \times  \dfrac{1}{5} + {\left( 2 - 2\right) }^{2} \times  \dfrac{3}{5} + {\left( 3 - 2\right) }^{2} \times  \dfrac{1}{5} = \dfrac{2}{5}$
                
                设乙正确作答题数为 $Y$ ，则 $Y \sim  B\left( {3,\dfrac{2}{3}}\right)$ ， $E\left( Y\right)  = 3 \times  \dfrac{2}{3} = 2$ ，
                
                $D\left( Y\right)  = 3 \times  \dfrac{2}{3} \times  \left( {1 - \dfrac{2}{3}}\right)  = \dfrac{2}{3}$
                
                因为 $E\left( X\right)  = E\left( Y\right)$ ,所以甲乙两人实力相当,
                
                而 $D\left( X\right)  < D\left( Y\right)$ ,甲比乙更稳定,因此派甲去较为合适.
            \end{enumerate}
        \end{jieda}
        \item 为选拔奥运会射击选手，对甲、乙两名射手进行选拔测试.已知甲、乙两名射手在一次射击中的得分为两个相互独立的随机变量 $X,Y$ ,甲、乙两名射手在每次射击中击中的环数均大于 6 环,且甲射中 10,9,8,7 环的概率分别为 ${0.5},{3a},a,{0.1}$ ,乙射中 10,9,8 环的概率分别为0.3,0.3,0.2.
        \begin{enumerate}
            \item 求 $X,Y$ 的概率分布
            \item 求 $X,Y$ 的数学期望与方差，以此比较甲、乙的射击技术并从中选拔一人.
        \end{enumerate}
        \begin{jieda}
            \begin{enumerate}
                \item 依题意， ${0.5} + {3a} + a + {0.1} = 1$ ，解得 $a = {0.1}$ ，因为乙射中 10, 9, 8 环的概率分别为0.3,0.3,0.2，所以乙射中 7 环的概率为 $1 - \left( {{0.3} + {0.3} + {0.2}}\right)  = {0.2}$，故$X$ 的概率分布为:
                
                \begin{center}
                
                \begin{tabular}{|c|c|c|c|c|}
                \hline
                $X$ & 10 & 9 & 8 & 7 \\
                \cline{1-5}
                $P$ & 0.5 & 0.3 & 0.1 & 0.1 \\
                \cline{1-5}
                \hline
                \end{tabular}
                
                \end{center}
                
                $Y$ 的概率分布为:
                
                \begin{center}
                
                \begin{tabular}{|c|c|c|c|c|}
                \hline
                $Y$ & 10 & 9 & 8 & 7 \\
                \cline{1-5}
                $P$ & 0.3 & 0.3 & 0.2 & 0.2 \\
                \cline{1-5}
                \hline
                \end{tabular}
                
                \end{center}

                \item 由 (1) 可得$E\left( X\right)  = {10} \times  {0.5} + 9 \times  {0.3} + 8 \times  {0.1} + 7 \times  {0.1} = {9.2}$ (环)，$E\left( Y\right)  = {10} \times  {0.3} + 9 \times  {0.3} + 8 \times  {0.2} + 7 \times  {0.2} = {8.7}$ (环),
                
                $D\left( X\right)  = {\left( {10} - {9.2}\right) }^{2} \times  {0.5} + {\left( 9 - {9.2}\right) }^{2} \times  {0.3} + {\left( 8 - {9.2}\right) }^{2} \times  {0.1} + {\left( 7 - {9.2}\right) }^{2} \times  {0.1} = {0.96}$ ,
                
                $D\left( Y\right)  = {\left( {10} - {8.7}\right) }^{2} \times  {0.3} + {\left( 9 - {8.7}\right) }^{2} \times  {0.3} + {\left( 8 - {8.7}\right) }^{2} \times  {0.2} + {\left( 7 - {8.7}\right) }^{2} \times  {0.2} = {1.21}$ ,
                
                由于 $E\left( X\right)  > E\left( Y\right)$ ,说明甲平均射中的环数比乙高，又因为 $D\left( X\right)  < D\left( Y\right)$ ,说明甲射中的环数比乙集中,比较稳定，所以，甲比乙的技术好，故应选拔甲射手参加奥运会.
            \end{enumerate}
        \end{jieda}
        \item 为回馈顾客，某商场拟通过摸球兑奖的方式对 500 位顾客进行奖励，规定:每位顾客从一个装有 4 个标有面值的球的袋中一次性随机摸出 2 个球，球上所标的面值之和为该顾客所获的奖励额.
        \begin{enumerate}
            \item 若袋中所装的 4 个球中有 1 个所标的面值为 45 元，其余 3 个均为 15 元，求顾客所获的奖励额为 60 元的概率
            \item 商场对奖励总额的预算是 30000 元，为了使顾客得到的奖励总额尽可能符合商场的预算且每位顾客所获的奖励额相对均衡，请从如下两种方案中选择一种，并说明理由.方案一:袋中的 4 个球由 2 个标有面值 15 元和 2 个标有面值 45 元的两种球组成; 方案二: 袋中的 4 个球由 2 个标有面值 20 元和 2 个标有面值 40 元的两种球组成.
        \end{enumerate}
        \begin{jieda}
            \begin{enumerate}
                \item 设顾客的奖励额为 $X$ ,依题意得 $P\left( {X = {60}}\right)  = \dfrac{C_{1}^{1}C_{3}^{1}}{C_{4}^{2}} = \dfrac{1}{2}$
                \item 根据方案一，设顾客的奖励额为 $X$ ,其可能取值为 30, 60 ,90
                
                $P\left( {{X}_{1} = {30}}\right)  = \dfrac{C_{2}^{2}}{C_{4}^{2}} = \dfrac{1}{6},P\left( {{X}_{1} = {60}}\right)  = \dfrac{C_{2}^{1}C_{2}^{1}}{C_{4}^{2}} = \dfrac{4}{6} = \dfrac{2}{3},P\left( {{X}_{1} = {90}}\right)  = \dfrac{C_{2}^{2}}{C_{4}^{2}} = \dfrac{1}{6}$
                
                $E\left( {X}_{1}\right)  = {30} \times  \dfrac{1}{6} + {60} \times  \dfrac{2}{3} + {90} \times  \dfrac{1}{6} = {60}$
                
                $D\left( {X}_{1}\right)  = {\left( {30} - {60}\right) }^{2} \times  \dfrac{1}{6} + {\left( {60} - {60}\right) }^{2} \times  \dfrac{2}{3} + {\left( {90} - {60}\right) }^{2} \times  \dfrac{1}{6} = {300}$
                
                根据方案二,设顾客的奖励额为 ${X}_{2}$ ,其可能取值为40,60,80
                
                $P\left( {{X}_{2} = {40}}\right)  = \dfrac{C_{2}^{2}}{C_{4}^{2}} = \dfrac{1}{6},P\left( {{X}_{2} = {60}}\right)  = \dfrac{C_{2}^{2}C_{2}^{1}}{C_{4}^{2}} = \dfrac{4}{6} = \dfrac{2}{3},P\left( {{X}_{2} = {80}}\right)  = \dfrac{C_{2}^{2}}{C_{4}^{2}} = \dfrac{1}{6}$
                
                $E\left( {X}_{2}\right)  = {40} \times  \dfrac{1}{6} + {60} \times  \dfrac{2}{3} + {80} \times  \dfrac{1}{6} = {60}$
                
                $D\left( {X}_{2}\right)  = {\left( {40} - {60}\right) }^{2} \times  \dfrac{1}{6} + {\left( {60} - {60}\right) }^{2} \times  \dfrac{2}{3} + {\left( {80} - {60}\right) }^{2} \times  \dfrac{1}{6} = \dfrac{400}{3}$
                
                商场对奖励总额的预算是 30000 元，故每个顾客平均奖励额最多为 60，两方案均符合要求，但方案二奖励的方差比方案一小，所以应选择方案二.
            \end{enumerate}
        \end{jieda}
        \item 某短视频软件经过几年的快速发展，深受人们的喜爱，该软件除了有娱乐属性外，也可通过平台推送广告. 某公司为了宣传新产品，现有以下两种宣传方案:
        
        方案一:投放该平台广告，据市场调研，其收益 $X$ 分别为 0 元，20 万元，40 万元，且 $P\left( {X = {20}}\right)  = {0.3}$ ， 期望 $E\left( X\right)  = {30}$ .
        
        方案二: 投放传统广告,据市场调研,其收益 $Y$ 分别为 10 万元,20 万元,30 万元,其概率依次为0.3,0.4,0.3.

        \begin{enumerate}
            \item 请写出方案一的分布列，并求方差 $D\left( X\right)$
            \item 请你根据所学的统计知识给出建议，该公司宣传应该投放哪种广告？并说明你的理由.
        \end{enumerate}
        \begin{jieda}
            \begin{enumerate}
                \item 设 $P\left( {X = 0}\right)  = a,P\left( {X = {40}}\right)  = b$ ,

                依题意得 $a + b + {0.3} = 1$，又 $E\left( X\right)  = 0 \times  a + {20} \times  {0.3} + {40b} = {30}$，解得: $a = {0.1}$ ， $b = {0.6}$，故$X$ 的分布列为
                
                \begin{center}
                
                \begin{tabular}{|c|c|c|c|}
                \hline
                $X$ & 0 & 20 & 40 \\
                \cline{1-4}
                $P$ & 0.1 & 0.3 & 0.6 \\
                \cline{1-4}
                \hline
                \end{tabular}
                
                \end{center}
                
                则 $D\left( X\right)  = {\left( 0 - {30}\right) }^{2} \times  {0.1} + {\left( {20} - {30}\right) }^{2} \times  {0.3} + {\left( {40} - {30}\right) }^{2} \times  {0.6} = {180}$ .
                \item 由题得 $Y$ 的分布列为
                
                \begin{center}
                
                \begin{tabular}{|c|c|c|c|}
                \hline
                $Y$ & 10 & 20 & 30 \\
                \cline{1-4}
                $P$ & 0.3 & 0.4 & 0.3 \\
                \cline{1-4}
                \hline
                \end{tabular}
                
                \end{center}
                
                则 $E\left( Y\right)  = {10} \times  {0.3} + {20} \times  {0.4} + {30} \times  {0.3} = {20}$ ,
                
                $D\left( Y\right)  = {\left( {10} - {20}\right) }^{2} \times  {0.3} + {\left( {20} - {20}\right) }^{2} \times  {0.4} + {\left( {30} - {20}\right) }^{2} \times  {0.3} = {60}.$
                
                由 $E\left( X\right)  > E\left( Y\right)$ 可知采用平台广告投放期望收益较大,又 $D\left( X\right)  > D\left( Y\right)$ ,说明平台广告投放的风险较高.
                
                综上所述, 如果公司期望高收益, 选择平台广告; 如果公司期望收益稳定, 选择传统广告.
            \end{enumerate}
        \end{jieda}
    \end{enumerate}
    \item 根据二项分布求概率
    \begin{enumerate}
        \item 已知随机变量 $X \sim  B\left( {4,p}\right)$ ，若 $P\left( {X = 2}\right)  = \dfrac{8}{27}$ ，则 $p =$\blkda{$\dfrac{1}{3}$ 或 $\dfrac{2}{3}$}
        \begin{jieda}
            因为 $X \sim  B\left( {4,p}\right)$ ,则 $P\left( {X = 2}\right)  = {{C}}_{4}^{2}{p}^{2}{\left( 1 - p\right) }^{2} = 6{p}^{2}{\left( 1 - p\right) }^{2} = \dfrac{8}{27}$
            
            且 $0 < p < 1$ ,整理可得 $p\left( {1 - p}\right)  = \dfrac{2}{9}$ ,解得 $p = \dfrac{1}{3}$ 或 $\dfrac{2}{3}$ .
        \end{jieda}
        \item 已知随机变量 $X \sim  B\left( {4,p}\right)$ ，其中 $0 < p < 1$ ，若 $P\left( {X \leq  3}\right)  = \dfrac{15}{16}$ ， 则 $P\left( {X = 3}\right)  =$\blkda{$\dfrac{1}{4}$}
        \begin{jieda}
            由二项分布的知识得 $P\left( {X \leq  3}\right)  = 1 - P\left( {X = 4}\right)  = 1 - {p}^{4} = \dfrac{15}{16}$，得 ${p}^{4} = \dfrac{1}{16}$ ，又 $0 < p < 1$ ，所以 $p = \dfrac{1}{2}$
            
            所以 $P\left( {X = 3}\right)  = {{C}}_{4}^{3} \times  {\left( \dfrac{1}{2}\right) }^{3} \times  \left( {1 - \dfrac{1}{2}}\right)  = \dfrac{1}{4}$
        \end{jieda}
        
        \item 如图所示，已知一质点在外力的作用下，从原点 $O$ 出发，每次向左移动的概率为 $\dfrac{3}{4}$ . 向右移动的概率为 $\dfrac{1}{4}$ . 若该质点每次移动一个单位长度，设经过 5 次移动后，该质点位于 $X$ 的位置,则 $P\left( {X > 0}\right)  =$\blkda{$\dfrac{53}{512}$}
        
        \begin{center}
        \includegraphics[width=0.3\textwidth]{bodyxb/src/bo_d67fmln7aajc739ac57g_0_144_1698_400_61_0.jpg}
        \end{center}
        \begin{jieda}
            由题意可知,当 $X > 0$ 时, $X$ 的可能取值为 1,3,5,且 $\dfrac{X+5}{2} \sim  B\left( {5,\dfrac{1}{4}}\right)$
            
            $P\left( {X > 0}\right)  = P\left( {X = 5}\right)  + P\left( {X = 3}\right)  + P\left( {X = 1}\right) = {\left( \dfrac{1}{4}\right) }^{5} + {{C}}_{5}^{1}{\left( \dfrac{1}{4}\right) }^{4} \cdot  \left( \dfrac{3}{4}\right)  + {{C}}_{5}^{2}{\left( \dfrac{1}{4}\right) }^{3} \cdot  {\left( \dfrac{3}{4}\right) }^{2} = \dfrac{53}{512}$
        \end{jieda}
        
        \item 已知随机变量 $X \sim  B\left( {3,p}\right)$ . 若 $\dfrac{1}{2} \leq  p < 1$ ，则 $P\left( {X \geq  \dfrac{3}{2}}\right)$ 的取值范围是\blkda{$\left\lbrack  {\dfrac{1}{2},1}\right)$}
        \begin{jieda}
            由 $X \sim  B\left( {3,p}\right)$ ,则 $P\left( {X = k}\right)  = {{C}}_{3}^{k}{p}^{k}{\left( 1 - p\right) }^{3 - k},k = 0,1,2,3$
            
            $P\left( {X \geq  \dfrac{3}{2}}\right)  = P\left( {X = 2}\right)  + P\left( {X = 3}\right)  = {{C}}_{3}^{2}{p}^{2}{\left( 1 - p\right) }^{1} + {{C}}_{3}^{3}{p}^{3} =  - 2{p}^{3} + 3{p}^{2}$，且 $\dfrac{1}{2} \leq  p < 1$ ,
            
            构造函数 $y =  - 2{x}^{3} + 3{x}^{2},\dfrac{1}{2} \leq  x < 1$ ,其导函数为 ${y}^{\prime } =  - 6{x}^{2} + {6x}$ ,
            
            由于 $\dfrac{1}{2} \leq  x < 1,{y}^{\prime } > 0$ ,故函数 $y =  - 2{x}^{3} + 3{x}^{2}$ 在区间 $\left\lbrack  {\dfrac{1}{2},1}\right)$ 上单调递增；
            
            当 $x = \dfrac{1}{2}$ 时，取最小值 $\dfrac{1}{2}$ ，当 $x = 1$ 时，函数值为 1，所以 $P\left( {X \geq  \dfrac{3}{2}}\right)  \in  \left\lbrack  {\dfrac{1}{2},7}\right)$ 
        \end{jieda}
    \end{enumerate}
    \item 服从二项分布的随机变量概率最大问题
    \begin{enumerate}
        \item 某学校在假期组织 30 位学生前往北京、上海、广州、深圳、杭州、苏州 、成都、重庆 8 个城市参加研学活动.每个学生可自由选择 8 个城市中的任意 1 个(不要求每个城市必须要有学生选择). 若每位学生选择去每个城市的概率都相等且互不影响，则有\blkda{7}个学生选择前往北京或上海研学的概率最大.
        \begin{jieda}
            设有 $X$ 个学生选择前往北京或上海研学，由题意可得每个学生选择前往北京或上海研学的概率 $P = \dfrac{2}{8} = \dfrac{1}{4}$，则 $X \sim  B\left( {{30},\dfrac{1}{4}}\right)$
            
            设有 $k$ 个学生选择前往北京或上海研学的概率最大，则 $\left\{  \begin{array}{l} P\left( {X = k}\right)  \geq  P\left( {X = k + 1}\right) \\  P\left( {X = k}\right)  \geq  P\left( {X = k - 1}\right)  \end{array}\right.$，即$\left\{\begin{array}{l}
                 {{C}}_{30}^{k} \cdot  {\left( \dfrac{1}{4}\right) }^{k} \cdot  {\left( \dfrac{3}{4}\right) }^{{30} - k} \geq  {{C}}_{30}^{k + 1} \cdot  {\left( \dfrac{1}{4}\right) }^{k + 1} \cdot  {\left( \dfrac{3}{4}\right) }^{{29} - k}  \\
                 {{C}}_{30}^{k} \cdot  {\left( \dfrac{1}{4}\right) }^{k} \cdot  {\left( \dfrac{3}{4}\right) }^{{30} - k} \geq  {{C}}_{30}^{k - 1} \cdot  {\left( \dfrac{1}{4}\right) }^{k - 1} \cdot  {\left( \dfrac{3}{4}\right) }^{{31} - k} 
            \end{array}\right.$
            
            即 $\left\{  \begin{array}{l} 
            \dfrac{30!}{k!(30-k)!} \times 3 \geq \dfrac{30!}{(k+1)!(29-k)!} \\
            \dfrac{30!}{k!(30-k)!} \geq \dfrac{30!}{(k-1)!(31-k)!} \times 3
            \end{array}\right.$，即 $\left\{  \begin{array}{l} 
            \dfrac{k+1}{(30-k)} \geq \dfrac{1}{3} \\
            \dfrac{(31-k)}{k} \geq 3
            \end{array}\right.$
            
            解得 $\dfrac{27}{4} \leq  k \leq  \dfrac{31}{4}$，又 $k \in  {\mathbf{N}}^{ * }$ ,所以 $k = 7$，所以有 7 个学生选择前往北京或上海研学的概率最大.
        \end{jieda}
        
        \item $X \sim  B\left( {5,\dfrac{1}{4}}\right)$ ,则 $P\left( {\xi  = k}\right)  = {{C}}_{5}^{k}{\left( \dfrac{1}{4}\right) }^{k}{\left( \dfrac{3}{4}\right) }^{5 - k}$ 取最大值时 $k$ 的值为\blkda{1}
        \begin{jieda}
            由已知 $P\left( {\xi  = k}\right)  = {{C}}_{5}^{k}{\left( \dfrac{1}{4}\right) }^{k}{\left( \dfrac{3}{4}\right) }^{5 - k},k = 0,1,2,3,4,5$ ,

            所以由 $\left\{  \begin{array}{l} P\left( {\xi  = k}\right)  \geq  P\left( {\xi  = k + 1}\right) \\  P\left( {\xi  = k}\right)  \geq  P\left( {\xi  = k - 1}\right)  \end{array}\right.$
            
            得:
            
            ${{C}}_{5}^{k}{\left( \dfrac{7}{4}\right) }^{k}{\left( \dfrac{3}{4}\right) }^{5 - k} \geq  {{C}}_{5}^{k + 1}{\left( \dfrac{7}{4}\right) }^{k + 1}{\left( \dfrac{3}{4}\right) }^{4 - k}$
            
            ${{C}}_{5}^{k}{\left( \dfrac{1}{4}\right) }^{k}{\left( \dfrac{3}{4}\right) }^{5 - k} \geq  {{C}}_{5}^{k - 1}{\left( \dfrac{1}{4}\right) }^{k - 1}{\left( \dfrac{3}{4}\right) }^{6 - k}$ ,
            
            解得 $\dfrac{1}{2} \leq  k \leq  \dfrac{3}{2}$ ，又因为 $k \in  {\mathbf{N}}^{ * }$ ，所以 $k = 1$ .
        \end{jieda}

        \item 已知随机变量 $X \sim  B\left( {n,{0.5}}\right)$ ，当且仅当 $k = 4$ 时， $P\left( {X = k}\right)$ 取得最大值，则 $n =$\blkda{8}
        \begin{jieda}
            由题得 $P\left( {X = k}\right)  = {{C}}_{n}^{k}{\left( \dfrac{1}{2}\right) }^{n},k = 0,1,\cdots ,n$
            
            由题知在 ${{C}}_{n}^{0}{\left( \dfrac{1}{2}\right) }^{n},{{C}}_{n}^{1}{\left( \dfrac{1}{2}\right) }^{n},\cdots ,{{C}}_{n}^{n}{\left( \dfrac{1}{2}\right) }^{n}$ 中,最大值只有 ${{C}}_{n}^{4}{\left( \dfrac{1}{2}\right) }^{n}$，即在 ${{C}}_{n}^{0},{{C}}_{n}^{1},\cdots ,{{C}}_{n}^{n}$ 中,最大值只有 ${{C}}_{n}^{4}$ ,由二项式系数的对称性可知 $n = 8$ .
        \end{jieda}
        
        \item 如图是一块高尔顿板的示意图，在一块木板上钉着若干排相互平行但相互错开的圆柱形小木钉，小木钉之间留有适当的空隙作为通道，前面挡有一块玻璃. 将小球从顶端放入，小球下落的过程中，每次碰到小木钉后都等可能地向左或向右落下，最后落入底部的格子中. 记格子从左到右的编号分别为 $0,1,2,\cdots ,{10}$ ,用 $X$ 表示小球最后落入格子的号码,若 $P\left( {X = k}\right)  \leq  P\left( {X = {k}_{0}}\right)$ , 则 ${k}_{0} =$\blkda{5}

        \begin{center}
        \includegraphics[width=0.3\textwidth]{bodyxb/src/bo_d67fmln7aajc739ac57g_3_140_989_387_395_0.jpg}
        \end{center}
        \begin{jieda}
            每下落一层向左或向右落下等可能,概率均为 $\dfrac{1}{2}$，每一层均要乘以 $\dfrac{1}{2}$ ,共做 10 次选择，故 $X$ 服从二项分布, $X \sim  B\left( {{10},\dfrac{1}{2}}\right)$，由题知在 ${{C}}_{10}^{0}{\left( \dfrac{1}{2}\right) }^{10},{{C}}_{10}^{1}{\left( \dfrac{1}{2}\right) }^{10},\cdots ,{{C}}_{10}^{10}{\left( \dfrac{1}{2}\right) }^{10}$ 中,最大值只有 ${{C}}_{10}^{k_0}{\left( \dfrac{1}{2}\right) }^{10}$，即在 ${{C}}_{10}^{0},{{C}}_{10}^{1},\cdots ,{{C}}_{10}^{10}$ 中,最大值只有 ${{C}}_{10}^{k_0}$ ,由二项式系数的对称性可知 $k_0 = 5$ .
        \end{jieda}
    \end{enumerate}
    \item 二项分布的均值与方差
    \begin{enumerate}
        \item 已知随机变量 $X \sim  B\left( {4,\dfrac{1}{3}}\right)$ ，下列表达式正确的是\dotfill{(\kh{C})}
        \xx{$P\left( {X = 2}\right)  = \dfrac{4}{81}$}{$E\left( {{3X} + 1}\right)  = 4$}{$D\left( {{3X} + 1}\right)  = 8$}{$D\left( X\right)  = \dfrac{4}{9}$}
        \begin{jieda}
            因为 $X \sim  B\left( {4,\dfrac{1}{3}}\right)$ ,所以 $E\left( X\right)  = 4 \times  \dfrac{1}{3} = \dfrac{4}{3},D\left( X\right)  = 4 \times  \dfrac{1}{3} \times  \left( {1 - \dfrac{1}{3}}\right)  = \dfrac{8}{9}$
            
            因此 $E\left( {{3X} + 1}\right)  = {3E}\left( X\right)  + 1 = 3 \times  \dfrac{4}{3} + 1 = 5$ ,
            
            $D\left( {{3X} + 1}\right)  = {3}^{2} \cdot  D\left( X\right)  = 9 \times  \dfrac{8}{9} = 8,$
            
            因此选项 B、D 不正确，选项 C 正确；
            
            又因为 $P\left( {X = 2}\right)  = {{C}}_{4}^{2}{\left( \dfrac{1}{3}\right) }^{2}{\left( 1 - \dfrac{1}{3}\right) }^{2} = \dfrac{8}{27}$ 所以选项 A 不正确
        \end{jieda}
        \item 设随机变量 $X \sim  B\left( {{12},p}\right)$ ，若 $E\left( X\right)  \leq  4$ ，则 $D\left( X\right)$ 的最大值为\blkda{$\dfrac{8}{3}$}
        \begin{jieda}
            随机变量 $X \sim  B\left( {{12},p}\right)$ ,由 $E\left( X\right)  \leq  4$ 得 $0 < {12p} \leq  4$ ,解得 $0 < p \leq  \dfrac{1}{3}$，$D\left( X\right)  = {12p}\left( {1 - p}\right)  =  - {12}{p}^{2} + {12p}$

            易知二次函数 $y =  - {12}{x}^{2} + {12x}$ 的开口向下，对称轴为 $x = \dfrac{1}{2}$ 所以在 $\left( {0,\dfrac{1}{3}}\right\rbrack$ 上单调递增，
            
            于是 $p = \dfrac{1}{3}$ 时 $D\left( X\right)$ 取得最大值,即最大值为 ${12p}\left( {1 - p}\right)  = {12} \times  \dfrac{1}{3} \times  \dfrac{2}{3} = \dfrac{8}{3}$
        \end{jieda}
        \item 下图是一块高尔顿板示意图:在一块木块上钉着若干排互相平行但相互错开的圆柱形小木钉，小木钉之间留有适当的空隙作为通道，前面挡有一块玻璃，将小球从顶端放入，小球在下落过程中，每次碰到小木钉后都等可能地向左或向右落下，最后落入底部的格子中， 格子从左到右分别编号为 0,1,2,3,4,5 用 $X$ 表示小球落入格子的号码，则下面计算错误的是\dotfill{(\kh{B})}
        
        \begin{center}
        \includegraphics[width=0.2\textwidth]{bodyxb/src/bo_d67fmln7aajc739ac57g_5_154_472_269_327_0.jpg}
        \end{center}
        \xx{$P\left( {X = 0}\right)  = \dfrac{1}{32}$}{$P\left( {X = 5}\right)  = \dfrac{1}{64}$}{$E\left( X\right)  = \dfrac{5}{2}$}{$D\left( X\right)  = \dfrac{5}{4}$}
        \begin{jieda}
            设 $A =$ “向右下落”，则 $A =$ “向左下落”， $P\left( A\right)  = P\left( \bar{A}\right)  = \dfrac{1}{2}$ ，
            
            因为小球最后落入格子的号码 $X$ 等于事件 $A$ 发生的次数，
            
            而小球下落的过程中共碰撞小木钉 5 次，所以 $X \sim  B\left( {5,\dfrac{1}{2}}\right)$ ，
            
            对于 A: $P\left( {X = 0}\right)  = {\left( 1 - \dfrac{1}{2}\right) }^{5} = \dfrac{1}{32}$ ,故 A 正确;
            
            对于 B: $P\left( {X = 5}\right)  = {\left( \dfrac{1}{2}\right) }^{5} = \dfrac{1}{32}$ ,故 B 错误;
            
            对于 ${C} : E\left( X\right)  = 5 \times  \dfrac{1}{2} = \dfrac{5}{2}$ 故 ${C}$ 正确；
            
            对于 ${D} : D\left( X\right)  = 5 \times  \dfrac{1}{2}\left( {1 - \dfrac{1}{2}}\right)  = \dfrac{5}{4}$ . 故 ${D}$ 正确
        \end{jieda}
        
        \item 一只小虫从数轴上的原点出发爬行，若一次爬行过程中，小虫等概率地向前或向后爬行 1 个单位,设爬行 $n$ 次后小虫所在位置对应的数为随机变量 ${\xi }_{n}$ ,则下列说法错误的是\dotfill{(\kh{C})}
        \xx{$E\left( {\xi }_{n}\right)  = 0$}{$D\left( {\xi }_{n}\right)  = n$}{$P\left( {{\xi }_{100} = 0}\right)  < P\left( {{\xi }_{100} = 2}\right)$}{$P\left( {{\xi }_{102} = 0}\right)  < P\left( {{\xi }_{100} = 0}\right)$}

        \begin{jieda}
            对于$A, B$，设$X_n = \dfrac{{\xi }_{n}+n}{2}$，则$X_n \sim B\left(n, \dfrac{1}{2}\right)$，$\left\{\begin{array}{l}
                 E(X_n) = \dfrac{n}{2} \\
                 D(X_n) = \dfrac{n}{4}
            \end{array}\right.$，故$\left\{\begin{array}{l}
                 E(\xi_n) = 0 \\
                 D(\xi_n) = n
            \end{array}\right.$，故$A, B$正确

            对于 ${C}$ ,因为 $P\left( {{\xi }_{100} = 0}\right)  = {{C}}_{100}^{50}{\left( \dfrac{1}{2}\right) }^{100},P\left( {{\xi }_{100} = 2}\right)  = {{C}}_{100}^{51}{\left( \dfrac{1}{2}\right) }^{100}$ ,
            
            所以 $\dfrac{P\left( {{\xi }_{100} = 0}\right) }{P\left( {{\xi }_{100} = 2}\right) } = \dfrac{{{C}}_{100}^{50}}{{{C}}_{100}^{51}} = \dfrac{51}{50} > 1$ ,所以 $P\left( {{\xi }_{100} = 0}\right)  > P\left( {{\xi }_{100} = 2}\right)$ ,所以 ${C}$ 错误
            
            对于 ${D}$ ，因为 $P\left( {{\xi }_{100} = 0}\right)  = {{C}}_{100}^{50}{\left( \dfrac{1}{2}\right) }^{100},P\left( {{\xi }_{102} = 0}\right)  = {{C}}_{102}^{51}{\left( \dfrac{1}{2}\right) }^{102}$ ，
            
            所以 $\dfrac{P\left( {{\xi }_{100} = 0}\right) }{P\left( {{\xi }_{102} = 0}\right) } = \dfrac{4{{C}}_{100}^{50}}{{{C}}_{102}^{51}} = \dfrac{4 \times  {51} \times  {51}}{{102} \times  {101}} = \dfrac{102}{101} > 1$ ,
            
            所以 $P\left( {{\xi }_{102} = 0}\right)  < P\left( {{\xi }_{100} = 0}\right)$ ,所以 D 正确
        \end{jieda}
    \end{enumerate}
    \item 二项分布的实际应用
    \begin{enumerate}
        \item 已知某计算机网络的服务器有三台设备，只要有一台能正常工作， 计算机网络就不会断掉. 如果三台设备各自能正常工作的概率都为 0.8 , 它们之间互相不影响. 设能正常工作的设备数为 $X$ .
        \begin{enumerate}
            \item 求 $X$ 的分布列
            \item 求 $E\left( X\right)$ 和 $D\left( X\right)$
            \item 求计算机网络不会断掉的概率
        \end{enumerate}
        \begin{jieda}
            \begin{enumerate}
                \item 由题意得 $X$ 的可能取值为0,1,2,3，且 $X \sim  B\left( {3,0,8}\right)$
                
                $P\left( {X = 0}\right)  = {{C}}_{3}^{0} \times  {0.8}^{0} \times  {\left( 1 - {0.8}\right) }^{3} = {0.008}$
                
                $P\left( {X = 1}\right)  = {{C}}_{3}^{1} \times  {0.8}^{1} \times  {\left( 1 - {0.8}\right) }^{2} = {0.096}$
                
                $P\left( {X = 2}\right)  = {{C}}_{3}^{2} \times  {0.8}^{2} \times  {\left( 1 - {0.8}\right) }^{1} = {0.384}$
                
                $P\left( {X = 3}\right)  = {{C}}_{3}^{3} \times  {0.8}^{3} \times  {\left( 1 - {0.8}\right) }^{0} = {0.512}$
                
                所以 $X$ 的分布列如下.
                
                \begin{center}
                \begin{tabular}{|c|c|c|c|c|}
                \hline
                $X$ & 0 & 1 & 2 & 3 \\
                \cline{1-5}
                $P$ & 0.008 & 0.096 & 0.384 & 0.512 \\
                \cline{1-5}
                \hline
                \end{tabular}
                \end{center}
                
                \item 因为 $X \sim  B\left( {3,0,8}\right)$ ,所以 $E\left( X\right)  = 3 \times  {0.8} = {2.4},D\left( X\right)  = 3 \times  {0.8} \times  \left( {1 - {0.8}}\right)  = {0.48}$ .
                
                \item 要使得计算机网络不会断掉，也就是要求能正常工作的设备至少有一台，即 $X \geq  1$ ，因此所求概率为 $P\left( {X \geq  1}\right)  = 1 - P\left( {X < 1}\right)  = 1 - P\left( {X = 0}\right)  = 1 - {0.008} = {0.992}$ .
            \end{enumerate}
        \end{jieda}
        \item 某工厂生产的产品分为一等品、二等品和三等品.已知生产一件产品为一等品、二等品、三等品的概率分别为 ${P}_{1},{P}_{2},{P}_{3}$ ,且 ${P}_{1} + {P}_{2} + {P}_{3} = 1$ . 从该工厂生产的产品中随机抽取 $n$ 件,设其中一等品的数量为 $X$ ,二等品的数量为 $Y$ .
        \begin{enumerate}
            \item 已知 $X$ 的数学期望 $E\left( X\right)  = 4$ ， $X$ 的方差 $D\left( X\right)  = {2.4}$ ，求 ${P}_{1}$ 的值
            \item 若 $n = {20}$ ,且 $P\left( {Y = 6}\right)  = P\left( {Y = 8}\right)$ ,求 ${P}_{2}$ 的值
            \item 已知 ${P}_{1} = {0.4},{P}_{2} = {0.3}$ ，在抽取的 $n$ 件商品中，一等品和二等品的数量之和为 $M.M$ 的数学期望是否有最大值, 若有, 求出最大值; 若没有, 说明理由
        \end{enumerate}
        \begin{jieda}
            \begin{enumerate}
                \item 由题意知 $X \sim  B\left( {n,{P}_{1}}\right)$ ,则 $\left\{  \begin{matrix} E\left( X\right)  = n{P}_{1} = 4 \\  D\left( X\right)  = n{P}_{1}\left( {1 - {P}_{1}}\right)  = {2.4} \end{matrix}\right.$，解得 $n = {10},{P}_{1} = {0.4}$ 
                
                \item 由 $Y \sim  B\left( {{20},{P}_{2}}\right)$ ，则 $P\left( {Y = k}\right)  = {{C}}_{20}^{k}{P}_{2}^{k}{\left( 1 - {P}_{2}\right) }^{{20} - k}\left( {k = 0,1,2,\cdots ,{20}}\right)$ ，由 $P\left( {Y = 6}\right)  = P\left( {Y = 8}\right)$ 得, ${{C}}_{20}^{6} \cdot  {P}_{2}^{6} \cdot  {\left( 1 - {P}_{2}\right) }^{14} = {{C}}_{20}^{8} \cdot  {p}_{2}^{8} \cdot  {\left( 1 - {P}_{2}\right) }^{12}$，化简可得 $\dfrac{{{C}}_{20 }^{6}}{{{C}}_{20 }^{8}} = \dfrac{{P}_{2}^{2}}{{\left( 1 - {P}_{2}\right) }^{2}}$ ,即 $\dfrac{4}{13} = \dfrac{{P}_{2}^{2}}{{\left( 1 - {P}_{2}\right) }^{2}}$ ,解得 ${P}_{2} = \dfrac{2\sqrt{13} - 4}{9}$ 
                
                \item 由题意知, $M \sim  B\left( {n,{P}_{1} + {P}_{2}}\right)$ ,又 ${P}_{1} = {0.4},{P}_{2} = {0.3}$，所以 $M \sim  B\left( {n,{0.7}}\right)$ ,则 $E\left( M\right)  = {0.7n}$ ,当 $n$ 增大时, $E\left( M\right)$ 也增大.
                
                所以,当 $n \rightarrow   + \infty ,E\left( M\right)  \rightarrow   + \infty$ ,故 $M$ 的数学期望没有最大值.
                
                但在实际情境中， $n$ 的取值是有限的，比如取工厂的总产量时， $E\left( M\right)$ 取最大值.
            \end{enumerate}
        \end{jieda}
        \item 某批产品由一等品、二等品、三等品及次品构成. 随机抽取 20 件, 统计情况如下表:
        
        \begin{center}
        \begin{tabular}{|c|c|c|c|c|}
        \hline
        等级 & 一等品 & 二等品 & 三等品 & 次品 \\
        \cline{1-5}
        件数 & 10 & 6 & 3 & 1 \\
        \cline{1-5}
        \hline
        \end{tabular}
        
        \end{center}
        
        以频率估计概率.
        \begin{enumerate}
            \item 若从这批产品中任取一件，试估计“其为一等品或二等品”的概率；
            
            \item 在抽取的 20 件产品中，随机抽取 3 件，求其中“既有一等品又有二等品”的概率；
            
            \item 若改进技术后，产品为一等品的概率提升为 $\dfrac{3}{5}$ 从这批产品中随机抽取 10 件产品，其中恰好有 5 件一等品的概率记为 ${p}_{1}$ ，恰好有 7 件一等品的概率记为 ${p}_{2}$ ，请直接写出 ${p}_{1}$ 与 ${p}_{2}$ 的大小关系.
        \end{enumerate}

        \begin{jieda}
            \begin{enumerate}
                \item 设 $A =$ “若从这批产品中任取一件，其为一等品或二等品”，$P\left( A\right)  = \dfrac{{10} + 6}{20} = \dfrac{16}{20} = \dfrac{4}{5}$
                
                \item 设 $A{ = }^{\prime \prime }$ 在抽取的 20 件产品中，随机抽取 3 件，既有一等品又有二等品”，$P\left( B\right)  = \dfrac{{10} \times  6 \times  4 + {{C}}_{10}^{2} \times  6 + {10} \times  {{C}}_{6}^{2}}{{{C}}_{20}^{3}} = \dfrac{11}{19}$
                
                \item ${p}_{1} < {p}_{2}$ ,因为 ${p}_{1} = {{C}}_{10}^{5}{\left( \dfrac{3}{5}\right) }^{5}{\left( \dfrac{2}{5}\right) }^{5},{p}_{2} = {{C}}_{10}^{7}{\left( \dfrac{3}{5}\right) }^{7}{\left( \dfrac{2}{5}\right) }^{3}$，所以 $\dfrac{{p}_{1}}{{p}_{2}} = \dfrac{42}{45} < 1$ .
            \end{enumerate}
        \end{jieda}
        
        \item 一次课外活动举行篮球投篮趣味比赛，选手在连续投篮时，第一次投进得 1 分, 并规定: 若某次投进, 则下一次投进的得分是本次得分的两倍; 若某次未投进, 则该次得 0 分, 且下一次投进得 1 分. 已知某同学连续投篮 $n$ 次,总得分为 $X$ ,每次投进的概率为 $\dfrac{1}{3}$ ,且每次投篮相互独立.
        \begin{enumerate}
            \item 当 $n = {30}$ 时，判断 $E\left( X\right)$ 与 10 的大小，并说明理由;
            
            \item 当 $n = 3$ 时，求 $X$ 的概率分布列和数学期望；
            
            \item 记 $X = 3$ 的概率为 ${P}_{n}\left( {n \geq  2,n \in  {N}^{ * }}\right)$ ,求 ${P}_{n}$ 的表达式.
        \end{enumerate}
        \begin{jieda}
            \begin{enumerate}
                \item $E\left( X\right)  > {10}$ . 理由如下:
                
                记该同学投篮 30 次投进次数为 $\xi$ ,则 $\xi  \sim  B\left( {{30},\dfrac{1}{3}}\right)$ .
                
                若每次投进都得 1 分,则得分的期望为 $E\left( \xi \right)  = {30} \times  \dfrac{1}{3} = {10}$ .
                
                由题中比赛规则可知连续投进时, 得分翻倍,
                
                故实际总得分的期望 $E\left( X\right)$ 必大于每次都得 1 分的数学期望.
                
                所以 $E\left( X\right)  > {10}$ .
                
                \item $X$ 的可能取值为0,1,2,3,7，其中
                
                $P\left( {X = 0}\right)  = {\left( \dfrac{2}{3}\right) }^{3} = \dfrac{8}{27}P\left( {X = 1}\right)  = {{C}}_{3}^{1} \times  \dfrac{1}{3} \times  {\left( \dfrac{2}{3}\right) }^{2} = \dfrac{12}{27}$
                
                $P\left( {X = 2}\right)  = {\left( \dfrac{1}{3}\right) }^{2} \times  \dfrac{2}{3} = \dfrac{2}{27}P\left( {X = 3}\right)  = {{C}}_{2}^{1} \times  {\left( \dfrac{1}{3}\right) }^{2} \times  \dfrac{2}{3} = \dfrac{4}{27}$
                
                $P\left( {X = 7}\right)  = {\left( \dfrac{1}{3}\right) }^{3} = \dfrac{1}{27}$
                
                所以 $X$ 的概率分布列为
                
                \begin{center}
                \begin{tabular}{|c|c|c|c|c|c|}
                \hline
                $X$ & 0 & 1 & 2 & 3 & 7 \\
                \cline{1-6}
                $P$ & $\dfrac{8}{27}$ & $\dfrac{12}{27}$ & $\dfrac{2}{27}$ & $\dfrac{4}{27}$ & $\dfrac{17}{27}$ \\
                \cline{1-6}
                \hline
                \end{tabular}
                
                \end{center}
                
                所以 $E\left( X\right)  = 0 \times  \dfrac{8}{27} + 1 \times  \dfrac{12}{27} + 2 \times  \dfrac{2}{27} + 3 \times  \dfrac{4}{27} + 7 \times  \dfrac{1}{27} = \dfrac{35}{27}$
                
                (3)投篮 $n$ 次得分为 3 分，有如下两种情形:
                
                情形一，恰好两次投进，且两次相邻；
                
                情形二，恰好三次投进，且任意两次都不相邻.
                
                当 $2 \leq  n \leq  4$ 时,情形二不可能发生,所以 ${P}_{n} = {{C}}_{n - 1}^{1}{\left( \dfrac{1}{3}\right) }^{2} \times  {\left( \dfrac{2}{3}\right) }^{n - 2} = \dfrac{n - 1}{4} \times  {\left( \dfrac{2}{3}\right) }^{n}$ ;
                
                当 $n \geq  5$ 时,情形一发生的概率为 ${{C}}_{n - 1}^{1}{\left( \dfrac{1}{3}\right) }^{2} \times  {\left( \dfrac{2}{3}\right) }^{n - 2} = \dfrac{n - 1}{4} \times  {\left( \dfrac{2}{3}\right) }^{n}$ ,
                
                情形二发生是指，将$(n - 3)$次未投进的投篮排成一列，共有$(n - 2)$个空位，选择其中 3 个空位作为投进的投篮,
                
                所以概率为 ${{C}}_{n - 2}^{3}{\left( \dfrac{1}{3}\right) }^{3} \times  {\left( \dfrac{2}{3}\right) }^{n - 3} = \dfrac{{n}^{3} - 9{n}^{2} + {26n} - {24}}{48} \times  {\left( \dfrac{2}{3}\right) }^{n}$ ,
                
                所以 ${P}_{n} = \dfrac{n - 1}{4} \times  {\left( \dfrac{2}{3}\right) }^{n} + \dfrac{{n}^{3} - 9{n}^{2} + {26n} - {24}}{48} \times  {\left( \dfrac{2}{3}\right) }^{n} = \dfrac{{n}^{3} - 9{n}^{2} + {38n} - {36}}{48} \times  {\left( \dfrac{2}{3}\right) }^{n}$ .
                
                综上所述， ${P}_{n} = \left\{  \begin{matrix} \dfrac{n - 1}{4} \times  {\left( \dfrac{2}{3}\right) }^{n},n = 2,3,4 \\  \dfrac{{n}^{3} - 9{n}^{2} + {38n} - {36}}{48} \times  {\left( \dfrac{2}{3}\right) }^{n},n \geq  5,n \in  {\mathbf{N}}^{ * } \end{matrix}\right.$ .
            \end{enumerate}
        \end{jieda}
    \end{enumerate}
    \item 超几何分布的判断
    \begin{enumerate}
        \item 下列随机事件中的随机变量 $X$ 服从超几何分布的是\blkda{\ding{172}\ding{174}}
        \begin{dingautolist}{172}
            \item 从 7 男 3 女共 10 名学生干部中选出 5 名学生干部，记选出女生的人数为 $X$
            \item 将一枚硬币连抛 3 次，正面向上的次数 $X$
            \item 从一批含有 13 件正品、2 件次品的产品中,不放回地任取 3 件，取得的次品数为 $X$
            \item 某射手的命中率为 0.8 ，现对目标射击 1 次，记命中目标的次数为 $X$
            \item 盒中有 4 个白球和 3 个黑球,每次从中摸出 1 球且放回, $X$ 是首次摸出黑球时的总次数
        \end{dingautolist}
        \begin{jieda}
            \begin{dingautolist}{172}
                \item 超几何分布
                \item 二项分布
                \item 超几何分布
                \item 两点分布
                \item 几何分布
            \end{dingautolist}
        \end{jieda}
        
        \item 一个袋中有 6 个同样大小的黑球，编号为 1, 2, 3, 4, 5, 6, 还有 4 个同样大小的白球, 编号为 7, 8, 9, 10.现从中任取 4 个球, 有如下几种变量:
        \begin{dingautolist}{172}
            \item $X$ 表示取出的最大号码
            \item $X$ 表示取出的最小号码
            \item $X$ 表示取出的白球个数
            \item 取出一个黑球记 2 分，取出一个白球记 1 分， $X$ 表示取出的 4 个球的总得分减去 4 的差
        \end{dingautolist}
        这四种变量中服从超几何分布的是\blkda{\ding{174}\ding{175}}
    \end{enumerate}
    \item 求超几何分布的概率
    \begin{enumerate}
        \item 数学老师从 6 道题中随机抽 3 道让同学检测, 规定至少要解答正确 2 道题才能及格. 某同学只能正确求解其中的 4 道题, 则该同学能及格的概率为\blkda{$\dfrac{4}{5}$}
        \begin{jieda}
            由题意知抽取 3 道题该同学不及格的情况只有: 只对一道题一种情况，则只答对一道题的概率为 $P = \dfrac{C_{4}^{1}C_{2}^{2}}{C_{6}^{3}} = \dfrac{1}{5}$ ，所以该同学及格的概率为 $\dfrac{4}{5}$
        \end{jieda}
        \item 设袋中有 8 个红球，4 个白球，若从袋中任取 4 个球，则其中至多 3 个红球的概率为\blkda{$1 - \dfrac{{{C}}_{8}^{4}}{{{C}}_{12}^{4}}$}
        \begin{jieda}
            从袋中任取 4 个球,其中红球的个数 $X$ 服从参数为 $N = {12},M = 8,n = 4$ 的超几何分布, 故至多有 3 个红球的概率为 $P\left( {X \leq  3}\right)  = 1 - P\left( {X = 4}\right)  = 1 - \dfrac{C_{8}^{4}}{C_{12}^{4}}$ .
        \end{jieda}
        \item 某学习小组共 12 人，其中有 5 名是“三好学生”，现从该小组中任选 5 人参加竞赛，用 $\xi$ 表示这 5 人中“三好学生”的人数，则下列概率中等于 $\dfrac{{{C}}_{7}^{5}{{C}}_{5}^{0} + {{C}}_{5}^{1}{{C}}_{7}^{4}}{{{C}}_{12}^{5}}$ 的是\dotfill{(\kh{B})}
        \xx{$P\left( {\xi  = 1}\right)$}{$P\left( {\xi  \leq  1}\right)$}{$P\left( {\xi  \geq  1}\right)$}{$P\left( {\xi  \leq  2}\right)$}
        \begin{jieda}
            由题意可得$P\left( {\xi  = 0}\right)  = \dfrac{C_{7}^5C_{5}^{0}}{C_{12}^5},P\left( {\xi  = 1}\right)  = \dfrac{C_{7}^{4}C_{5}^{1}}{C_{12}^5},P\left( {\xi  = 2}\right)  = \dfrac{C_{7}^{3}C_{5}^{2}}{C_{12}^5}$ ,
            
            $P\left( {\xi  \leq  1}\right)  = P\left( {\xi  = 0}\right)  + P\left( {\xi  = 1}\right)  = \dfrac{{{C}}_{7}^{5}{{C}}_{5}^{6} + {{C}}_{5}^{7}{{C}}_{7}^{4}}{{{C}}_{12}^{5}}$ ,
            
            $P\left( {\xi  \geq  1}\right)  = P\left( {\xi  = 2}\right)  + P\left( {\xi  = 1}\right)  = \dfrac{{{C}}_{7}^{3}{{C}}_{5}^{2} + {{C}}_{5}^{7}{{C}}_{7}^{4}}{{{C}}_{12}^{5}},P\left( {\xi  \leq  2}\right)  = 1$
        \end{jieda}
        \item 一个班级共有 30 名学生，其中有 10 名女生， 现从中任选三人代表班级参加学校开展的某项活动,假设选出的 3 名代表中的女生人数为变量 $X$ ,男生的人数为变量 $Y$ ，则 $P\left( {X = 2}\right)  + P\left( {Y = 2}\right)$ 等于\blkda{$\dfrac{{{C}}_{10}^{2}{{C}}_{20}^1 + {{C}}_{10}^1{{C}}_{20}^{2}}{{{C}}_{30}^{3}}$}
        \begin{jieda}
            因为 $P\left( {X = 2}\right)  = \dfrac{{{C}}_{10}^{2}{{C}}_{20}^{1}}{{{C}}_{30}^{3}},P\left( {Y = 2}\right)  = \dfrac{{{C}}_{10}^{1}{{C}}_{20}^{2}}{{{C}}_{30}^{3}}$，所以 $P\left( {X = 2}\right)  + P\left( {Y = 2}\right)  = \dfrac{{{C}}_{10}^{2}{{C}}_{20}^{1} + {{C}}_{10}^{1}{{C}}_{20}^{2}}{{{C}}_{30}^{3}}$ .
        \end{jieda}
    \end{enumerate}
    \item 超几何分布的均值与方差
    \begin{enumerate}
        \item 某地盛行糕点有 $n$ 种,该地的糕点店从中准备了 $m\left( {m < n}\right)$ 种糕点供顾客选购. 已知某顾客喜好的糕点有 $k\left( {k < n}\right)$ 种，则当其随机进入一家糕点店时，会发现该店中有若干种糕点符合其喜好. 记随机变量 $X$ 为该顾客发现符合其喜好的糕点的种数,则 $E\left( X\right)  =$\blkda{$\dfrac{km}{n}$}
        \item 2024 年 5 月中国邮政发行了《巢湖》特种邮票 3 枚，巢湖是继《太湖》(5 枚)、《鄱阳湖》(3 枚)、《洞庭湖》(4 枚)后，第四个登上特种邮票的五大淡水湖. 现从 15 枚邮票中随机抽取 2 枚，记抽取邮票《巢湖》的枚数为 $X$ ，则 $E\left( X\right)  =$\blkda{$\dfrac{2}{5}$}
        \begin{jieda}
            $E[X] = 2 \times \dfrac{3}{15} =  \dfrac{2}{5}$
        \end{jieda}
        \item 2024 年“与辉同行”直播间开播，董宇辉领衔 7 位主播从“心” 出发，其中男性 5 人，女性 3 人，现需排班晚 8:00 黄金档，随机抽取两人，则男生人数的期望为\blkda{$\dfrac{5}{4}$}
        
        \item 一个袋子中装有 6 个红球和 4 个白球，假设每个球被摸到的可能性是相等的. 从袋子中摸出 2 个球,其中白球的个数为 $X$ ,则 $X$ 的数学期望是
        \blkda{$\dfrac{4}{5}$}
    \item 已知 6 件产品中有 2 件次品，4 件正品，检验员从中随机抽取 3 件进行检测，记取到的正品数为 $X$ ，则下列结论正确的是\dotfill{(\kh{D})}
    \xx{$E\left( {{2X} - 1}\right)  = \dfrac{4}{3}$}{$D\left( X\right)  = \dfrac{1}{5}$}{$E\left( X\right)  = 1$}{$D\left( {{2X} - 1}\right)  = \dfrac{8}{5}$}
    \begin{jieda}
        根据题意可知, $X$ 可能取1,2,3,且服从超几何分布，故$E\left( X\right)  = 3 \times \dfrac{4}{6} = 2$
        
        $D\left( X\right)  = 3 \times \dfrac{4}{6} \times \left(1-\dfrac{4}{6}\right)\times \dfrac{6-3}{6-1} = \dfrac{2}{5}.$
        
        $E\left( {{2X} - 1}\right)  = {2E}\left( X\right)  - 1 = 2 \times  2 - 1 = 3$，$D\left( {{2X} - 1}\right)  = {4D}\left( X\right)  = \dfrac{8}{5}$
    \end{jieda}
    
    \item 口袋中有相同的黑色小球 $n$ 个，红、白、蓝色的小球各一个，从中任取 4 个小球. $\xi$ 表示当 $n = 3$ 时取出黑球的数目， $\eta$ 表示当 $n = 4$ 时取出黑球的数目. 则下列结论成立的是\dotfill{(\kh{A})}
    \xx{$E\left( \xi \right)  < E\left( \eta \right) ,D\left( \xi \right)  < D\left( \eta \right)$}{$E\left( \xi \right)  > E\left( \eta \right) ,D\left( \xi \right)  < D\left( \eta \right)$}{$E\left( \xi \right)  < E\left( \eta \right) ,D\left( \xi \right)  > D\left( \eta \right) $}{$E\left( \xi \right)  > E\left( \eta \right) ,D\left( \xi \right)  > D\left( \eta \right)$}
    \begin{jieda}
        $E\left( \xi \right)  = 2,D\left( \xi \right) =\dfrac{2}{5}$，$E\left( \eta \right)  = \dfrac{16}{7}, D\left( \eta \right)  = \dfrac{24}{49}$
    \end{jieda}

    
        \item 已知 20 条试题中有 8 条选择题, 甲无放回地依次从中抽取 5 条题, 乙有放回地依次从中抽取 5 条题,甲、乙每次均抽取一条试题,抽出的 5 条题中选择题的条数分别为 ${\xi }_{1},{\xi }_{2}$ , ${\xi }_{1},{\xi }_{2}$ 的期望分别为 $E\left( {\xi }_{1}\right) ,E\left( {\xi }_{2}\right)$ ,方差分别为 $D\left( {\xi }_{1}\right) ,D\left( {\xi }_{2}\right)$ ,则\dotfill{(\kh{A})}
        \xx{$E\left( {\xi }_{1}\right)  = E\left( {\xi }_{2}\right) ,D\left( {\xi }_{1}\right)  < D\left( {\xi }_{2}\right)$}{$E\left( {\xi }_{1}\right)  = E\left( {\xi }_{2}\right) ,D\left( {\xi }_{1}\right)  > D\left( {\xi }_{2}\right)$}{$E\left( {\xi }_{1}\right)  < E\left( {\xi }_{2}\right) ,D\left( {\xi }_{1}\right)  < D\left( {\xi }_{2}\right)$}{$E\left( {\xi }_{1}\right)  < E\left( {\xi }_{2}\right) ,D\left( {\xi }_{1}\right)  > D\left( {\xi }_{2}\right)$}
        \begin{jieda}
            由题意可知, ${\xi }_{1}$ 的可能取值为 $0,1,2,3,4,5,{\xi }_{2}$ 的可能取值为0,1,2,3,4,5,
            
            随机变量${\xi }_{1}$ 服从超几何分布，随机变量 ${\xi }_{2}$ 服从二项分布，
            
            根据超几何分布的均值方差公式得:
            
            $n = 5,N = {20},M = 8$ ,即 $E\left( {\xi }_{1}\right)  = \dfrac{nM}{N} = \dfrac{5 \times  8}{20} = 2$ ,
            
            $D\left( {\xi }_{1}\right)  = \dfrac{nM}{N}\left( {1 - \dfrac{M}{N}}\right) \dfrac{N - n}{N - 1} = \dfrac{5 \times  8}{20} \times  \left( {1 - \dfrac{8}{20}}\right)  \times  \dfrac{{20} - 5}{{20} - 1} = 2 \times  \dfrac{3}{5} \times  \dfrac{15}{19} = \dfrac{18}{19}$
            
            根据超二项分布的均值方差公式得:
            
            $n = 5,p = \dfrac{8}{20} = \dfrac{2}{5}$ ,即 $E\left( {\xi }_{2}\right)  = {np} = 5 \times  \dfrac{2}{5} = 2$
            
            $D\left( {\xi }_{2}\right)  = {np}\left( {1 - p}\right)  = 5 \times  \dfrac{2}{5} \times  \dfrac{3}{5} = \dfrac{6}{5}$
            
            所以 $E\left( {\xi }_{1}\right)  = E\left( {\xi }_{2}\right) ,D\left( {\xi }_{1}\right)  < D\left( {\xi }_{2}\right)$ .
        \end{jieda}
        
        \item 一个不透明的袋子有 10 个除颜色不同外，大小、质地完全相同的球, 其中有 6 个黑球, 4 个白球. 现进行如下两个试验, 试验一: 逐个不放回地随机摸出 3 个球, 记取到白球的个数为 ${X}_{1}$ ,期望方差分别为 $E\left( {X}_{1}\right) ,D\left( {X}_{1}\right)$ ; 试验二: 逐个有放回地随机摸出 3 个球,记取到白球的个数为 ${X}_{2}$ ,期望和方差分别为 $E\left( {X}_{2}\right) ,D\left( {X}_{2}\right)$ ,则下列判断正确的是\dotfill{(\kh{A})}
        \xx{$E\left( {X}_{1}\right)  = E\left( {X}_{2}\right) ,D\left( {X}_{1}\right)  < D\left( {X}_{2}\right)$}{$E\left( {X}_{1}\right)  = E\left( {X}_{2}\right) ,D\left( {X}_{1}\right)  > D\left( {X}_{2}\right)$}{$E\left( {X}_{1}\right)  > E\left( {X}_{2}\right) ,D\left( {X}_{1}\right)  > D\left( {X}_{2}\right)$}{$E\left( {X}_{1}\right)  < E\left( {X}_{2}\right) ,D\left( {X}_{1}\right)  < D\left( {X}_{2}\right)$}
    
    \item 为营造浓厚的全国文明城市创建氛围，积极响应创建全国文明城市号召，提高对创城行动的责任感和参与度，学校号召师生利用周末参与创城志愿活动.高二(1)班某小组有男生 4 人，女生 2 人，现从中随机选取 2 人作为志愿者参加活动.
    \begin{enumerate}
        \item 求在有女生参加活动的条件下，恰有一名女生参加活动的概率
        \item 记参加活动的女生人数为 $X$ ，求 $X$ 的分布列及期望 $E\left( X\right)$ 、方差 $D\left( X\right)$
    \end{enumerate}
    \begin{jieda}
        \begin{enumerate}
            \item 设“有女生参加活动”为事件 $A$ ，“恰有一名女生参加活动”为事件 $B$ .
            
            则 $P\left( {AB}\right)  = \dfrac{{{C}}_{4}^{1}{{C}}_{2}^{1}}{{{C}}_{6}^{2}} = \dfrac{8}{15},P\left( A\right)  = \dfrac{{{C}}_{4}^{1}{{C}}_{2}^{1} + {{C}}_{2}^{2}}{{{C}}_{6}^{2}} = \dfrac{3}{5}$ ，所以 $P\left( {B \mid  A}\right)  = \dfrac{P\left( {AB}\right) }{P\left( A\right) } = \dfrac{\dfrac{8}{15}}{\dfrac{3}{5}} = \dfrac{8}{9}$ 
            \item 依题意知 $X$ 服从超几何分布，且 $P\left( {X = k}\right)  = \dfrac{C_{2}^{k}C_{4}^{2 - k}}{C_{6}^{2}}\left( {k = 0,1,2}\right)$ ，
            
            $P\left( {X = 0}\right)  = \dfrac{{{C}}_{4}^{2}}{{{C}}_{6}^{2}} = \dfrac{2}{5},P\left( {X = 1}\right)  = \dfrac{{{C}}_{4}^{1} \cdot  {{C}}_{2}^{1}}{{{C}}_{6}^{2}} = \dfrac{8}{15},P\left( {X = 2}\right)  = \dfrac{{{C}}_{2}^{2}}{{{C}}_{6}^{2}} = \dfrac{1}{15},$
            
            所以 $X$ 的分布列为:
            
            \begin{center}
            \begin{tabular}{|c|c|c|c|}
            \hline
            $X$ & 0 & 1 & 2 \\
            \cline{1-4}
            $P$ & $\dfrac{2}{5}$ & $\dfrac{8}{15}$ & $\dfrac{1}{15}$ \\
            \cline{1-4}
            \hline
            \end{tabular}
            \end{center}
            
            $E\left( X\right)  = 0 \times  \dfrac{2}{5} + 1 \times  \dfrac{8}{15} + 2 \times  \dfrac{1}{15} = \dfrac{2}{3},D\left( X\right)  = \dfrac{2}{5} \times  {\left( 0 - \dfrac{2}{3}\right) }^{2} + \dfrac{8}{15} \times  {\left( 1 - \dfrac{2}{3}\right) }^{2} + \dfrac{1}{15} \times  {\left( 2 - \dfrac{2}{3}\right) }^{2} = \dfrac{16}{45}$
        \end{enumerate}
    \end{jieda}
    \item 某大学的武术协会有 10 名同学，成员构成如下表所示. 表中部分数据不清楚,只知道从这 10 名同学中随机抽取 1 名同学,该名同学的专业为数学的概率为 $\dfrac{2}{5}$ .
    
    \begin{center}
    \begin{tabular}{|c|c|c|c|c|}
    \hline
    性别 & 中文 & 数学 & 英语 & 体育 \\
    \cline{1-5}
    男 & $a$ & $b$ & 1 & 1 \\
    \cline{1-5}
    女 & 1 & 1 & 1 & 1 \\
    \cline{1-5}
    \hline
    \end{tabular}
    \end{center}
    
    现从这 10 名同学中随机选取 3 名同学参加该市的武术比赛(每名同学被选到的可能性相等).
    \begin{enumerate}
        \item 求$a$、$b$的值
        
        \item 求选出的 3 名同学恰为专业互不相同的男生的概率
        
        \item 设 $Y$ 为选出的 3 名同学中是女生或专业为数学的人数，求随机变量 $Y$ 的分布列、均值及方差.
    \end{enumerate}
    \begin{jieda}
        \begin{enumerate}
            \item 由题意得 $\dfrac{b + 1}{10} = \dfrac{2}{5}$ 解得 $b = 3$ .
            
            由 $a + 3 + 1 + 1 + 1 + 1 + 1 + 1 = {10}$ ,得解得 $a = 1$ .
            
            \item 所求的概率为 $P = \dfrac{{{C}}_{3}^{1}{{C}}_{3}^{2} + {{C}}_{3}^{3}}{{{C}}_{10}^{3}} = \dfrac{9 + 1}{120} = \dfrac{1}{12}$ .
            
            \item 由已知，这 10 名同学中是女生或者专业为数学的人数为 $7,Y$ 的可能取值为0,1,2,3.
            
            $P\left( {Y = 0}\right)  = \dfrac{C_{3}^{3}}{C_{10}^{3}} = \dfrac{1}{120},P\left( {Y = 1}\right)  = \dfrac{C_{7}^{1}C_{3}^{2}}{C_{10}^{3}} = \dfrac{21}{120} = \dfrac{7}{40}.$
            
            $P\left( {Y = 2}\right)  = \dfrac{C_{7}^{2}C_{3}^{1}}{C_{10}^{3}} = \dfrac{63}{120} = \dfrac{21}{40},P\left( {Y = 3}\right)  = \dfrac{C_{7}^{3}}{C_{10}^{3}} = \dfrac{35}{120} = \dfrac{7}{24}$
            
            所以 $Y$ 的分布列为
            
            \begin{center}
            \begin{tabular}{|c|c|c|c|c|}
            \hline
            $Y$ & 0 & 1 & 2 & 3 \\
            \cline{1-5}
            $P$ & $\dfrac{1}{120}$ & $\dfrac{7}{40}$ & $\dfrac{21}{40}$ & $\dfrac{7}{24}$ \\
            \cline{1-5}
            \hline
            \end{tabular}
            \end{center}
            
            均值为 $E\left( Y\right)  = 0 \times  \dfrac{1}{120} + 1 \times  \dfrac{7}{40} + 2 \times  \dfrac{21}{40} + 3 \times  \dfrac{7}{24} = \dfrac{21}{10}$ ,
            
            方差为 $D\left( Y\right)  = {\left( 0 - \dfrac{21}{10}\right) }^{2} \times  \dfrac{1}{120} + {\left( 1 - \dfrac{21}{10}\right) }^{2} \times  \dfrac{7}{40} + {\left( 2 - \dfrac{21}{10}\right) }^{2} \times  \dfrac{21}{40} + {\left( 3 - \dfrac{21}{10}\right) }^{2} \times  \dfrac{7}{24} = \dfrac{49}{100}$
        \end{enumerate}
    \end{jieda}
    \end{enumerate}
    \item 二项分布与超几何分布的综合应用
    \begin{enumerate}
    \item 某校设计了一个实验学科的实验考查方案；考生从 6 道备选题中一次性随机抽取 3 题，按照题目要求独立完成全部实验操作，规定:至少正确完成两题便可通过，已知 6 道备选题中甲生有 4 题能正确完成，2 题不能完成；考生乙每题正确完成的概率都是 $\dfrac{2}{3}$ ，求:
    \begin{enumerate}
        \item 分别写出甲、乙两考生正确完成题数的概率分布列，并计算数学期望
        \item 试用统计知识分析比较两考生的实验操作能力
    \end{enumerate}
    \begin{jieda}
        \begin{enumerate}
            \item 设考生甲正确完成实验操作的题数为 $\xi$ ，则 $\xi$ 的可能取值是 1，2，3，

            $P\left( {\xi  = 1}\right)  = \dfrac{{{c}}_{4}^{7}{{c}}_{2}^{2}}{{{c}}_{6}^{3}} = \dfrac{1}{5},P\left( {\xi  = 2}\right)  = \dfrac{{{c}}_{4}^{2}{{c}}_{2}^{1}}{{{c}}_{6}^{3}} = \dfrac{3}{5},P\left( {\xi  = 3}\right)  = \dfrac{{{c}}_{4}^{3}{{c}}_{2}^{0}}{{{c}}_{6}^{3}} = \dfrac{1}{5},$
            
            所以 $\xi$ 的分布列为:
            
            \begin{center}
            \begin{tabular}{|c|c|c|c|}
            \hline
            $\xi$ & 1 & 2 & 2 \\
            \cline{1-4}
            $P$ & $\dfrac{1}{5}$ & $\dfrac{3}{5}$ & $\dfrac{1}{5}$ \\
            \cline{1-4}
            \hline
            \end{tabular}
            \end{center}
            
            则 $E\left( \xi \right)  = 1 \times  \dfrac{1}{5} + 2 \times  \dfrac{3}{5} + 3 \times  \dfrac{1}{5} = 2$ ;
            
            设考生乙正确完成实验操作的题数为 $\eta$ ，易知 $\eta  \sim  B\left( {3,\dfrac{2}{3}}\right)$ ，
            
            所以 $P\left( {\eta  = 0}\right)  = {{C}}_{3}^{0}{\left( 1 - \dfrac{2}{3}\right) }^{3} = \dfrac{1}{27},P\left( {\eta  = 1}\right)  = {{C}}_{3}^{1}{\left( \dfrac{2}{3}\right) }^{1}{\left( 1 - \dfrac{2}{3}\right) }^{2} = \dfrac{2}{9}$ ,
            
            $P\left( {\eta  = 2}\right)  = {{C}}_{3}^{2}{\left( \dfrac{2}{3}\right) }^{2}{\left( 1 - \dfrac{2}{3}\right) }^{1} = \dfrac{4}{9},P\left( {\eta  = 3}\right)  = {{C}}_{3}^{3}{\left( \dfrac{2}{3}\right) }^{3} = \dfrac{8}{27}$
            
            所以 $\eta$ 的分布列为:
            
            \begin{center}
            \begin{tabular}{|c|c|c|c|c|}
            \hline
            $\eta$ & 0 & 1 & 2 & 3 \\
            \cline{1-5}
            $P$ & $\dfrac{1}{27}$ & $\dfrac{2}{9}$ & $\dfrac{4}{9}$ & $\dfrac{8}{27}$ \\
            \cline{1-5}
            \hline
            \end{tabular}
            \end{center}
            
            所以 $E\left( \eta \right)  = 3 \times  \dfrac{2}{3} = 2$ .
        \end{enumerate}
        \item 由(1)知 $E\left( \xi \right)  = E\left( \eta \right)  = 2$ ，
        
        $D\left( \xi \right)  = {\left( 1 - 2\right) }^{2} \times  \dfrac{1}{5} + {\left( 2 - 2\right) }^{2} \times  \dfrac{3}{5} + {\left( 3 - 2\right) }^{2} \times  \dfrac{1}{5} = \dfrac{2}{5}$
        
        $D\left( \eta \right)  = 3 \times  \dfrac{2}{3} \times  \left( {1 - \dfrac{2}{3}}\right)  = \dfrac{2}{3}$
        
        $P\left( {\xi  \geq  2}\right)  = \dfrac{3}{5} + \dfrac{1}{5} = \dfrac{4}{5},P\left( {\eta  \geq  2}\right)  = \dfrac{4}{9} + \dfrac{8}{27} = \dfrac{20}{27}$
        
        所以 $D\left( \xi \right)  < D\left( \eta \right) ,P\left( {\xi  \geq  2}\right)  > P\left( {\eta  \geq  2}\right)$ ,
        
        故从正确完成实验操作的题数的均值方面分析, 两人水平相当;
        
        从正确完成实验操作的题数的方差方面分析, 甲的水平更稳定;
        
        从至少正确完成 2 题的概率方面分析, 甲通过的可能性更大.
    \end{jieda}
    \item 某市为了解党员对党史知识的掌握情况，在该市随机选取 100 名党员进行调查，其中不同年龄段的人数分布如下表:
    
    \begin{center}
    \begin{tabular}{|c|c|c|c|c|c|c|}
    \hline
    年龄 & 小于 30 & [30,40) & [40,50) & [50,60) & [60,70) & 大于等于 70 \\
    \cline{1-7}
    人数 & 5 & 10 & 25 & 35 & 15 & 10 \\
    \cline{1-7}
    \hline
    \end{tabular}
    \end{center}
    \begin{enumerate}
        \item 已知年龄在 $\lbrack {30},{40})$ 的女性党员有 3 人，现从该年龄段这 10 人中随机选择 3 人进行座谈，用 $X$ 表示这 3 人中女性党员的人数，求 $X$ 的分布列和数学期望
        \item 若用样本的频率近似代替概率，现从该市党员中随机抽取 20 名党员进行调查，用 $Y$ 表示其中年龄在 $\lbrack {30},{50})$ 内的人数，求当 $P\left( {Y = k}\right) \left( {k = 0,1,\cdots ,{20}}\right)$ 取最大值时 $k$ 的值
    \end{enumerate}
    \begin{jieda}
        \begin{enumerate}
            \item 由题意得 $X$ 服从超几何分布,其中 $n = 3,N = {10},M = 3$ ,
            
            所以 $P\left( {X = 0}\right)  = \dfrac{{{C}}_{7}^{3}}{{{C}}_{10}^{3}} = \dfrac{7}{24},P\left( {X = 1}\right)  = \dfrac{{{C}}_{7}^{2}{{C}}_{3}^{1}}{{{C}}_{10}^{3}} = \dfrac{21}{40}$ ,
            
            $P\left( {X = 2}\right)  = \dfrac{{{C}}_{7}^{1}{{C}}_{3}^{2}}{{{C}}_{10}^{3}} = \dfrac{7}{40},P\left( {X = 1}\right)  = \dfrac{{{C}}_{3}^{3} = }{{{C}}_{10}^{3}} = \dfrac{1}{120},$
            
            所以 $X$ 的分布列为
            
            \begin{center}
            \begin{tabular}{|c|c|c|c|c|}
            \hline
            $X$ & 0 & 1 & 2 & 3 \\
            \cline{1-5}
            $P$ & $\dfrac{4}{27}$ & $\dfrac{21}{40}$ & $\dfrac{7}{40}$ & $\dfrac{1}{120}$ \\
            \cline{1-5}
            \hline
            \end{tabular}
            \end{center}
            
            所以 $E\left( X\right)  = 0 \times  \dfrac{7}{24} + 1 \times  \dfrac{21}{40} + 2 \times  \dfrac{7}{40} + 3 \times  \dfrac{1}{120} = {0.9}$

            \item 由题意可得从该市党员中随机抽取 1 名党员，其年龄在 $[30,50)$ 的概率为 $\dfrac{{10} + {25}}{100} = {0.35}$
            
            则 $Y \sim  B\left( {{20},{0.35}}\right)$ ,即 $P\left( {Y = k}\right)  = {{C}}_{20}^{k} \times  {0.35}^{k} \times  {0.65}^{{20} - k}\left( {k = 0,1,2,\cdots ,{20}}\right)$ ,
            
            所以 $P\left( {Y = k + 1}\right)  = {{C}}_{20}^{k + 1} \times  {0.35}^{k + 1} \times  {0.65}^{{19} - k}$ ,
            
            令 $\dfrac{P\left( {Y = k - 1}\right) }{P\left( {Y = k}\right) } = \dfrac{{{C}}_{20}^{k - 1} \times  {0.35}^{k - 1} \times  {0.65}^{{21} - k}}{{{C}}_{20}^{k} \times  {0.35}^{k} \times  {0.65}^{{20} - k}} = \dfrac{k}{{21} - k} \times  \dfrac{13}{7} \leq  1$ ，得 $k \leq  \dfrac{147}{20}$ ，
            
            令 $\dfrac{P\left( {Y = k + 1}\right) }{P\left( {Y = k}\right) } = \dfrac{{{C}}_{20}^{k + 1} \times  {0.35}^{k + 1} \times  {0.65}^{{19} - k}}{{{C}}_{20}^{k} \times  {0.35}^{k} \times  {0.65}^{{20} - k}} = \dfrac{{20} - k}{k + 1} \times  \dfrac{7}{13} \leq  1$ ，得 $k \geq  \dfrac{127}{20}$ ，
            
            所以 $\dfrac{127}{20} \leq  k \leq  \dfrac{147}{20}$ ,
            
            因为 $k \in  {Z}$ ,所以当 $k = 7$ 时, $P\left( {Y = k}\right)$ 最大.
        \end{enumerate}
    \end{jieda}
    \end{enumerate}
    \item 正态密度函数
    \begin{enumerate}
        \item 设随机变量 $X$ 的正态分布密度函数为 $f\left( x\right)  = \frac{1}{2\sqrt{\pi }} \cdot  {{e}}^{-\frac{{\left( x + 3\right) }^{2}}{4}},x \in  \left( {-\infty , + \infty }\right)$ ，则参数 $\mu$ ， $\sigma$ 的值分别是\dotfill{(\kh{D})}
        \xx{$\mu  = 3,\sigma  = 2$}{$\mu  =  - 3,\sigma  = 2$}{$\mu  = 3,\sigma  = \sqrt{2}$}{$\mu  =  - 3,\sigma  = \sqrt{2}$}
        \begin{jieda}
            由正态分布密度函数表达式知 $\mu  =  - 3,\sigma  = \sqrt{2}$.
        \end{jieda}
        
        \item 设随机变量 $X$ 服从正态分布,且相应的概率密度函数为 $f\left( x\right)  =  \dfrac{1}{\sqrt{6\pi }}{{e}}^{-\frac{{x}^{2} - {4x} + 4}{6}}$ ,则\dotfill{(\kh{C})}
        \xx{$\mu  = 2,\sigma  = 3$}{$\mu  = 3,\sigma  = 2$}{$\mu  = 2,\sigma  = \sqrt{3}$}{$\mu  = 3,\sigma  = \sqrt{3}$}
        \begin{jieda}
            由正态分布密度函数 $f\left( x\right)  = \dfrac{1}{\sqrt{6\pi }}{{e}}^{-\frac{{x}^{2} - {4x} + 4}{6}} = \dfrac{1}{\sqrt{2\pi } \cdot  \sqrt{3}}{{e}}^{-\frac{{\left( x - 2\right) }^{2}}{2 \times  {\left( \sqrt{3}\right) }^{2}}}$ ,可得 $\mu  = 2,\sigma  = \sqrt{3}$.
        \end{jieda}

        \item 已知随机变量 $X$ 的正态密度函数为 $f\left( x\right)  = \frac{1}{\sqrt{8\pi }}{{e}}^{-\frac{{x}^{2}}{8}}\left( {x \in  R}\right)$ ,则其均值和标准差分别是\blkda{0 和 2}
        \begin{jieda}
            由 $f\left( x\right)$ 的形式,知 $X$ 的均值和标准差分别为 0 和 2.
        \end{jieda}
        
        \item 给出下列函数: 
        \begin{dingautolist}{172}
            \item $f\left( x\right)  = \dfrac{1}{\sqrt{2\pi }\sigma }{{e}}^{-\frac{{\left( x + \mu \right) }^{2}}{2{\sigma }^{2}}}$
            \item $f\left( x\right)  = \dfrac{1}{\sqrt{2\pi }}{{e}}^{-\frac{{\left( x - \mu \right) }^{2}}{4}}$
            \item $f\left( x\right)  = \dfrac{1}{\sqrt{2} \cdot  \sqrt{2\pi }}{{e}}^{-\frac{{x}^{2}}{4}}$
            \item $f\left( x\right)  = \dfrac{1}{\sqrt{\pi }}{{e}}^{-{\left( x - \mu \right) }^{2}}$
        \end{dingautolist}
        其中 $\mu  \in  \left( {-\infty , + \infty }\right)$ ， $\sigma  > 0$，则可以作为正态分布密度函数的个数有\blkda{\ding{172}\ding{174}\ding{175}}
    \end{enumerate}
    \item 正态曲线的特点
    \begin{enumerate}
        \item 已知三个随机变量的正态密度函数 ${f}_{i}\left( x\right) \left( {x \in  R,i = 1,2,3}\right)$ 的图象如图所示，则\dotfill{(\kh{D})}
        
        \begin{center}
        \includegraphics[width=0.3\textwidth]{bodyxb/src/bo_d67fmj77aajc739ac550_1_143_1245_384_281_0.jpg}
        \end{center}

        \xx{${\mu }_{1} < {\mu }_{2} = {\mu }_{3},{\sigma }_{1} = {\sigma }_{2} > {\sigma }_{3}$}{${\mu }_{1} > {\mu }_{2} = {\mu }_{3},{\sigma }_{1} = {\sigma }_{2} < {\sigma }_{3}$}{${\mu }_{1} = {\mu }_{2} < {\mu }_{3},{\sigma }_{1} < {\sigma }_{2} = {\sigma }_{3}$}{${\mu }_{1} < {\mu }_{2} = {\mu }_{3},{\sigma }_{1} = {\sigma }_{2} < {\sigma }_{3}$}
        \begin{jieda}
            因为正态密度函数 ${f}_{2}\left( x\right)$ 和 ${f}_{3}\left( x\right)$ 的图象关于同一条直线对称,所以 ${\mu }_{2} = {\mu }_{3}$ .
            
            又 ${f}_{2}\left( x\right)$ 的图象的对称轴在 ${f}_{1}\left( x\right)$ 的图象的对称轴的右边,所以 ${\mu }_{1} < {\mu }_{2} = {\mu }_{3}$ .
            
            因为 $\sigma$ 越大,曲线越“矮胖”. $\sigma$ 越小,曲线越“瘦高”,
            
            由图可知，正态密度函数 ${f}_{1}\left( x\right)$ 和 ${f}_{2}\left( x\right)$ 的图象一样 “瘦高”， ${f}_{3}\left( x\right)$ 的图象明显 “矮胖”.
            
            所以 ${\sigma }_{1} = {\sigma }_{2} < {\sigma }_{3}$ .
        \end{jieda}

        
        \item 已知正态分布 $N\left( {1,{\sigma }^{2}}\right)$ 的正态密度曲线如图所示, $X \sim  N\left( {1,{\sigma }^{2}}\right)$ , 则下列选项中，不能表示图中阴影部分面积的是 
        \dotfill{(\kh{C})}
        
        \begin{center}
        \includegraphics[width=0.2\textwidth]{bodyxb/src/bo_d67fmj77aajc739ac550_2_142_1184_345_186_0.jpg}
        \end{center}
        \xx{$\dfrac{1}{2} - P\left( {X \leq  0}\right)$}{$\dfrac{1}{2} - P\left( {X \geq  2}\right)$}{$\dfrac{1}{2} - P\left( {1 \leq  X \leq  2}\right)$}{$\dfrac{1}{2}P\left( {X \leq  2}\right)  - \dfrac{1}{2}P\left( {X \leq  0}\right)$}
        \begin{jieda}
            正态分布 $N\left( {1,{\sigma }^{2}}\right)$ 的正态密度曲线关于直线 $x = 1$ 对称，可得图中阴影部分可表示为 $P\left( {0 \leq  X \leq  1}\right)  = P\left( {X \leq  1}\right)  - P\left( {X \leq  0}\right)  = \dfrac{1}{2} - P\left( {X \leq  0}\right)  = \dfrac{1}{2} - P\left( {X \geq  2}\right)$ , 故选项 A, B 正确;
            
            对 ${C}$ : 由对称性可得 $\dfrac{1}{2} - P\left( {1 \leq  X \leq  2}\right)  = P\left( {X \geq  2}\right)  = P\left( {X \leq  0}\right)$ ,故选项 ${C}$ 错误;
            
            对 ${D}$ : 由对称性可得 $P\left( {0 \leq  X \leq  1}\right)  = P\left( {1 \leq  X \leq  2}\right)$ ,
            
            所以图中阴影部分面积可表示为 $P\left( {0 \leq  X \leq  1}\right)  = \dfrac{1}{2}\left\lbrack  {P\left( {X \leq  2}\right)  - P\left( {X \leq  0}\right) }\right\rbrack$ ,故选项 ${D}$ 正确.
        \end{jieda}
        
        \item 某市高中数学统考中，甲、乙、丙三所学校的数学成绩分别服从正态分布 $N\left( {{\mu }_{1},{\sigma }_{1}}\right)$，$N\left( {{\mu }_{2},{\sigma }_{2}}\right) $，$N\left( {{\mu }_{3},{\sigma }_{3}}\right)$ ，其正态分布的密度曲线如图所示，则\dotfill{(\kh{D})}
        
        \begin{center}
        \includegraphics[width=0.3\textwidth]{bodyxb/src/bo_d67fmj77aajc739ac550_2_142_407_424_280_0.jpg}
        \end{center}

        \xx{${\mu }_{1} = {\mu }_{2} > {\mu }_{3}$}{${\mu }_{1} = {\mu }_{2} < {\mu }_{3}$}{${\mu }_{1} > {\mu }_{2} = {\mu }_{3}$}{${\mu }_{1} < {\mu }_{2} = {\mu }_{3}$}
        \begin{jieda}
            观察曲线知, ${\mu }_{1} < {\mu }_{2} = {\mu }_{3}$ 
        \end{jieda}
        
        \item 设随机变量 $X \sim  N\left( {{\mu }_{1},{\sigma }_{1}^{2}}\right) ,Y \sim  N\left( {{\mu }_{2},{\sigma }_{2}^{2}}\right)$ ,这两个正态分布密度曲线如图所示, 则\dotfill{(\kh{D})}
        
        \begin{center}
        \includegraphics[width=0.3\textwidth]{bodyxb/src/bo_d67fmj77aajc739ac550_3_143_407_413_284_0.jpg}
        \end{center}
        \xx{${\mu }_{1} > {\mu }_{2}$}{${\sigma }_{1} < {\sigma }_{2}$}{$P\left( {X \leq  {\mu }_{1}}\right)  > P\left( {Y \geq  {\mu }_{2}}\right)$}{$P\left( {X \leq  {\mu }_{2}}\right)  > P\left( {Y \leq  {\mu }_{1}}\right)$}
        \begin{jieda}
            $X$ 的密度曲线的对称轴在 $Y$ 的密度曲线的对称轴的左边,即 ${\mu }_{1} < {\mu }_{2}$ .
            
            $X$ 的密度曲线较为分散, $Y$ 的密度曲线较为集中,即 ${\sigma }_{1} > {\sigma }_{2}$ ,故 ${{AB}}$ 错误;
            
            因为 $P\left( {X \leq  {\mu }_{1}}\right)  = {0.5},P\left( {Y \geq  {\mu }_{2}}\right)  = {0.5}$ ,所以 ${C}$ 错误;
            
            因为 $P\left( {x \leq  {\mu }_{2}}\right)  > {0.5},P\left( {Y \leq  {\mu }_{1}}\right)  < {0.5}$ ,所以 ${D}$ 正确;
        \end{jieda}
    \end{enumerate}
    \item 利用正态曲线的对称性求概率
    \begin{enumerate}
        \item 若随机变量 $X \sim  N\left( {{60},{\sigma }^{2}}\right)$ ,且 $P\left( {X \leq  {40}}\right)  = {0.2}$ ,则 $P\left( {X \leq  {80}}\right)  =$\blkda{$0.8$}
        \begin{jieda}
            因为 $P\left( {X \leq  {40}}\right)  = {0.2}$ ,所以 $P\left( {X \geq  {80}}\right)  = {0.2}$ ,所以 $P\left( {X \leq  {80}}\right)  = 1 - {0.2} = {0.8}$ .
        \end{jieda}
        \item 已知随机变量 $X$ 服从正态分布 $N\left( {3,{\sigma }^{2}}\right) ,P\left( {X \leq  4}\right)  = {0.84}$ , 则 $P\left( {2 < X \leq  4}\right)  =$\blkda{$0.68$}
        \begin{jieda}
            由题意得 $P\left( {3 < X \leq  4}\right)  = {0.84} - {0.5} = {0.34}$，由正态曲线的对称性知 $P\left( {2 < X \leq  3}\right)  = P\left( {3 < X \leq  4}\right)  = {0.34}$ ,
            
            所以 $P\left( {2 < X \leq  4}\right)  = {0.34} \times  2 = {0.68}$ .
        \end{jieda}
        
        \item 若随机变量 $X \sim  N\left( {6,1}\right)$ ，且 $P\left( {5 < X \leq  7}\right)  = a,P\left( {4 < X \leq  8}\right)  = b$ ,则 $P\left( {4 < X \leq  7}\right)$ 等于\blkda{$\dfrac{b + a}{2}$}
        \begin{jieda}
            随机变量 $X \sim  N\left( {6,1}\right)$ ,且 $P\left( {5 < X \leq  7}\right)  = a,P\left( {4 < X \leq  8}\right)  = b$
            
            由正态密度曲线的对称性可知, $P\left( {4 < X \leq  5}\right)  = \dfrac{b - a}{2}$ ,
            
            所以 $P\left( {4 < X \leq  7}\right)  = \dfrac{b - a}{2} + a = \dfrac{b + a}{2}$ .
        \end{jieda}
        \item 已知随机变量 $X - N\left( {\mu ,{\sigma }^{2}}\right)$ ,若其对应的正态密度函数 $f\left( x\right)$ 满足 $f\left( {2 - x}\right)  = f\left( x\right)$ ,且 $P\left( {X \leq  0}\right)  = {0.1}$ ,则 $P\left( {1 \leq  X \leq  2}\right)  =$\blkda{$0.4$}
        \begin{jieda}
            由 $f\left( {2 - x}\right)  = f\left( x\right)$ ,则正态密度函数 $f\left( x\right)$ 关于 $x = 1$ 对称,即 $\mu  = 1$，则 $P\left( {1 \leq  X \leq  2}\right)  = P\left( {0 \leq  X \leq  1}\right)  = {0.5} - P\left( {X \leq  0}\right)  = {0.5} - {0.1} = {0.4}$ .
        \end{jieda}
        \item 设随机变量 $X$ 服从标准正态分布 $X \sim  N\left( {0,1}\right)$ ,那么对于任意 $a$ ,记 $\Phi \left( a\right)  =  P\left( {X < a}\right)$ ,已知 $\Phi \left( a\right)  = {0.7}$ ,则 $P\left( {\left| X\right|  < a}\right)  = $\blkda{$0.4$} .
        \begin{jieda}
            $P\left( {\left| X\right|  < a}\right)  = P\left( {-a < X < a}\right)  = 1 - {2P}\left( {X > a}\right)  = 1 - 2\left\lbrack  {1 - \Phi \left( a\right) }\right\rbrack = 1 - 2 \times  \left( {1 - {0.7}}\right)  = {0.4}$ .
        \end{jieda}
    \end{enumerate}

    \item 根据正态曲线的对称性求参数
    \begin{enumerate}
        \item 已知随机变量 $X$ 服从正态分布 $N\left( {1,{\sigma }^{2}}\right)$ ,若 $P\left( {X \leq  a}\right)  = {0.3}$ ,且 $P(a \leq  X \leq  a + 2) = {0.4}$ ,则 $a =$ \blkda{0}
        \begin{jieda}
            由题意知随机变量 $X$ 服从正态分布 $N\left( {1,{\sigma }^{2}}\right) ,P\left( {X \leq  a}\right)  = {0.3}$
            
            如图所示,结合 $P\left( {a \leq  X \leq  a + 2}\right)  = {0.4}$ ,得 $P\left( {X \geq  a + 2}\right)  = {0.3}$ ,
            
            \begin{center}
            \includegraphics[width=0.3\textwidth]{bodyxb/src/bo_d67fmj77aajc739ac550_4_141_1870_375_267_0.jpg}
            \end{center}
            
            
            可知 $a,a + 2$ 关于 $x = 1$ 对称,所以 $a + a + 2 = 2 \times  1$ ,解得 $a = 0$
        \end{jieda}
        
        \item 设随机变量 $X \sim  N\left( {1,{5}^{2}}\right)$ ,且 $P\left( {X \leq  0}\right)  = P\left( {X \geq  a - 2}\right)$ ,则实数 $a$ 的值为\blkda{4}
        \begin{jieda}
            因为 $X \sim  N\left( {1,{5}^{2}}\right)$ ,所以正态分布密度曲线的对称轴为 $x = 1$，因为 $P\left( {X \leq  0}\right)  = P\left( {X \geq  a - 2}\right)$ ,所以 $\dfrac{0 + a - 2}{2} = 1 \Longrightarrow  a = 4$ .
        \end{jieda}
        
        \item 已知随机变量 $X$ 服从正态分布 $N\left( {4,{\sigma }^{2}}\right)$ ,若 $P(X \leq  1 +  {2a}) + P\left( {X \leq  1 - a}\right)  = 1$ ,则 $a =$\blkda{6}
        \begin{jieda}
            因为 $P\left( {X \leq  1 + {2a}}\right)  + P\left( {X \leq  1 - a}\right)  = 1$ ,所以$P\left( {X \leq  1 + {2a}}\right)  = 1 - P\left( {X \leq  1 - a}\right)  = P\left( {X > 1 - a}\right)$ ,因为随机变量服从正态分布 $X \sim  N\left( {4,{\sigma }^{2}}\right)$ ,所以 $1 + {2a} + 1 - a = 2 \times  4$ ,解得: $a = 6$ .
        \end{jieda}
        \item 已知 $\xi  \sim  N\left( {2,{5}^{2}}\right)$ ,且 $P\left( {\xi  \leq  1}\right)  = P\left( {\xi  \geq  a + 1}\right)$ ,则 $\dfrac{1}{x} + \dfrac{4}{a - x}(0 <  x < a)$ 的最小值为\blkda{$\dfrac{9}{2}$}
        \begin{jieda}
            因为 $\xi  \sim  N\left( {2,{5}^{2}}\right)$ ,且 $P\left( {\xi  \leq  t}\right)  = P\left( {\xi  \geq  a + t}\right)$ ,则 $\dfrac{1 + a + t}{2} = 2$ ,解得 $a = 2$
            
            因为 $0 < x < 2$ ,则 $0 < 2 - x < 2$，所以 $\dfrac{1}{x} + \dfrac{4}{2 - x} = \dfrac{1}{2}\left\lbrack  {x + \left( {2 - x}\right) }\right\rbrack  \left( {\dfrac{1}{x} + \dfrac{4}{2 - x}}\right)  = \dfrac{1}{2}\left( {5 + \dfrac{2 - x}{x} + \dfrac{4x}{2 - x}}\right) \geq  \dfrac{1}{2}\left( {5 + 2\sqrt{\dfrac{2 - x}{x} \cdot  \dfrac{4x}{2 - x}}}\right)  = \dfrac{9}{2}$
            
            当且仅当 $\left\{  \begin{array}{l} \dfrac{2 - x}{x} = \dfrac{4x}{2 - x} \\  0 < x < 2 \end{array}\right.$ 时,即当 $x = \dfrac{2}{3}$ 时,等号成立
        \end{jieda}
    \end{enumerate}
    \item 标准正态分布与${3\sigma }$ 原则
    \begin{enumerate}
        \item 随机变量 $\xi$ 服从正态分布 $N\left( {1,{\sigma }^{2}}\right)$ ,随机变量 $\eta$ 服从标准正态分布$N\left( {0,1}\right)$ ,若 $P\left( {\eta  < 1}\right) = a$ ,则 $P\left( {1 < \xi  < 1 + \sigma }\right)  =$\blkda{$a - \dfrac{1}{2}$} . (用字母 $a$ 表示)
        \begin{jieda}
            随机变量 $\eta$ 服从标准正态分布 $N\left( {0,1}\right)$ ,根据对称性可知 $P\left( {\eta  < 0}\right)  = \dfrac{1}{2}$
            
            因为 $P\left( {\eta  < 1}\right)  = a$ ,所以 $P\left( {0 < \eta  < 1}\right)  = a - \dfrac{1}{2}$ ,即 $P\left( {\mu  < \eta  < \mu  + \sigma }\right)  = a - \dfrac{1}{2}$ ,
            
            随机变量 $\xi$ 服从正态分布 $N\left( {1,{\sigma }^{2}}\right)$，故 $P\left( {1 < \xi  < 1 + \sigma }\right)  = a - \dfrac{1}{2}$
        \end{jieda}
        
        \item 现实世界中的很多随机变量服从正态分布，例如反复测量某一个物理量,其测量误差 $X$ 通常被认为服从正态分布. 若某物理量做 $n$ 次测量,测量结果的误差 $X \sim  N\left( {0,\dfrac{2}{n}}\right)$ , 要控制 $\left| X\right|  \geq  \dfrac{1}{2}$ 的概率不大于 0.0027，至少要测量的次数为\blkda{$72$}(参考数据: $P\left( {\mu  - {3\sigma } \leq  X \leq  \mu  + {3\sigma }}\right)  = {0.9973}$ )
        \begin{jieda}
            因为 $X \sim  N\left( {0,\dfrac{2}{n}}\right)$ ,所以 $\mu  = 0,\sigma  = \sqrt{\frac{2}{n}}$，根据题意得 $P\left( {\left| X\right|  \geq  \dfrac{1}{2}}\right)  \leq  {0.0027}$ ,则 $P\left( {\left| X\right|  < \dfrac{1}{2}}\right)  \geq  1 - {0.0027} = {0.9973}$，即 $P\left( {-\dfrac{1}{2} < X < \dfrac{1}{2}}\right)  \geq  {0.9973}$，因为 $\mu  = 0$ ，所以 $P\left( {-{3\sigma } \leq  X \leq  {3\sigma }}\right)  = {0.9973}$，所以 ${3\sigma } \leq  \dfrac{1}{2}$ 所以 $\sqrt{\dfrac{2}{n}} \leq  \dfrac{1}{6}$ 解得 $n \geq  {72}$，所以至少要测量的次数为 72 次
        \end{jieda}
        \item 学校有 1000 名学生生参加了“希望杯”数学竞赛，此次竞赛成绩 $X$ 服从正态分布 $N\left( {{75},{100}}\right)$ ，估计竞赛成绩在 65 分到 85 分之间的人数约为\blkda{$683$}人.(参考数据 $P\left( {\mu  - \sigma  < X < \mu  + \sigma }\right)  = {0.6826},P\left( {\mu  - {2\sigma } < X < \mu  + {2\sigma }}\right)  = {0.9544},P(\mu  - {3\sigma } < X <  \mu  + {3\sigma }) = {0.9974})$
        \begin{jieda}
            由于竞赛成绩 $X$ 服从正态分布 $N\left( {{75},{100}}\right)$，所以 $\mu  = {75},\sigma  = {10}$，所以 $P\left( {{65} < X < {85}}\right)  = P\left( {{75} - {10} < X < {75} + {10}}\right)  = {0.6826}$，故该校 1000 名学生竞赛成绩在 65 分到 85 分之间的人数约为: $1000 \times 0.6826=682.6 \approx 683$
        \end{jieda}
        \item 某校高二级学生参加期末调研考试的数学成绩 $X$ 服从正态分布 $N\left( {{95},{8}^{2}}\right)$ ,将考试成绩从高到低,按照 16\%,34\%,34\%,16\%的比例分为 $A,B,C,D$ 四个等级. 若小明的数学成绩为 105 分, 则属于等级 \blkda{A}(附: $P\left( {\mu  - \sigma  < X < \mu  + \sigma }\right)  \approx  {0.68},P\left( {\mu  - {2\sigma } < X < \mu  + {2\sigma }}\right)  \approx  {0.95},P\left( {\mu  - {3\sigma } < X < \mu  + {3\sigma }}\right)  \approx$ 0.99)
        \begin{jieda}
            数学测试成绩服从正态分布 $N\left( {{95},{8}^{2}}\right)$ ,则 $\mu  = {95},\sigma  = 8$，由于 $A,D$ 等级的概率之和为 ${16}\%  + {16}\%  = {32}\%  = 1 - P\left( {\mu  - \sigma  < X < \mu  + \sigma }\right)$
            
            所以 $P\left( {X < \mu  - \sigma }\right)  = P\left( {X > \mu  + \sigma }\right)  = \dfrac{1 - P\left( {\left| {X - \mu }\right|  \leq  \sigma }\right) }{2} = {0.16}$
            
            $P\left( {X < {87}}\right)  = P\left( {X > {103}}\right)  = \dfrac{1 - P\left( {\left| {X - \mu }\right|  \leq  \sigma }\right) }{2} = {0.16}$
            
            ,而 $P\left( {\mu  - \sigma  < X < \mu }\right)  = P\left( {\mu  < X < \mu  + \sigma }\right)  = {0.34}$ ,即 $P\left( {{87} < X < {95}}\right)  = P\left( {{95} < X < {103}}\right)  = {0.34}$ , 故 $X > {103}$ 为 $A$ 等级, ${95} < X < {103}$ 为 $B$ 等级, ${87} < X < {95}$ 为 $C$ 等级, $X < {87}$ 为 $D$ 等级, 故 105 分为 $A$ 等级.
        \end{jieda}
        
        \item 某次高二质量抽测中,学生的数学成绩 $X$ 服从正态分布 $N$ (96,144). 已知参加本次考试的学生约有 10000 人，如果小明在这次考试中数学成绩为 120 分，则小明的数学成绩在本次抽测的名次大约是第\blkda{$228$}名. 附:若 $X \sim  N\left( {\mu ,{\sigma }^{2}}\right)$ ,则 $P\left( {\mu  - \sigma  < X < \mu  + \sigma }\right)  = {0.6827},P\left( {\mu  - {2\sigma } < X < \mu  + {2\sigma }}\right)  = {0.9545}$
        \begin{jieda}
            由 $X \sim  N\left( {{96},{144}}\right) ,\mu  + {2\sigma } = {96} + {24} = {120},\mu  - {2\sigma } = {96} - {24} = {72}$，则 $P\left( {{72} < X < {120}}\right)  = {0.9545}$ ,故 $P\left( {X \geq  {120}}\right)  = \dfrac{1 - {0.9545}}{2} = {0.02275}$，${10000} \times  {0.02275} = {227.5} \approx  {228}$ ，故小明的数学成绩在本次抽测的名次大约是第 228 名.
        \end{jieda}
    \end{enumerate}
    \item 正态分布的应用
    \begin{enumerate}
        \item 小明上学有时坐公交车，有时骑自行车，他记录多次数据，分析得到: 坐公交车平均用时 30min，样本方差为 36；骑自行车平均用时 34min，样本方差为 4，假设坐公交车用时 $X \sim  N\left( {{30},{6}^{2}}\right)$ ,骑自行车用时 $Y \sim  N\left( {{34},{2}^{2}}\right)$ ,则\dotfill{(\kh{B})}
        \xx{$P\left( {X \leq  {38}}\right)  > P\left( {Y \leq  {38}}\right)$}{$P\left( {X \leq  {34}}\right)  > P\left( {Y \leq  {34}}\right)$}{如果有 38 分钟可用，小明应选择坐公交车}{如果有 34 分钟可用，小明应选择自行车}
        \begin{jieda}
            因为 $X \sim  N\left( {{30},{6}^{2}}\right) ,Y \sim  N\left( {{34},{2}^{2}}\right)$，将 $X,Y$ 化为标准正态分布,则 $\dfrac{X - {30}}{6} \sim  N\left( {0,1}\right) ,\dfrac{Y - {34}}{2} \sim  N\left( {0,1}\right)$
            
            因为 $\dfrac{38 - {30}}{6} < \dfrac{38 - {34}}{2}$ ,所以 $P\left( {X \leq  {38}}\right)  < P\left( {Y \leq  {38}}\right)$ ,故 A 错误;
            
            又 $P\left( {X \leq  {34}}\right)  > \dfrac{1}{2},P\left( {Y \leq  {34}}\right)  = \dfrac{1}{2},P\left( {X \leq  {34}}\right)  > P\left( {Y \leq  {34}}\right)$ ,故 B 正确;
            
            因为 $P\left( {X \leq  {38}}\right)  < P\left( {Y \leq  {38}}\right)$ ,所以如果有 38 分钟可用,小明应选择自行车,故 ${C}$ 错误;
            
            因为 $P\left( {X \leq  {34}}\right)  > P\left( {Y \leq  {34}}\right)$ ,所以如果有 34 分钟可用,小明应选择坐公交车,故 ${D}$ 错误.
        \end{jieda}
        \item 某省高考改革试点方案规定:2023 年高考总成绩由语文、数学、 外语三门统考科目和思想政治、历史、地理、物理、化学、生物六门选考科目组成，将每门选考科目的考生原始成绩从高到低划分为 $A,B + ,B,C + ,C,D + ,D,E$ 共 8 个等级,参照正态分布原则,确定各等级人数所占比例分别为 $3\% ,7\% ,{16}\% ,{24}\% ,{24}\% ,{16}\% ,7\% ,3\%$ ,选考科目成绩计入考生总成绩时,将 $A$ 至 $E$ 等级内的考生原始成绩，依照等比例转换法则，分别转换到[91,100],[81,90],[71,80],[61,70],[51,60],[41,50],[31,40],[21,30]八个分数区间，得到考生的等级成绩，如果该省某次高考模拟考试物理科目的原始成绩 $X \sim  N$ (50,256)，那么 $B$ +等级的原始分最低大约为\blkda{$71$}(参考数据:  \ding{172}  若 $X \sim  N\left( {\mu ,{\sigma }^{2}}\right) ,Y = \dfrac{X - \mu }{\sigma }$ ,则 $Y \sim  N\left( {0,1}\right)$ ;  \ding{173}  当 $Y \sim  N\left( {0,1}\right)$ 时, $P\left( {Y \leq  {1.3}}\right)  \approx  {0.9}$ )
        \begin{jieda}
            由题意可得,将学生成绩从高到低排名,排在 $B$ +等级最低分后面的学生约为学生总数的 90\%. 因为原始成绩 $X \sim  N\left( {{50},{256}}\right)$ ,所以 $\mu  = {50},\sigma  = {16}$ .令 $Y = \dfrac{X - \mu }{\sigma }$ ,则 $Y \sim  N\left( {0,1}\right)$ ; 又当 $Y \sim  N\left( {0,1}\right)$ 时, $P\left( {Y \leq  {1.3}}\right)  \approx  {0.9}$，所以 $\dfrac{X - {50}}{16} \leq  {1.3}$ ，解得 $X \leq  {70.8}$ ，所以 $B +$ 等级的原始分最低大约为 71 .
        \end{jieda}
        \item 某校高二年级下学期期末考试数学试卷满分为 150 分，90 分以上 (含 90 分) 为及格. 阅卷结果显示，全年级 1200 名学生的数学成绩近似服从正态分布，试卷的难度系数 (难度系数 $=$ 平均分/满分) 为 0.49,标准差为 22,则该次数学考试及格的人数约为\blkda{272}人附: 若 $X \sim  N\left( {\mu ,{\sigma }^{2}}\right)$ , 记 $p\left( k\right)  = P\left( {\mu  - {k\sigma } \leq  X \leq  \mu  + {k\sigma }}\right)$ ,则 $p\left( {0.75}\right)  \approx  {0.547},p\left( 1\right)  \approx  {0.683}$ .
        \begin{jieda}
            由题得 $\mu  = {0.49} \times  {150} = {73.5},\sigma  = {22}$ 
            
            $\because p\left( k\right)  = P\left( {\mu  - {k\sigma } \leq  X \leq  \mu  + {k\sigma }}\right) ,p\left( {0.75}\right)  \approx  {0.547}$
            
            $\therefore P\left( {{57} \leq  X \leq  {90}}\right) = p\left( {0.75}\right)  \approx  {0.547}$ ,
            
            $P\left( {X \geq  {90}}\right)  = {0.5} \times  \left( {1 - {0.547}}\right)  = {0.2265},$
            
            $\therefore$ 该校及格人数为 ${0.2265} \times  {1200} \approx  {272}$ (人)
        \end{jieda}
        \item 在某次数学考试中,考生的成绩 $X$ 服从正态分布 $X \sim  N$ (95,225).
        \begin{enumerate}
            \item 试求考试成绩 $X$ 位于区间 $\left\lbrack  {{65},{125}}\right\rbrack$ 内的概率
            \item 若这次考试共有 3000 名考生，试估计考试成绩位于区间 $[80,110]$ 内的考生人数
        \end{enumerate}
        (参考数据: $P\left( {\mu  - \sigma  \leq  X \leq  \mu  + \sigma }\right)  \approx  {0.6827},P\left( {\mu  - {2\sigma } \leq  X \leq  \mu  + {2\sigma }}\right)  \approx  {0.9545}$ )
        
        \begin{jieda}
            \begin{enumerate}
                \item $\because X \sim  N\left( {{95},{225}}\right)$
                
                $\therefore \mu  = {95},\sigma  = {15}$ .
                
                $\because \mu  - {2\sigma } = {95} - 2 \times  {15} = {65}$ ,
                
                $\mu  + {2\sigma } = {95} + 2 \times  {15} = {125}.$
                
                且 $P\left( {\mu  - {2\sigma } \leq  X \leq  \mu  + {2\sigma }}\right)  \approx  {0.9545}$ ,
                
                $\therefore P\left( {{65} \leq  X \leq  {125}}\right)  \approx  {0.9545}$ .
                
                \item $\because \mu  - \sigma  = {95} - {15} = {80}$ ,
                
                $\mu  + \sigma  = {95} + {15} = {110},$
                
                且 $P\left( {\mu  - \sigma  \leq  X \leq  \mu  + \sigma }\right)  \approx  {0.6827}$ ,
                
                $\therefore P\left( {{80} \leq  X \leq  {110}}\right)  \approx  {0.6827}$ ,
                
                $\therefore$ 考试成绩位于区间 $\left\lbrack  {{80},{110}}\right\rbrack$ 内的考生人数为 ${3000} \times  {0.6827} \approx  {2048}$ (人).
            \end{enumerate}
        \end{jieda}
        
        \item 辽宁、广东、河北、湖北、湖南、江苏、福建、重庆等八省市从 2021 年起全部采用 "3 + 1 + 2 "的新高考模式. "3 "指的是语文、数学、外语，这三门科目考试参加统一高考，由教育部考试中心统一命题，以原始成绩计入考生总成绩； “1”指的是物理和历史中的一科，考生必须从物理和历史两个科目中选择一科，由各省自主命题，以原始成绩计入考生总成绩. 为了让考生更好的适应新高考模式，某省几个地市进行了统一的高考适应性考试. 在所有人考考生中有 30000 人选考物理，考后物理成绩 $X$ (满分 100 分) 服从正态分布 $N\left( {{55},{10}^{2}}\right)$ .
        \begin{enumerate}
            \item 分别估计成绩在 $\left\lbrack  {{45},{65}}\right\rbrack$ 和 75 分以上者的人数；(运算过程中精确到 0.0001 ，最后结果保留为整数)

            附 1: $P\left( {\mu  - \sigma  \leq  X \leq  \mu  + \sigma }\right)  \approx  {0.6827},P\left( {\mu  - {2\sigma } \leq  X \leq  \mu  + {2\sigma }}\right)  \approx  {0.9545},P(\mu  - {3\sigma } \leq  X \leq  \mu  +  {3\sigma }) \approx  {0.9973}$
            \item 本次考试物理成绩 $X$ 服从正态分布 $N\left( {\mu ,{\sigma }^{2}}\right)$ . 令 $\eta  = \dfrac{X - \mu }{\sigma }$ ，则 $\eta  \sim  N\left( {0,1}\right)$ ，若本次考试物理成绩的前$25\%$划定为优秀等级，试估计物理优秀等级划线分大约为多少分？(结果保留为整数)
            附 2: 若 $\eta  \sim  N\left( {0,1}\right)$ ,则 $P\left( {\eta  < {0.8}}\right)  \approx  {0.75}$ .
        \end{enumerate}
        \begin{jieda}
            \begin{enumerate}
                \item 因为 $X \sim  N\left( {{55},{10}^{2}}\right)$ ,所以 $\mu  = {55},\sigma  = {10}$ ,
                
                又 $P\left( {\mu  - \sigma  \leq  X \leq  \mu  + \sigma }\right)  \approx  {0.6827}$ ,所以 $P\left( {{45} \leq  X \leq  {65}}\right)  \approx  {0.6827}$ ,
                
                所以成绩在 $\left\lbrack  {{45},{65}}\right\rbrack$ 的人数约为 ${30000} \times  {0.6827} = {20481}$ 人,
                
                由正态分布曲线的对称性可得: $P\left( {{35} \leq  X \leq  {75}}\right)  = P\left( {\mu  - {2\sigma } \leq  X \leq  \mu  + {2\sigma }}\right)  \approx  {0.9545}$ ,
                
                则 $P\left( {X > {75}}\right)  = \dfrac{1 - P\left( {{35} \leq  X \leq  {75}}\right) }{2} \approx  {0.0228}$ ,
                
                所以估计 75 分以上的人数约为 ${30000} \times  {0.0228} = {684}$ 人;
                
                \item 设该划线分为 $m$ ，由 $X \sim  N\left( {{55},{10}^{2}}\right)$ ，得 $\mu  = {55},\sigma  = {10}$ ，
                
                令 $\eta  = \dfrac{X - \mu }{\sigma } = \dfrac{X - {55}}{10} \geq  \dfrac{m - {55}}{10}$ ,
                
                由题意因为 $\eta  \sim  N\left( {0,1}\right) ,P\left( {\eta  < {0.8}}\right)  \approx  {0.75}$ ,
                
                所以 $P\left( {\eta  \geq  {0.8}}\right)  \approx  {0.25}$ ,所以 $\dfrac{m - {55}}{10} \approx  {0.8}$ ,
                
                所以 $m \approx  {63}$ .
            \end{enumerate}
        \end{jieda}
    \end{enumerate}
    \item 正态分布与其他知识综合
    \begin{enumerate}
        \item 为加大自然生态系统和环境保护力度，加强企业对尊重自然、顺应自然、保护自然的生态文明理念，某市对化工企业的排污情况进行调查，并出台相应的整治措施. 相关部门对 1000 家化工企业所排污水的质量及周围空气质量进行了综合检测, 得分情况如频率分布直方图所示.
        
        \begin{center}
        \includegraphics[width=0.3\textwidth]{bodyxb/src/bo_d67fmj77aajc739ac550_12_139_1619_452_371_0.jpg}
        \end{center}
        \begin{enumerate}
            \item 计算该市化工企业的平均得分 $\bar{x}$ (同一组中的数据以这组数据的中间值为代表)
            \item 已知化工企业的得分情况 $X$ 近似服从正态分布 $N\left( {\mu ,{\sigma }^{2}}\right)$ ，其中 $\mu  = \bar{x},{\sigma }^{2} = {16}^{2}$ ，则得分在 [67,83] 内的企业大约有多少家
            \item 按照(2)中概率分布随机抽取 100 家化工企业，分数不低于 19 分的企业有多少家时概率最大.
        \end{enumerate}
        参考数据: 若随机变量 $X$ 服从正态分布 $N\left( {\mu ,{\sigma }^{2}}\right)$ ,则 $P\left( {\mu  - \sigma  < X < \mu  + \sigma }\right)  = {0.6826},P(\mu  - {2\sigma } < X <  \mu  + {2\sigma }) = {0.9544},P\left( {\mu  - {3\sigma } < X < \mu  + {3\sigma }}\right)  = {0.9974}.$
        \begin{jieda}
            \begin{enumerate}
                \item 该市被调查的化工企业的污染情况得分的平均值为 $x = \left( {{10} \times  {0.0050} + {30} \times  {0.0125} + {50} \times  }\right.  {0.0150} + {70} \times  {0.0100} + {90} \times  {0.0075}) \times  {20} = {51}.$

                \item 由(1)知化工企业的得分情况 $X \sim  N\left( {{51},{16}^{2}}\right)$ . 因为 ${67} = {51} + {16},{83} = {51} + {16} \times  2$，所以 $P\left( {{67} \leq  X \leq  {83}}\right)  = P\left( {0 \leq  X \leq  {83}}\right)  - P\left( {0 \leq  X < {67}}\right) = P\left( {0 \leq  X \leq  {51} + {16} \times  2}\right)  - P\left( {0 \leq  X < {51} + {16}}\right) = \dfrac{1}{2}\left\lbrack  {P\left( {\mu  - {2\sigma } \leq  X \leq  \mu  + {2\sigma }}\right)  - P\left( {\mu  - \sigma  \leq  X \leq  \mu  + \sigma }\right) }\right\rbrack   \approx  \dfrac{{0.9544} - {0.6826}}{2} = {0.1359}.$ 可得所求企业大约有 ${0.1359} \times  {1000} \approx  {136}$ 家.
                
                \item 由 (2) 得 $P\left( {X \geq  {19}}\right)  = P\left( {X \geq  \mu  - {2\sigma }}\right)  \approx  {0.5} + \dfrac{0.9544}{2} = {0.9772}$ ,
                
                所以每家企业得分不低于 19 分的概率为 0.9772 .
                
                则得分不低于 19 分的企业数 $\xi  \sim  B\left( {{100},{0.9772}}\right)$ .
                
                其中恰有 $k$ 家企业得分不低于 19 分的概率为 $P\left( {\xi  = k}\right)  = {{C}}_{100}^{k}{p}^{k}{\left( 1 - p\right) }^{{100} - k}$ ,
                
                令 $\dfrac{P\left( {\xi  = k}\right) }{P\left( {\xi  = k - 1}\right) } = \dfrac{\left( {{101} - k}\right) p}{\left( {1 - p}\right) k} > 1,\dfrac{P\left( {\xi  = k}\right) }{P\left( {\xi  = k + 1}\right) } = \dfrac{\left( {k + 1}\right) \left( {1 - p}\right) }{\left( {{100} - k}\right) p} > 1$ ,
                
                可得 ${101p} - 1 < k < {101p},{101p} = {98.6972}$ ,解得 $k = {98}$ ,
                
                故在走访的 100 家化工企业中, 分数不低于 19 分的企业有 98 家时概率最大.
            \end{enumerate}
        \end{jieda}
        \item 某试点高校校考过程中笔试通过后才能进入面试环节. 2022 年报考该试点高校的学生的笔试成绩 $X$ 近似服从正态分布 $N\left( {\mu ,{\sigma }^{2}}\right)$ . 其中， $\mu$ 近似为样本平均数， ${\sigma }^{2}$ 近似为样本方差 ${s}^{2}$ . 已知 $\mu$ 的近似值为 ${76.5},s$ 的近似值为 5.5,以样本估计总体.
        \begin{enumerate}
            \item 假设有 84.135\% 的学生的笔试成绩高于该校预期的平均成绩，求该校预期的平均成绩大约是多少？
            
            \item 若笔试成绩高于 76.5 进入面试，若从报考该试点高校的学生中随机抽取 10 人，设其中进入面试学生数为$\xi$ ,求随机变量 $\xi$ 的期望.
            
            \item 现有甲、乙、丙、丁四名学生进入了面试，且他们通过面试的概率分别为 $\dfrac{1}{3}\text{ 、 }\dfrac{1}{3}\text{ 、 }\dfrac{1}{2}\text{ 、 }\dfrac{1}{2}$ 设这 4 名学生中通过面试的人数为 $X$ ,求随机变量 $X$ 的分布列和数学期望.
        \end{enumerate}
        
        参考数据: 若 $X \sim  N\left( {\mu ,{\sigma }^{2}}\right)$，则: $P\left( {\mu  - \sigma  < X \leq  \mu  + \sigma }\right)  \approx  {0.6827}$，$P\left( {\mu  - {2\sigma } < X \leq  \mu  + {2\sigma }}\right)  \approx  {0.9545}$ ，$P\left( {\mu  - {3\sigma } < X \leq  \mu  + {3\sigma }}\right)  \approx  {0.9973}.$
        \begin{jieda}
            \begin{enumerate}
                \item 由 $P\left( {X > \mu  - \sigma }\right)  = \dfrac{1}{2} + \dfrac{P\left( {\mu  - \sigma  < X \leq  \mu  + \sigma }\right) }{2} \approx  {0.84135}$ ,又 $\mu$ 的近似值为 76.5, $\sigma  = s$ 的近似值为 5.5，所以该校预期的平均成绩大约是 ${76.5} - {5.5} = {71}$ (分)
                
                \item 由 $\mu  \approx  {76.5}$ ,可得 $P\left( {\xi  > {76.5}}\right)  = \dfrac{1}{2}$ ，
                
                即从所有报考该试点高校的学生中随机抽取 1 人，
                
                该学生笔试成绩高于 76.5 的概率为 $\dfrac{1}{2}$
                
                所以随机变量 $\xi$ 服从二项分布 $\xi  \sim  \left( {{10},\dfrac{1}{2}}\right)$ ,故 $E\left( \xi \right)  = {10} \times  \dfrac{1}{2} = 5$
                
                \item $X$ 的可能取值为 0,1,2,3,4,
                
                $P\left( {X = 0}\right)  = {{C}}_{2}^{0} \times  {\left( 1 - \dfrac{1}{3}\right) }^{2} \times  {{C}}_{2}^{0} \times  {\left( 1 - \dfrac{1}{2}\right) }^{2} = \dfrac{1}{9},$
                
                $P\left( {X = 1}\right)  = {{C}}_{2}^{1} \times  \dfrac{1}{3} \times  \left( {1 - \dfrac{1}{3}}\right)  \times  {{C}}_{2}^{0} \times  {\left( 1 - \dfrac{1}{2}\right) }^{2} + {{C}}_{2}^{0} \times  {\left( 1 - \dfrac{1}{3}\right) }^{2} \times  {{C}}_{2}^{1} \times  \dfrac{1}{2} \times  \left( {1 - \dfrac{1}{2}}\right)  = \dfrac{1}{3}$
                
                $P\left( {X = 2}\right)  = {{C}}_{2}^{2} \times  {\left( \dfrac{1}{3}\right) }^{2} \times  {{C}}_{2}^{0} \times  {\left( 1 - \dfrac{1}{2}\right) }^{2} + {{C}}_{2}^{1} \times  \dfrac{1}{3} \times  \left( {1 - \dfrac{1}{3}}\right)  \times  {{C}}_{2}^{1} \times  \dfrac{1}{2} \times  \left( {1 - \dfrac{1}{2}}\right)  + {{C}}_{2}^{0} \times  {\left( 1 - \dfrac{1}{3}\right) }^{2} \times  {{C}}_{2}^{2} \times  {\left( \dfrac{1}{2}\right) }^{2} = \dfrac{13}{36}$ ,
                
                $P\left( {X = 3}\right)  = {{C}}_{2}^{2} \times  {\left( \dfrac{1}{3}\right) }^{2} \times  {{C}}_{2}^{1} \times  \dfrac{1}{2} \times  \left( {1 - \dfrac{1}{2}}\right)  + {{C}}_{2}^{1} \times  \dfrac{1}{3} \times  \left( {1 - \dfrac{1}{3}}\right)  \times  {{C}}_{2}^{2} \times  {\left( \dfrac{1}{2}\right) }^{2} = \dfrac{1}{6}$ ,
                
                $P\left( {X = 4}\right)  = {{C}}_{2}^{2} \times  {\left( \dfrac{1}{3}\right) }^{2} \times  {{C}}_{2}^{2} \times  {\left( \dfrac{1}{2}\right) }^{2} = \dfrac{1}{36}$
                
                所以 $X$ 的分布列为
                
                \begin{center}
                \begin{tabular}{|c|c|c|c|c|c|}
                \hline
                $X$ & 0 & 1 & 2 & 3 & 4 \\
                \hline
                $P$ & $\dfrac{1}{9}$ & $\dfrac{1}{3}$ & $\dfrac{13}{36}$ & $\dfrac{1}{6}$ & $\dfrac{1}{3}$ \\
                \cline{1-6}
                \hline
                \end{tabular}
                \end{center}
                
                所以 $E\left( X\right)  = 0 \times  \dfrac{1}{9} + 1 \times  \dfrac{1}{3} + 2 \times  \dfrac{13}{36} + 3 \times  \dfrac{1}{6} + 4 \times  \dfrac{1}{36} = \dfrac{5}{3}$
            \end{enumerate}
        \end{jieda}
        \item 某小区开展了两会知识问答活动，现将该小区参与该活动的 240 位居民的得分(满分 100 分)进行了统计，得到如下的频率分布直方图.
        
        \begin{center}
        \includegraphics[width=0.3\textwidth]{bodyxb/src/bo_d67fmj77aajc739ac550_15_139_533_426_343_0.jpg}
        \end{center}
        \begin{enumerate}
            \item 若此次知识问答的得分 $X$ 服从 $N\left( {\mu ,{\mathcal{E}}^{2}}\right)$ ，其中 $\mu$ 近似为参与本次活动的 240 位居民的平均得分(同一组中的数据用该组区间的中点值代表)，求 $P\left( {{66} < X \leq  {90}}\right)$ 的值
            \item 本次活动，制定了如下奖励方案:参与本次活动得分低于 $\mu$ 的居民获得一次抽奖机会，参与本次活动得分不低于 $\mu$ 的居民获得两次抽奖机会，每位居民每次有 $\dfrac{2}{3}$ 的机会抽中一张 10 元的话费充值卡，有 $\dfrac{1}{3}$ 的机会抽中一张 20 元的话费充值卡，假设每次抽奖相互独立，假设该小区居民王先生参与本次活动，求王先生获得的话费充值卡的总金额 $Y$ 的概率分布列,并估计本次活动需要准备的话费充值卡的总金额.
        \end{enumerate}
        参考数据: $P\left( {\mu  - \sigma  < X \leq  \mu  + \sigma }\right)  \approx  {0.6827},P\left( {\mu  - {2\sigma } < X \leq  \mu  + {2\sigma }}\right)  \approx  {0.9545},P(\mu  - {3\sigma } < X \leq  \mu  + {3\sigma }) \approx  {0.9973}$ .
        \begin{jieda}
            \begin{enumerate}
                \item 依题意， $\mu  = {55} \times  {0.1} + {65} \times  {0.3} + {75} \times  {0.3} + {85} \times  {0.2} + {95} \times  {0.1} = {74}$ , 所以 $X \sim  N\left( {{74},{8}^{2}}\right)$
                
                故 $P\left( {{66} < X \leq  {90}}\right)  = P\left( {\mu  - \sigma  < X \leq  \mu  + {2\sigma }}\right)$
                
                $= \dfrac{P\left( {\mu  - \sigma  < X \leq  \mu  + \sigma }\right)  + P\left( {\mu  - {2\sigma } < X \leq  \mu  + {2\sigma }}\right) }{2} = \dfrac{{0.6827} + {0.9545}}{2} = \dfrac{1.6372}{2}  = {0.8186}$ .
                
                \item 参与活动的每位居民得分低于 74 分的概率为 ${0.1} + {0.3} + \dfrac{4}{10} \times  {0.3} = \dfrac{13}{25}$
                
                得分不低于 74 分的概率为 $\dfrac{12}{25}$
                
                $Y$ 的所有可能取值分别为10,20,30,40.
                
                $P\left( {Y = {10}}\right)  = \dfrac{13}{25} \times  \dfrac{2}{3} = \dfrac{26}{75},P\left( {Y = {20}}\right)  = \dfrac{13}{25} \times  \dfrac{1}{3} + \dfrac{12}{25} \times  {\left( \dfrac{2}{3}\right) }^{2} = \dfrac{87}{225} = \dfrac{29}{75}$
                
                $P\left( {Y = {30}}\right)  = \dfrac{12}{25} \times  2 \times  \dfrac{2}{3} \times  \dfrac{1}{3} = \dfrac{48}{225} = \dfrac{16}{75},P\left( {Y = {40}}\right)  = \dfrac{12}{25} \times  {\left( \dfrac{1}{3}\right) }^{2} = \dfrac{12}{225} = \dfrac{4}{75}$
                
                所以 $Y$ 的概率分布为
                
                \begin{center}
                \begin{tabular}{|c|c|c|c|c|}
                \hline
                Y & 10 & 20 & 30 & 40 \\
                \cline{1-5}
                $P$ & $\dfrac{26}{75}$ & $\dfrac{29}{75}$ & $\dfrac{16}{75}$ & $\dfrac{4}{75}$ \\
                \cline{1-5}
                \hline
                \end{tabular}
                \end{center}
                
                所以 $E\left( Y\right)  = {10} \times  \dfrac{26}{75} + {20} \times  \dfrac{29}{75} + {30} \times  \dfrac{16}{75} + {40} \times  \dfrac{4}{75} = \dfrac{296}{15}$ ,
                
                所以本次活动需要准备的话费充值卡的总金额为 ${240} \times  \dfrac{296}{15} = {4736}$ 元.
            \end{enumerate}
        \end{jieda}
        
        \item 目前，教师职业越来越受青睐，考取教师资格证成为不少人的就业规划之一. 当前, 中小学教师资格考试分笔试和面试两部分, 笔试通过后才能进入面试环节. 已知某市 2024 年共有 10000 名考生参加了中小学教师资格考试的笔试，笔试成绩 $\xi  \sim  N(60, 10^2)$，只有笔试成绩高于 70 分的考生才能进入面试环节.
        \begin{enumerate}
            \item 利用正态分布的知识，估计该市报考中小学教师资格的 10000 名笔试考生中，进入面试的人数(结果只保留整数)
            \item 现有甲、乙、丙 3 名考生进入了面试，且他们通过面试的概率分别为 $\dfrac{3}{4},\dfrac{2}{3},\dfrac{1}{2}$ ，设这 3 名考生中通过面试的人数为 $X$ ,求随机变量 $X$ 的分布列和数学期望.
        \end{enumerate}
        参考数据: 若 $X \sim  N\left( {\mu ,{\sigma }^{2}}\right)$ ,则 $P\left( {\mu  - \sigma  \leq  X \leq  \mu  + \sigma }\right)  \approx  {0.6827},P\left( {\mu  - {2\sigma } \leq  X \leq  \mu  + {2\sigma }}\right)  \approx  {0.9545}$ , $P\left( {\mu  - {3\sigma } \leq  X \leq  \mu  + {3\sigma }}\right)  \approx  {0.9973}.$
        \begin{jieda}
            \begin{enumerate}
                \item 由题意可知 $\mu  = {60},\sigma  = {10}$，则 $P\left( {\xi  > {70}}\right)  = P\left( {\xi  > \mu  + \sigma }\right)  = \dfrac{1 - P\left( {\mu  - \sigma  \leq  \xi  \leq  \mu  + \sigma }\right) }{2} \approx  \dfrac{1 - {0.6827}}{2} = {0.15865}$，则共 ${10000} \times  {0.15865} = {1586.5}$ ,即 1586 人进入面试.
                
                \item 由题意可知，随机变量 $X$ 的可能取值有 0,1,2,3,
                
                甲、乙、丙 3 名考生没通过面试的概率分别为 $1 - \dfrac{3}{4} = \dfrac{1}{4},1 - \dfrac{2}{3} = \dfrac{1}{3},1 - \dfrac{1}{2} = \dfrac{1}{2}$ ,
                
                则 $P\left( {X = 0}\right)  = \dfrac{1}{4} \times  \dfrac{1}{3} \times  \dfrac{1}{2} = \dfrac{1}{24}$ ,
                
                $P\left( {X = 1}\right)  = \dfrac{3}{4} \times  \dfrac{1}{3} \times  \dfrac{1}{2} + \dfrac{1}{4} \times  \dfrac{2}{3} \times  \dfrac{1}{2} + \dfrac{1}{4} \times  \dfrac{1}{3} \times  \dfrac{1}{2} = \dfrac{1}{4}$
                
                $P\left( {X = 2}\right)  = \dfrac{3}{4} \times  \dfrac{2}{3} \times  \dfrac{1}{2} + \dfrac{3}{4} \times  \dfrac{1}{3} \times  \dfrac{1}{2} + \dfrac{1}{4} \times  \dfrac{2}{3} \times  \dfrac{1}{2} = \dfrac{11}{24}$
                
                $P\left( {X = 3}\right)  = \dfrac{3}{4} \times  \dfrac{2}{3} \times  \dfrac{1}{2} = \dfrac{1}{4}$
                
                故随机变量 $X$ 的分布列为:
                
                \begin{center}
                \begin{tabular}{|c|c|c|c|c|}
                \hline
                $X$ & 0 & 1 & 2 & 3 \\
                \cline{1-5}
                $P$ & $\dfrac{1}{24}$ & $\dfrac{1}{4}$ & $\dfrac{11}{24}$ & $\dfrac{1}{4}$ \\
                \cline{1-5}
                \hline
                \end{tabular}
                \end{center}
                
                故 $E\left( X\right)  = 0 \times  \dfrac{1}{24} + 1 \times  \dfrac{1}{4} + 2 \times  \dfrac{11}{24} + 3 \times  \dfrac{1}{4} = \dfrac{23}{12}$
            \end{enumerate}
        \end{jieda}
    \end{enumerate}
\end{enumerate}